\documentclass[10pt]{book}
\usepackage[a5paper,%This enables easy large-letter editions (A4)
bindingoffset=5mm,%having an inner margin offset prevents text falling into the gutter
inner=20mm,
outer=20mm,
top=20mm,
bottom=20mm,
twoside,
heightrounded]{geometry}

%Basic English formatting:
	\frenchspacing %Prevents the silly double space after each sentence
	\usepackage{pifont} %Needed for \ding characters (for the crosses mainly)
	\usepackage{ebgaramond} %Standard open source font
	\setlength{\parindent}{1em}
	
%To ensure language compatibility.
	\usepackage[utf8]{inputenc}
	\usepackage[greek,latin,british]{babel}

%To change tracking and tighten up spacing:
	\usepackage{microtype}

%Makes sure odd pages are on the right when moving between page number types.
	\usepackage[strict]{changepage}
	\usepackage{scrextend}
	
%URL
	\usepackage[hidelinks]{hyperref}	

%Page Headers
	\usepackage{fancyhdr}
	\pagestyle{fancy}
	\fancyhf{}
	\renewcommand{\headrulewidth}{0pt} %Prevent line at top of page
	\fancyfoot[C]{\thepage}
%Prevents chapter from resetting page layout
	\fancypagestyle{plain}{
		\fancyhf{}
		\fancyfoot[C]{}
		\renewcommand{\headrulewidth}{0pt}}

%Change section format:
\setcounter{secnumdepth}{0} %Remove section numbering
	\usepackage{titlesec}
	\titleformat{\section} % Command to modify
		{\Large \normalfont \centering} % Format
  		{} % No label (removes number)
  		{0pt} % Spacing between label and title (irrelevant here since no label)
  		{}
	
	\titleformat{\subsection}
	{\normalsize\bfseries\scshape\centering}
  	{}
  	{0pt}
	{}
	
	\titleformat{\subsubsection}
	{\normalsize\centering\itshape}
  	{}
  	{0pt}
	{}

%Colour \& Images

	%Image Configuration
	\usepackage{graphicx}
	\graphicspath{{resources/images}}
	\usepackage{wrapfig}
	\usepackage{epsfig}

	%Enables text under image without caption.
	\usepackage{caption}
	\captionsetup[figure]{labelformat=empty,justification=centering}
	\usepackage{float}

	%Colour
	\usepackage[svgnames,dvipsnames]{xcolor}

	%Changes to the Table of Contents
	\addto\captionsbritish{\renewcommand{\contentsname}{Table of Contents}}
	\setcounter{tocdepth}{1} %Only chapters & sections are added to the TOC
	\usepackage{tocloft}
	\setlength{\cftbeforetoctitleskip}{-8ex}

	%For quotations in the Foreword and Editor's Preface:
	\usepackage{quoting}

	%For the signature in the Preface \& Margins
	\usepackage{ragged2e}

	%For the Chant Notation
	\usepackage[autocompile]{gregoriotex}
	%I'm not sure what this does.
	\gresetheadercapture{initial-style}{gresetinitiallines}{}

%Rubrics:
	%Colour for print:
	\definecolor{RubricRed}{cmyk}{0, 1, 1, 0.25}

	%For black words in rubrics.
	\newcommand{\blackrubric}[1]{{\color{black}#1}}
	
	%Heading Rubrics
	\newenvironment{inhead}{\par\centering\color{RubricRed}\small\smallskip\ignorespaces}
	{\par\normalsize\smallskip\ignorespacesafterend}
	
	%Normal Rubrics
	\newenvironment{rubric}{\par\smallskip\par\noindent\color{RubricRed}\small\ignorespaces\P}
		{\par\normalsize\smallskip\ignorespacesafterend}
		
	%For rubrics right after sections
	\newenvironment{secrubric}{\par\noindent\small\ignorespaces\color{RubricRed}\P}
	{\par\normalsize\smallskip\ignorespacesafterend}
	
	%Inline Rubric
	\newcommand{\inrub}[1]{{{\color{RubricRed}\small(#1)}}}
	
	%Determine Marginnote parameters
	\setlength{\marginparwidth}{15mm} % Set the width of the margin note
	\setlength{\marginparsep}{2mm} % Set the spacing between the margin note and the text
	\usepackage{marginnote}

	%Margin Rubric on Odd Pages
	
	\newcounter{alphacount} %Counter to convert number to letter
\newcommand{\margrub}[2]{\setcounter{alphacount}{#1}\textsuperscript{\alph{alphacount}}%
	\marginpar{\begin{minipage}{0.65\marginparwidth}
	\sloppy%
	\itshape%
	\scriptsize%
	\textsuperscript{\alph{alphacount}}%
	#2
	\fussy
	\end{minipage}}
  \hspace{-0.2em}}
 
  
  	%Margin Rubric on Even Pages
  \newcommand{\margrubleft}[2]{\setcounter{alphacount}{#1}\textsuperscript{\alph{alphacount}}%
	\marginpar{\RaggedLeft%
	\sloppy%
	\itshape%
	\scriptsize%
	\textsuperscript{\alph{alphacount}}%
	#2%
	\fussy}
  	\hspace{-0.2em}}
  
  	%Margin Rubric in Column Contexts (English Column)
  \newcommand{\margcolrub}[2]{\setcounter{alphacount}{#1}%
	\textsuperscript{\alph{alphacount}}%
	\Ifthispageodd{\reversemarginpar
	\marginpar{%
	\raggedright%
	\sloppy%
	\itshape%
	\scriptsize%
	    \textsuperscript{\alph{alphacount}}%
	    #2\hspace{-0.2em}%
	    \fussy}
	    }{\normalmarginpar
	    \marginpar{%
		\raggedleft%
		\sloppy%
		\itshape%
		\scriptsize%
		\textsuperscript{\alph{alphacount}}%
	    #2\hspace{-0.2em}%
	    \fussy}}}
  
  %Really just a placeholder in case it becomes relevant later in the latin side of a double column:
  \newcommand{\marglatincolrub}[2]{\setcounter{alphacount}{#1}\textsuperscript{\alph{alphacount}}\hspace{-0.2em}}
	
		%Margin Rubric without letter/number.
	  \newcommand{\margy}[1]{\marginpar{\RaggedRight%
		\sloppy%
		\itshape%
		\scriptsize%
		#1%
		\fussy}}
		
		%Margin Rubric without letter/number, on the left side (even numbered page).
	  	\newcommand{\margyleft}[1]{\marginpar{\RaggedLeft%
		\sloppy%
		\itshape%
		\scriptsize%
		#1%
		\fussy}}
  
  
		
	  
	%Kalendar \& Lectionary
	\usepackage{array}
	\usepackage{longtable} %Allows lectionary cells to go over multiple pages
	\newcolumntype{C}[1]{>{\centering\arraybackslash}p{#1}}
	\newcolumntype{L}[1]{>{\RaggedRight\arraybackslash}p{#1}}
	
		%Lectionary Evensong
		\newcommand{\evensong}[1]{Evensong}
	
	%Scriptural Citations in Mass
	\newcommand{\vr}[1]{{\footnotesize{ (#1)}}}
	
	%Disable \vr
	\newcommand{\novr}[1]{}
	
	%Invitatory
	\newcommand{\inv}[1]{{\color{RubricRed}#1.}}

%Propers:

	%Drop Caps (And specifies a specific folder for normal lettrine images):
	\usepackage{lettrine}
	\newcommand{\lett}[2]{\lettrine[image=true,realheight=true,lines=3]{dropcaps/#1.png}{{\lowercase{#2}}}}
	
	%If no second or third lines, and line amount might be doubtful, then this is used to ensure adequate spacing.
	\newcommand{\lettspace}[0]{\ifnum\prevgraf=1
		\vspace{2\baselineskip}
			\else
				\ifnum\prevgraf=2
					\vspace{1\baselineskip}
				\fi
			\fi}
	
	
	%Ensure that short paragraphs don't lead to text moving into drop cap.
	
	\newcounter{thirdlinecounter} %Counter to tell thirdline how many paragraphs precede it.
	\newcommand{\secondline}[1]{%
		\ifnum\prevgraf>2 %If there are over two lines in the previous paragraph (3), then we don't need to adjust anything, and it is put as is.
			#1\par
			\setcounter{thirdlinecounter}{3}
		\else
			\ifnum\prevgraf=2%If there are only two lines, then we only need to do an indent.
				\parindent=3.25em
				#1\par
            	\parindent=1em%Reset indent to normal
            	\setcounter{thirdlinecounter}{2}
			\else
				\ifnum\prevgraf=1%If the lett paragraph is only one line, then we need to cause the entirety of the next paragraph to be indented.
				\noindent
				\leftskip=3.25em
				#1\par
				\leftskip=0em%Reset indent to normal
				\setcounter{thirdlinecounter}{1}
					\else%In case something is not caught.
						#1
						\setcounter{thirdlinecounter}{3}
						\par
				\fi
			\fi
		\fi}
	
	%Make sure proper formatting for the third line.
	\newcommand{\thirdline}[1]{\addtocounter{thirdlinecounter}{\prevgraf}%
		\ifnum\thethirdlinecounter=2%If there are only two lines preceding (one line of the lett and one line of the secondline
    		\parindent=3.25em
    		#1\par
    		\parindent=1em
				\else
	    		#1\par
		\fi}
	
	\usepackage[super]{nth} %To add superscript to ordinal numeral.
	
	%Rank, Date, Name for Feast Days
	\newcommand{\festival}[3]{\par\subsection{#2. #3}
	
	\begin{inhead}
		#1
	\end{inhead}}
	
	%Headings for the Propers
	\newcommand{\feastday}[1]{\fancyhead[RO,LE]{\textit{#1}}} %Feast Name
	
	\newcommand{\feastdate}[2]{\fancyhead[RE,LO]{#1 #2}} %Feast Date
	
	%Psalm Numbers
	\newcommand{\psanum}[1]{{\small{#1}} }
	
	%For drawing boxes around unapproved text:
	\usepackage[most]{tcolorbox}
	
	%To make sure versicle and response start with indent.
	\usepackage{indentfirst}

	%To create hanging indent for Hymns.
	\usepackage{hanging}
	
	%Introit
	\newcommand{\introit}[0]{\subsubsection{Introit}}
	
	%Collect
	\newcommand{\collect}[0]{\subsubsection{Collect}}

	%Readings
	\newcommand{\readingcitation}[2]{\subsubsection{#1. \emph{#2}}}
	
	%Divine Name in Readings
	\newcommand{\divineName}[1]{\textsc{\lowercase{#1}}}
	
	%Gradual \& Alleluia
	\newcommand{\gradall}[2]{\gradual{#1}\par\alleluia{#2}}
	
	%Gradual \& Tract
	\newcommand{\gradtr}[2]{\gradual{#1}\par\tract{#2}}
	
	%Alleluia Verse
	\newcommand{\alleluia}[1]{\vspace{1ex}\par\noindent
		\emph{Alleluia.} #1}
	%Gradual
	\newcommand{\gradual}[1]{\vspace{1ex}\par\noindent
		\emph{Gradual.} #1}
		
	%Tract
	\newcommand{\tract}[1]{\vspace{1ex}\par\noindent
		\emph{Tract.} #1}
		
    %Offertory
    \newcommand{\offertory}[1]{\vspace{1ex}\par\noindent
		\emph{Offertory.} #1}
		
	%Secret
	\newcommand{\secret}[0]{\subsubsection{Secret}}

	%Communion
	\newcommand{\communion}[1]{\vspace{1ex}\par\noindent
    	\emph{Communion.} #1}
	
	%Postcommunion
	\newcommand{\postcommunion}[0]{\subsubsection{Postcommunion}}
		
	%Antiphons
	\newcommand{\antiphon}[2]{{\centering \textit{#1 Ant.} #2}}

\usepackage{paracol} %To provide balanced English-Latin dual paragraphs
\usepackage[grid]{multicol} %For non bilingual dual columns.
\setlength\columnseprule{0.4pt} %Controls the width of the line between columns

%Prevent lettrine misalignment:
\raggedbottom

%English-Latin Double Column
\newcommand{\elcol}[2]{\sloppy
	\begin{paracol}{2}
		\begin{leftcolumn}
			#1
		\end{leftcolumn}
		\begin{rightcolumn}
			\selectlanguage{latin}%
			#2
		\end{rightcolumn}
		\switchcolumn*
	\end{paracol}
\selectlanguage{british}
\fussy}

%English-Latin Non-Indented Column
\newcommand{\elcoldent}[2]{\sloppy
	\begin{paracol}{2}
		\begin{leftcolumn}
		\noindent%
			#1	
		\end{leftcolumn}
		\begin{rightcolumn}
                \selectlanguage{latin}%
			\noindent
			#2
		\end{rightcolumn}
		\switchcolumn*
	\end{paracol}
\selectlanguage{british}
\fussy}

%For difference between through the year and in eastertide:
\newcommand{\ttyie}[2]{\sloppy
	\begin{paracol}{2}
		\begin{leftcolumn}\noindent
			#1	
		\end{leftcolumn}
		\begin{rightcolumn}\noindent
			#2
		\end{rightcolumn}
		\switchcolumn*
	\end{paracol}
\selectlanguage{british}
\fussy}

%Expedites adding dual language Oremus/Let Us Pray
\newcommand{\oremus}[0]{\par
	\begin{paracol}{2}
		\begin{leftcolumn}
			\selectlanguage{british}\noindent%
			\centerline{\textit{Let us pray.}}
		\end{leftcolumn}
		\begin{rightcolumn}
			\selectlanguage{latin}\noindent%
			\centerline{\textit{Orémus.}}
		\end{rightcolumn}
		\switchcolumn*
	\end{paracol}
\selectlanguage{british}}

%Does the same as above, but for single language collects:
\newcommand{\letuspray}[0]{\par\noindent
	\centerline{\textit{Let us pray.}}
	\par}

\newcommand{\oremuslatin}[0]{\par\noindent
	\centerline{\textit{Orémus.}}
	\par}
	
	

	
\begin{document}
\begin{titlepage}
		\begin{center}
			{\Large{The Book of}}
			\par
			\vspace{3mm}
			{\Huge{\scshape Common Prayer}}
			\par
			\vspace{3mm}
			{\Large{and Administration of the}}
			\par
			\vspace{3.5mm}
			{\Huge{\scshape Sacraments,}}
			\par
			\vspace{3mm}
			{\large{and other}}
			\par
			\vspace{3.5mm}
			{\Huge{\scshape Rites and Ceremonies}}
			\par
			\vspace{3.5mm}
			{\large{of the}}
			\par
			\vspace{3mm}
			{\Huge{\scshape Church}}
			\par
			\vspace{3.5mm}
            {\large{Proposed for the Use in}}
			\par
			\vspace{3mm}
			{\Huge{\scshape the Antiochian}}
            \par
			\vspace{3mm}
            {\Huge{\scshape Western Rite Vicariate}}
            \par
			\vspace{2mm}
   		 	\centerline{\rule{0.5\textwidth}{0.4pt}}
			\par
			\vspace{4mm}
			{\LARGE{Together with the Psalter}}
			\par
			\vspace{3mm}
			{\large{or}}
			\par
			\vspace{3.5mm}
			{\LARGE{Psalms of David}}
   \par
   \end{center}
            
            \vfill
            
           \begin{figure}[H]
            	\centering
            	\includegraphics[scale=.15]{logo.eps}
            	\caption{\scshape Apologia Anglicana\\Boston, Massachusetts\\
				A. D. MMXXV}
            \end{figure}
		
	\end{titlepage}
\cleardoublepage
\frontmatter
\tableofcontents
\addtocontents{toc}{\vspace{-3ex}}%
\cleardoublepage
%Front Matter
%Roman numeral page numbers for the front material.
\setcounter{page}{1}
\frontmatter
\phantomsection
\addcontentsline{toc}{chapter}{Copyright}
\fancyhead[C]{{\Large Copyright}}
\noindent
The 2025 Proposed Book of Common Prayer ($\copyright$ 2025 Apologia Anglicana, LLC) is licenced under CC BY-SA 4.0. To view a copy of this license, visit \url{https://creativecommons.org/licenses/by-sa/4.0/}\\
\par
ISBN: 978-1-969461-04-0 (Personal Size Hardcover)

ISBN: 978-1-969461-05-7 (Altar Size Hardcover)

ISBN: 978-1-969461-06-4 (eBook)\\

Library of Congress Control Number: 2025920564\\
\par\noindent
Edited by J. P. `Augustine' Watson (Augustine@ApologiaAnglicana.org).\par\noindent
Front cover designed by Matthew Taylor (@semperadiuvans).

\begin{center}
	\textsc{Please note that this is not currently an official text of the Antiochian Western Rite Vicariate.}
\end{center}

\par\noindent
We are thankful to have such great giants upon whose shoulders we stand. These works are used in this Book according to fair use, where their content is not already in the public domain. What was used from the book, or in what way it was consulted, is described under each citation.\\

\textit{A Manual for Priests}. \nth{5} edition. Athens: Anglican Parishes Association, 2004.\par\noindent
\textit{The Burial rites.}\\

Davis, C. Lance, ed. \textit{The Anglican Office Book}. \nth{2} edition. Chester: Whithorn Press, 2023.\par\noindent
\textit{Lessons and Psalms in the Daily Office Lectionary, which are not provided in the 1928 Book of Common Prayer. The Invitatory antiphon for Septuagesimatide. Some of the text and Rubrics for the Minor Hours. Other Daily Office propers.}\\

\textit{The American Missal: Revised \& Expanded}. Glendale: Lancelot Andrewes Press, 2010.\par\noindent
\textit{The propers for Our Lady of Walsingham.}\\

The Armed Forces Division of the Protestant Episcopal Church. \textit{The Armed Forces Prayer Book}. New York: The Church Pension Fund, 1951.\par\noindent
\textit{The Prayer for the Armed Forces.}\\

%MANUAL ADJUSTMENT:
\clearpage
\textit{The English Missal for the Laity}. London: W. Knott \& Sons Limited, 1933.\par\noindent
\textit{Mass Rubrics. The Litany of Saints. The Thanksgiving for Holy Communion. Parts of the Mass not found in the Prayer Book Tradition. Propers not found in the Prayer Book Tradition nor used from other books.}\\

The General Synod of the Anglican Church of Canada. \textit{The Book of Common Prayer}. Toronto: Anglican Book Centre, 1962.\par\noindent
\textit{The Invitatory Antiphon for Lent.}\\

\textit{The Monastic Diurnal}. Oxford University Press, 1932.\par\noindent
\textit{Office Rubrics. Ranking of Feast Days. Marian Anthems. Office of the Dead. Sacrament of Penance. Preparation for Holy Communion.}\\

%Everything which could be said to be from here is found in the AOB.
%Hartzell, Paul, ed. The Prayer Book Office. \nth{2} edition. New York: Morehouse-Barlow Co., 1963.\par\noindent
%\textit{Lessons and Psalms in the Daily Office Lectionary, which are are not provided in the 1928 Book of Common Prayer.}
\textit{Orthodox Ritual}. St. Luke's Priory Press, 2016.\par\noindent
\textit{Translations for some of the prayers in the Sacraments of Baptism and Confirmation, as well as in other Pastoral Rites.}
\clearpage
\include{pages/Front/Dedication}
\phantomsection
\addcontentsline{toc}{chapter}{Foreword}
\fancyhead[C]{{\LARGE Foreword}}
\noindent
\lett{T}{he} publication of this edition of the Book of Common Prayer, carefully edited and directed toward those in the Antiochian Western Rite Vicariate, marks a contemplative moment of continued growth of the Orthodox Catholic faith in North America and in the territories of British Christendom. Much has already been written about the nearly universal acceptance of the Western Rite among our Holy Fathers in modern times, and so this foreword is not concerned with such a topic. The message I hope to leave with the faithful is the importance of beauty, and the importance of Orthodox Christianity being the Great Healer of nations.\\

The Orthodox Church has always emphasized the importance of beauty in the salvation of mankind, for as it says in the Epistle of James,

\begin{quoting}\noindent
	Every good gift and every perfect gift is from above, and cometh down from the Father of lights, with whom is no variableness, neither shadow of turning. ---James~1:17
\end{quoting}

In our homeland of North America, many of us have been led to Orthodoxy by the beauty of the English tradition of prayer which was given to our forefathers (according to the flesh) by way of our English and Scottish heritage. As the peoples of the British Isles arrived in the colonial New World, so too did their prayer tradition which was rooted in a deep connection to the blessed Dowry of Mary, that is, the Albion of Saint Alfred and the Three Holy Patrons: St. Edmund, St. Edward, and St. Gregory the Great. Even the Congregationalists, Puritans, and ``Low Churchmen'' could not avoid retaining an echo of the ancient Liturgy within their houses of prayer. \\

This deep connection to the England once known as Mary's Dowry is what allowed for beauty to thrive in a deeply desolate new world. Amidst the darkness of isolation and uncertainty, New Englanders and Southern Americans built lasting houses of prayer adorned with the beauty of Christ's Name and images of His saints. Within those halls, the men and women of an uncertain new world prayed the collects, hymns, and prayers of an ancient inheritance, albeit commingled with new and unfamiliar doctrines and liturgics that set a broad path against what was once dedicated to the very Narrow Way of Christ. In modern times, the Christians of the English tradition have found themselves near shipwreck and in danger of losing the beauty of their rich past. Though the valuable buildings, real-estate, and stained glass have passed on to ``the Religion of the Future,'' as our holy father St. Seraphim of Platina called it, the prayers and liturgical life of English Christians can be still be made firm in the life giving Water of Christ. This life giving water runs out from the side of the Living Christ, from His Church, which boldly proclaims the Catholic faith without shame.\\

The thirst for the living water of the ancient Orthodoxy of the Christian life was never fully stomped out by the errors of the dialectical tension of Roman Catholicism and Anglicanism. Much has been said before about the elegant writing of William Laud, Richard Hooker, Lancelot Andrewes, and George Herbert. And so I wish to offer a sometimes forgotten example from the Right Reverend William Seabury of Connecticut. Rt. Rev. Seabury was the first Angli-can bishop consecrated in the United States and, writing in favor of aligning the American Anglican ethos with that of the Scottish prayerbook, he wrote:

\begin{quoting}\noindent
	The grand fault in that [1662] office [of Holy Communion] is the deficiency of a more formal oblation of the elements, and of the invocation of the Holy Ghost to sanctify and bless them. The Consecration is made to consist merely in the Priest's laying his hands on the elements and pronouncing ``This is my body,'' etc., which words are not consecration at all, nor were they addressed by Christ to the Father, but were declarative to the Apostles . . . The efficacy of Baptism, of Confirmation, of Orders, is ascribed to the Holy Ghost, and His energy is implored for that purpose; and why he should not be invoked in the consecration of the Eucharist, especially as all the old Liturgies are full to the point, I cannot conceive. ---\emph{The Life and Correspondence of Samuel Seabury}, p. 354
\end{quoting}

We see in this early American bishop a desire to return to the ancient liturgies, to live within the uncreated energies experienced and preached by St. Paul, and to find what had been lost after the Great Schism and the deeply regrettable brother wars of the Reformation which led to the destruction, desolation, and murder of European Christianity.\\

Herein we find the purpose of the 2025 Proposed Book of Common Prayer: to bring forth the Good, to let go of all that which is contrary to the Truth, and to heal what is infirm. Following carefully the recommendations of the fathers of the Antiochian and Russian Orthodox Churches, the 2025 Proposed Book of Common Prayer has only sought to remain loyal to what the Orthodox Church has already observed and ruled. Recalling the conclusion to the Russian Observations on the American Prayerbook, it is worth restating their findings:

\begin{quoting}\noindent
	The committee, after reviewing these ``Observations,'' allowed in general the possibility that if Orthodox parishes, composed of former Anglicans, were organized in America, they might be allowed, at their desire, to perform their worship according to the ``Book of Common Prayer,'' but only on condition that the following corrections were made in the spirit of the Orthodox Church. ---\emph{Russian Observations upon the American Prayer Book} 2:VI
\end{quoting}

The 2025 Proposed Book of Common Prayer is not a restoration project, nor is it a historical recreation. Rather it is a spiritual surgery for those of the Anglican patrimony who seek communion with the One, Holy, Catholic, and Apostolic Church and to retain that which the ancient Church provided, influenced, and now attempts to sustain.\\

Continued prayers are also asked for all those of the Western Rite, so that a spirit of pretest, pride, or confusion may never yoke itself to the Latin and English presence within the Orthodox world. All things are submitted to the Holy Church so that Christ may judge, through His living Mystical Body.\\

\rightline{\emph{The Blessings of Our Lord Jesus Christ,}}

\vspace{1\baselineskip}

\rightline{The Reverend Father Justin Slaughter Doty}

\rightline{\emph{Priest, St. Euphemia Orthodox Christian Church}}

\rightline{\emph{Overseeing St. Botolph Orthodox Chapel}}

\phantomsection
\addcontentsline{toc}{chapter}{Editor's Preface}
\fancyhead[C]{{\LARGE Editor's Preface}}
\noindent
%\begin{quoting}
%	How dreadful is this place! this is none other but the house of God, and this is the gate of heaven. ---Opening Sentence for the Feast of Our Lady of Walsingham
%\end{quoting}

\lett{W}{e} rejoice on the occasion of the publication of this edition of the Book of Common Prayer. As Orthodox Catholics, we have been blessed to have received the Gospel and full membership within the Mystical Body of Christ. And by the undeserved gift bestowed upon us, we offer unto the Almighty Father `our worship and praise and thanksgiving'. In light of so great a gift, it is only fitting that we have a Book for so great a worship.

\vspace{0.5\baselineskip}

\emph{The Proposed 2025 Book of Common Prayer} is a revision of the 1928 Book of Common Prayer for use within the Orthodox Church, specifically the Antiochian Western Rite Vicariate. Where the 1928 BCP was deficient, it was changed or supplemented, either from the Prayer Book Tradition itself or other books, such as the Roman Missal and the Monastic Office.

\vspace{0.5\baselineskip}

For those familiar with the Anglican tradition, this Book will be very familiar. We took great care to not remove or change anything unless necessary or truly expedient. And where we changed or removed, we took even greater care to not follow our personal eccentricities or idiosyncrasies. Therefore, this Book stands as a true successor within the Prayer Book Tradition: both an integral expression of the English Rite as well as truly and wholly Orthodox.

\vspace{0.5\baselineskip}

%And for those unfamiliar with the English Rite, we believe that this will still provide great benefit: an accessible Book to offer true worship unto the Most Holy Trinity within the life of the Church.\\

We are thankful for all the support we have received. This Book simply would not be possible without the many laymen who spent their money, time, and energy to buy draft copies, use them in their prayer life, and offer corrections and suggestions. This is to say nothing of the priests and other clergy who offered their valuable time (of which they do not have much) to review the Book of Common Prayer and give indispensable feedback. We are especially thankful for the priests who took the time to study and use the Book of Common Prayer. And above all, thanksgiving must be given to the Right Reverend \textsc{John} of Worcester who oversees the Antiochian Western Rite Vicariate with the greatest charity and prudence, and who has pastorally cared for those involved with this project.

\vspace{0.5\baselineskip}

Please pray for this project, that this Book may, in the end, receive official approbation and full reception within the Vicariate. And please pray for our salvation as we pray for yours.\\

\rightline{On the Feast of Our Lady of Walsingham}

\vspace{1\baselineskip}

\rightline{\emph{Pax Christi,}}

\vspace{1\baselineskip}

\rightline{Mr. Augustine Watson, MSt}

\rightline{\emph{Chief Editor, Apologia Anglicana}}
\phantomsection
\addcontentsline{toc}{chapter}{Rubrics}
\fancyhead[C]{{\LARGE Rubrics}}
\subsection{Preface}
\lett{I}{n} the Christian life, the Holy Sacrifice of the Mass and the Daily Office serve as the highest modes of prayer. Due to their solemnity, it is important to attend carefully to their proper celebration.

	Through the Mass, the highest sacrifice is offered unto the Most Holy Trinity, and Christ is made substantially present upon the Altar. In the Office, the Christian obediently, faithfully, and joyfully observes the Sacred Scriptures' exhortation to offer the morning and evening sacrifice.
	
	The Christian who prays this Book, both layman and cleric, must keep the rubrics and their purpose within these prayers in mind. At the same time, however, the rubrics are to help the Christian offer his Sacrifice with loving care and faithfulness, not to bog him down with obligations. The Rites and Rubrics in this Book are laid out in a manner that should be accessible, especially after a few uses, lest `there be more business to find out what should be read, than to read it when it be found out.'
%It is in this mode of prayer that the Christian obediently, faithfully, and joyfully observes the Sacred Scripture's exhortation to offer the morning and evening sacrifice. And, through devotion to the Minor Hours, to `pray without ceasing'.\par
%It must be remembered, before a review of the rubrics, that the purpose of the Office is to facilitate this Sacrifice. Therefore, a simpler Office or an Office devoutly offered with imprecise rubrical distinctions, is immensely more pleasing than a despondent spirit.\par
%\emph{
%}
\subsection{General Rubrics}
\begin{description}
\item[Amen] When the word \emph{Amen} is in italics, it should be said by the Congregation in response to the Minister's prayer.
%This modifies the 1928 rubric, in line with the 1549 and 1662, especially with consideration of Orthodox doctrine.
\item[Antecommunion] It may occur that a community lacks a Priest on Sunday or another day where the Service of Holy Communion is desired. In such a case, Mattins and the Great Litany should first be said. Then, the Holy Mass may be said from the \emph{Collect for Purity} until the Sermon, inclusive, concluding with the Additional Collects (p. \pageref{AdditionalCollects}).\par
\textsc{Note,} Mattins, the Great Litany, and Antecommunion should all be led by the Minister in the same place as the Office, not at the Altar.
%From the Monastic Diurnal
\item[Collects] When several proper Collects (or Secrets, or Postcommunions) are said, only the first and the last receive the full ending.\par
	\textsc{Note,} In the Mass, the \emph{Lord be with you} is only said to introduce the first and second collects.
\item[Commemorations] When multiple Feast Days overlap on the same day, the Office will be of the Feast Day of highest rank. The other Feast Day(s) will then be commemorated (or, in some cases, suppressed or transferred). Commemoration consists of praying the highest-ranking Feast Day's collect first followed by the lower ranking Feast Day(s)' collect. In the Mass, this also applies to the Secret and Postcommunion.
%Collect Conclusion taken from the English Missal 1940 (it's not provided in the Lay Missal).
%\item[Conclusion of Collects] When the Collect is addressed to God the Father, the ending is `Through Jesus Christ thy Son our Lord; who liveth and reigneth with thee in the unity of the Holy Spirit, one God, world without end. Amen.'\par
%If in the beginning of the Collect the Son be mentioned: `Through the same Jesus Christ thy Son our Lord; who liveth and reigneth with thee in the unity of the Holy Spirit, one God, world without end. Amen.'\par
%If at the end of the Collect the Son be mentioned: `Who with thee liveth and reigneth in the unity of the Holy Spirit, one God, world without end. Amen.'\par
%If the Collect be addressed to the Son: `Who livest and reignest with God the Father in the unity of the Holy Spirit, one God, world without end. Amen.'\par
%When the Holy Spirit is mentioned in the Collect: `In the unity of the same Holy Spirit, etc.'
%Taking what I can from the 1933 Lay Missal:
\item[Conclusion of Collects] When the Collect is addressed to God the Father, the ending is `Through Jesus Christ thy Son our Lord, who liveth and reigneth with thee, in the unity of the Holy Ghost, God, throughout all ages, world without end. Amen.'\par
If in the beginning of the Collect the Son be mentioned: `Through the same Jesus Christ thy Son our Lord; who liveth and reigneth with thee in the unity of the Holy Ghost, God, throughout all ages, world without end. Amen.'\par
If at the end of the Collect the Son be mentioned: `Who liveth and reigneth with thee in the unity of the Holy Ghost, God, throughout all ages, world without end. Amen.'\par
If the Collect be addressed to the Son: `Who livest and reignest with the Father, in the unity of the Holy Ghost, God, throughout all ages, world without end. Amen.'\par
If the Collect be addressed to the Holy Spirit: `Who livest and reignest with the Father and the Son, God, throughout all ages, world without end. Amen.'\par
When the Holy Spirit is mentioned in the Collect: `In the unity of the same Holy Ghost, etc.'
\item[Dominus Vobiscum] The Minister may only proclaim the \emph{Dominus vobiscum} (The Lord be with you) if he be in Major Orders. Otherwise, he must instead lead with the \emph{Domine, exaudi orationem nostram} (O Lord, hear our prayer) wherever the \emph{Dominus vobiscum} appears, or omit it entirely if it would be said twice consecutively.
\item[Ferial Commemorations] In Advent \& Lent, the propers for the First Sunday of Advent and Ash Wednesday may be said when commemoration of their season's respective Feria is indicated.
\item[Gloria Patri] During Passiontide, the \emph{Gloria Patri} is omitted in the Invitatory of Mattins, Responsory of the Minor Hours, Aspersion of the People, Introit, and Lavabo. During the Sacred Triduum, the \emph{Gloria Patri} is wholly omitted.
\item[Latin] In much of this Book, Latin is provided alongside the English. The Latin may always be used instead of the English, unless otherwise determined by the Ordinary. When Latin is not provided, but great devotion exists for the Latin language, the Latin which serves as the base text (such as the Liber Precum Publicarum, Breviarium Monasticum, Missale, \& Rituale) may be used, with great concern for the catechesis and participation of the People.
\item[Liturgical Seasons] For the purpose of the rubrics, these festive seasons are referenced:
    \begin{itemize}
        \item Christmastide: From I Evensong of Christmas Day until I Evensong of Epiphany Day, exclusive.%\footnote{This is not to be confused with the `Christmas season', which can be said to end on Candlemas.}
        \item Epiphanytide: From I Evensong of Epiphany Day until I Evensong of Septuagesima Sunday, exclusive.
        \item Septuagesimatide (Pre-Lent): From I Evensong of Septuagesima Sunday until Mattins of Ash Wednesday, exclusive.
        \item Lent: From Mattins of Ash Wednesday until I Evensong of Passion Sunday, exclusive.
        \item Passiontide: From I Evensong of Passion Sunday until Mattins of Maundy Thursday, exclusive.
        \item Sacred Triduum: From Mattins of Maundy Thursday until I Evensong of Easter Sunday, exclusive.
        \item Eastertide (Paschaltide): From I Evensong of Easter Sunday until I Evensong of Trinity Sunday, exclusive.
        \item Trinitytide: From I Evensong of Trinity Sunday until I Evensong of the First Sunday of Advent, exclusive.
    \end{itemize}
\item[Minor Feast Days] Since only I and II Doubles are provided here in this Book (with some exceptions), the rubrics will be geared towards those. When the Office or Mass is celebrated with the rest of the Feast Days, the rubrics of the \emph{Prayer Book Hymnal} should be consulted.%\par
%The Mass propers and Office antiphons for the Feast Days not supplied in this Book may be found in the \textit{Book of Occasional Services}.\par
%For the rest of the Office propers, the Monastic Breviary should be consulted.

%Is this commentary really appropriate for the rubrics?:
% However, since these books have been published by Anglicans, great care must be taken to review the text in order to avoid any heterodox prayers.
\item[Mode of Recitation] Wherever it is indicated something be read, it may also be chanted, except for the Fore-Office and the Prayers after the Third Collect.
%\item[Plural \& Gendered Nouns and Pronouns] Sometimes there are words in the liturgy which change depending on number and gender, such as if the person being prayed for is a woman or if there are multiple people being prayed for. The Book, in accordance with English grammar, assumes the masculine singular, and if the prayer should be changed for either number or gender, the relevant words and letters are italicised.\par
%In cases where the change could only be due to number, such as prayers for Clerics (Subdeacons, Deacons, Priests, Bishops, etc.) which are---by divine law---always masculine, the relevant words are still italicised.
\end{description}

\subsection{Office Rubrics}
\begin{description}
%\item[Canticle Antiphons] %The antiphons for the \emph{Benedictus} and \emph{Magnificat} canticles are provided for those with special devotion to them.\par
%On Feast Days below Double, the antiphon may be said: as far as the {\dag} before, and repeated in full after, each Canticle.\par
%On Feast Days Double and higher, the antiphon may be said: entire before and after each Canticle.%\par
\item[Collects] The Collect to be said during the Office is provided in the relevant Proper. The Seasonal Prayers are to be said during Mass as indicated. The Seasonal Prayers (only its Collects) may be said as additional Prayers during the Office.
%The antiphons, in line with the Prayer Book Tradition, are optional.
\item[Ferial Days] The Ferial Office, that is, the Simple Office of the occurrent Season, is always said on weekdays in Lent, on Ember Days, and on Rogation Monday, unless a I or II Double occur, in which case the Office is of the Feast with Commemoration of the Feria.\par
On weekdays in Advent and between Septuagesima and Ash Wednesday, and on Common Vigils, the Ferial Office is always said unless a Double Feast or Octave occur, and then Commemoration is made of such Ferias.\par
If a Simple Octave Day occur on one of these Ferias, it is only commemorated.\par
In like manner throughout the year the Office is of the Feria on those weekdays on which there does not occur a Double Feast, an Octave, or the Office of St. Mary on Saturday.
\item[Gloria Patri] After every Psalm and the indicated Canticles, except in the Office of the Dead or where otherwise indicated, the \emph{Gloria Patri} should be said.\par
It is customary to reverence the Trinity by bowing one's head during the \emph{Gloria Patri}, instead of crossing oneself.
%\item[Great Litany] The petition `From all sedition, privy conspiracy, and rebellion' may be replaced by the version from the 1549 Book of Common Prayer.
\item[Lectionary] The Lessons and Psalms for the days of the Church Year are provided in the lectionary, according to the Rankings of Feast Days (p. \pageref{ranking}).\par
\textsc{Note,} The proper psalms for Second Class and lower Sundays are optional.
\item[Octaves] The Office of an Octave is said (or Commemoration thereof made, when it is hindered by a Feast or a Sunday) through eight continuous days in Octaves of Feasts of the I Class. Octaves of Feasts of the II Class, which are Simple Octaves, are kept only on the Octave Day itself, with the rite of a Simple, unless hindered by a more worthy Office; but no notice is taken of the days within the Octave. If any Octave is not ended before Ash Wednesday, the Vigil of Pentecost, or 17 December, no notice is taken of it thenceforth.\par
A Simple Octave Day occurring within an Octave is only commemorated.
\item[Octave Days] On the Octave Day, except for the proper Psalms and Lessons, the Office is said as on the day of the Feast, unless otherwise noted. On a Simple Octave Day the Office is Simple.
\par
On Sundays within Privileged Octaves, the Office is said as directed in the Proper of Season. On Sundays within Common Octaves, the Office is of the Sunday with Commemoration of the Octave.
%\item[Octave Day Lessons] Only on the Feast Day are the proper Psalms and Lessons of the Feast Day used, not also on the Octave Day. During the Octave, the Collect for the Feast is used, according to relevant rules for Octave rankings.
\item[Office of the Dead] The Office of the Dead may be prayed in addition to the Daily Office. It may also be prayed instead of the Daily Office on an ordinary Feria, a Memorial, or a Feast of Simple rank.\par
%CHECK:
The Office of the Dead may also replace the Daily Office of rank Double or lower on the day of death and the day of burial.
\item[Office of the Dead Antiphons] The Antiphons are not doubled except on 2 November; %14 November (at the Commemoration of All the Departed O.S.B., for Benedictines);
on the day of a burial; on the day after receiving tidings of a death; on the third, seventh, and thirtieth days; on the anniversary (even when transferred); and whenever the Office is celebrated solemnly.
\item[Office of the Dead Lessons] The lessons are always read without introduction or conclusion.
\item[Office of Our Lady on Saturday] On all Saturdays---except in Advent, Septuagesimatide, Lent, and on Ember Days---unless the Office be of a Double Feast (even transferred), or of an occurrent Octave or Vigil, or of a Sunday anticipated according to the rubrics, the Office is of St. Mary.
%\item[Proper Psalms \& Lesson] Proper Lessons are given for every day of the Church Year and for every Feast Day supplied in this Book (Prayer Book Holydays). Proper Psalms are supplied for every Sunday, Ember Days, and Prayer Book Holydays.\par
%Only Feast Days of II Double or higher receive Proper Psalms and Lessons, unless otherwise noted or supplied in this Book.
\item[Scripture Verse Numbers] When a lesson or psalm is listed with only one verse (such as John 1:45 or Psalm 102:15), the Minister reads the that verse and the verses following for the rest of the chapter.
\item[Suffrages] The Suffrages in the Minor Hours are not said on Feast Days Double or higher.
\item[Sunday Collects] %The Collect for a Sunday is said on the I Evensong, Mattins, and II Evensong of the Sunday, unless a Feast Day of higher rank is celebrated. It is also said throughout the week, unless a proper Collect is provided for the weekday.
If the Office be of a Feria without a proper Collect, then the Sunday Collect is read. Otherwise, the Collect for the Feast Day is read instead.
\item[Third Lesson] In accordance with pious custom, after the Second Canticle in the Major Hours, but before the Creed, it is fitting for a patristic lesson to be read from the Feast Day's Second or Third Nocturn in Monastic or Sarum Matins.
\item[Vigils] %Vigils, unlike modern Roman practice, are not anticipated celebrations of the following day. Rather, they are days in their own right, dedicated to fasting in preparation for the Feast Day.\par
The Office of a Vigil is only said at Mattins.\par
If a Vigil fall on a Sunday (except for the Vigils of Nativity \& Epiphany), it is said or commemorated on the preceding Saturday.
\end{description}
\subsection{Mass Rubrics}
\begin{description}
    \item[Chant] The Priest should be aware which parts of the liturgy, according to the solemnity of the Mass and the rank of the Feast Day, should be chanted.
    %\item[Commemorations] When Feast Days overlap, the highest ranking Feast Day is celebrated with commemoration of the lower ranking Feast Days, consisting of their Collect, Secret, and Postcommunion.
    \item[Conventual Masses] %The English Missal envisions that the Mass may be said in the context of a monastery, praying the full Monastic Breviary and having multiple Masses every day. Therefore, the conventual Mass is the higher Mass of the day, with large attendance, though the priests may celebrate a second Mass for a lower ranking Feast Day of that day, according to the rubrics.\par
    %Given that the Prayer Book Tradition does not envision many Masses being said in monastic fashion, rubrics from the English Missal regarding conventual and non-conventual Masses are largely irrelevant. For the ordinary parish church, it is safe to assume the day's Mass is a non-conventual Mass.
    Since this Book is oriented towards use in a parish church, it assumes the Mass being said is a `non-conventual Mass'.
    \item[Gloria in excelsis] The \emph{Gloria in excelsis} is said in the Mass on Feast Days of Double rank and higher. It is also said on all Sundays, except in Advent, Septuagesimatide, and Lent.\par
    \textsc{Note,} For special rubrics regarding Votive Masses and Ferias, the \emph{Ordo} or Missal Rubrics should be consulted.
    \item[Introduction of the Epistle] The Epistle is introduced according to manner specified in the rubric, preferably using the title of the book as given in the 1611 King James Bible.
    \item[Introit] The Introit for each Mass is provided. The Introit is said as written, followed by the \emph{Gloria Patri} (except in Passiontide) and the Introit repeated up to the `\emph{Ps}' or `\emph{Cant.}'
%    \item[Manual Acts] For the manual acts of clerics during the Mass, consultation and familiarity should be made with \emph{Ritual Notes} and the English Missal, especially during Christmastide and Holy Week.
    %Rubric from English Missal, reworded:
    \item[Nicene Creed] The Nicene Creed is said after the Gospel on all Sundays throughout the year, even if the Mass is of a Feast on which it would not otherwise be said. On all Feast Days of the Apostles, St. Joseph, Our Lady, the Holy Cross, and our Lord, and through their Octaves. From Christmas until the Octave of St. John, inclusive. On the Feast of the Epiphany and its Octave. Maundy Thursday. On Feasts of Angels, St. Mary Magdalene, St. Gregory, St. Ambrose, St. Augustine of Hippo, St. Jerome, St. Hilary, St. Isidore, Pope St. Leo I, St. Bede the Venerable, St. Peter Chrysologus, St. Athanasius, St. Basil, St. Cyril of Alexandria, St. Cyril of Jerusalem, St. Ephram Syrus, St. Gregory Nazianzen, St. John Chrysostom, St. John Damascene, St. Lawrence, St. John Baptist and his Octave Day. On All Hallows' Day and through its Octave, on the Dedication of St. Saviour, and of the Holy Apostles Peter and Paul. As well as on the anniversary of the dedication of a Church, and through its Octave; on the day of consecration of a Church, or of an Altar; on the Feasts of the Patron Saint of a Church, and on the day of a Saint where his Body or a considerable Relic is preserved; on the day of the creation and coronation of the Chief Bishop, and on the anniversary of such day; on the day, and the anniversary, of the election and consecration of a bishop.\par
    Also, the Creed is said in Votive Masses, which are celebrated solemnly for a grave matter or for a public Church cause, even if they are said in violet vestments on a Sunday.
    \item[Prefaces] The Prefaces with Solemn Chant are to be used in all Masses of at least Semidouble rank, and in Votive Masses for a grave and public cause.\par
    \textsc{Note,} When the Mass is of a Feast Day and has a proper Preface, its Preface is said in preference to the Season's.\par
    \textsc{Note,} When a Mass does not have a proper Preface, and a Feast Day with a proper Preface is commemorated, the Commemoration's Preface may be said, the Feasts of Our Lord always receiving priority.\par
    \textsc{Note,} When no proper Preface is said, the Priest proceeds immediately from \emph{Almighty, Everlasting God.} to \emph{Therefore, with Angels and Archangels}.
        %\begin{description}
            %\item[Advent] Every day in Advent, including Christmas Eve.
            %\item[Christmas] Christmas Day and its Octave, until Epiphany Eve, inclusive, and on the Feast of Corpus Christi and its Octave, including the Sunday within the Octave.
            %\item[Epiphany] Epiphany Day and its Octave.
            %\item[Septuagesima Sunday] Septuagesima Sunday.
            %\item[Sexagesima Sunday] Sexagesima Sunday.
            %\item[Quinquagesima Sunday] Quinquagesima Sunday.
            %\item[Lent] From Ash Wednesday until the Saturday before Passion Sunday, inclusive, unless the Mass be of the Lord and require another preface.
            %\item[Holy Rood] From Passion Sunday until Maundy Thursday, inclusive. And all Masses of the Holy Rood, the Passion of the Lord, or the Most Precious Blood, and within their Octaves, even when celebrated within the Octave of the Nativity of Our Lord.
            %\item[Easter] From Holy Saturday until Ascension Eve, inclusive.
            %\item[Ascension] From Ascension Thursday until the Friday after its Octave, inclusive. It is always said, even without a Commemoration, on the Friday after the Octave.
            %\item[Holy Ghost] From the Whitsunday Eve until Trinity Sunday, exclusive, in Votive Masses of the Holy Ghost (even when celebrated within the Octave of the Nativity of Our Lord).
            %\item[Trinity] On Sundays, from the \nth{3} Sunday after Trinity until the Sunday Next before Advent, inclusive. And also in Masses of the Most Holy Trinity, even when celebrated within the Octave of the Nativity of the Lord.
            %\item[Divine Compassion] Feast of the of the Compassion of Our Lord and its Octave Day.
            %\item[Of Our Lord] Feasts of the Purification, Annunciation, and Transfiguration.
            %\item[Blessed Virgin Mary] In Masses of the Blessed Virgin Mary, and their Octaves.
            %\item[St. Joseph] In Masses of St. Joseph, and their Octaves.
            %\item[Apostles \& Evangelists] In Masses of the Apostles or Evangelists and in a Mass of the Election and Enthronement of the Patriarch and their anniversaries.
            %\item[Requiem] All Masses of the Dead.
            %\item[Christ the King] Masses of Our Lord Jesus Christ the King, even when they are celebrated within the Octave of the Nativity of Our Lord.
            %\item[All Hallows] Feast of All Hallows and its Octave.
            %\item[Common] When no proper Preface is said, the Priest proceeds immediately from \emph{Almighty, Everlasting God.} to \emph{Therefore, with Angels and Archangels}.
            %\end{description}
    \item[Reciting the Propers] When a clerk of appropriate rank is not available, a clerk of higher rank shall fulfil his role, such as the Priest chanting the Gospel or a Deacon chanting the Epistle.
    %Translation of the following rubric is from the English Missal:
    \item[Requiem Mass] On the first day of each month (outside Advent, Lent, and Eastertide) not hindered by a Semidouble or higher Feast Day, the Mass is said generally for departed Priests, Benefactors, and others. But if no day occur in the week on which it can be said, a general commemoration of the departed shall be said, except during Lent and Eastertide.\par
    \textsc{Note,} The \emph{Dies irae} is said on All Souls' Day and on the day of the burial in all sung Masses. However, in other Masses, it may either be recited or omitted at the pleasure of the celebrant.
    \item[Seasonal Prayers] Throughout the year, as indicated in the rubrics, additional Collects, Secrets, and Postcommunions are to be said, assigned for the different seasons even within Octaves or on Vigils. On Marian Feast Days (or Feast Days which include her commemoration), when the second seasonal collect would be of St. Mary, it is instead of the Holy Ghost.\par
    \textsc{Note,} The propers chosen for the Collect should match for the Secret and Postcommunion.
    \item[Signs of the Cross] In the Mass, when {\ding{64}} appears, unless otherwise indicated or manifestly evident, it indicates the manual act of making the sign of the Cross.\par
    In the Canon, if there be a cross ({\ding{64}}) in reference to Bread or Body, it indicates a sign of the cross made over the Host. Likewise, if there be a cross ({\ding{64}}) in reference to Wine or Cup, it indicates a sign of the cross made over the Cup. Otherwise, it indicates a sign of the cross made over the Gifts generally, unless otherwise stated.\par
    When the rubrics tell the Priest to make the sign of the Cross on himself, unless otherwise indicated or manifestly evident, it indicates making the sign of the Cross from forehead to breast.
    \item[Votive Mass of Our Lady on Saturday] On Saturdays not hindered by a Double or Semidouble Feast, Octave, Vigil, Lenten Feria, or Ember Day, or by the Office of a Sunday anticipated, Mass is said of our Lady according to the season, as given in the Missal.\par
    %In Advent, however, if the office of our Lady be not said on Saturday, nevertheless the principal Mass is of her, with commemoration of Advent, unless it be an Ember Day or Vigil, as above.
\end{description}
\newpage
\phantomsection
\addcontentsline{toc}{chapter}{Rankings of Days}\label{ranking}
\fancyhead[C]{{\Large Ranking of Days}}
%\section{Ranking of Days}\label{ranking}
\noindent
\lett{F}{east} Days are ranked in order of priority. This is important, since it determines both what is prayed in the changeable parts of the Office (the Propers) and how to determine which Feast outranks the other(s) when Feast Days overlap on the same day.\par
The rankings for all of the Feast Days throughout the year are provided in the Kalendar (p. \pageref{kalendar}). However, there are some rankings which are not provided there or may need to be explained. For most people, it is expedient to order a yearly Ordo from the Antiochian Western Rite Vicariate which lays out the correct ordering for each day. However, to determine it without an Ordo, important rankings are described here.
\begin{itemize}
    \item Feast Days
        \begin{description}
            \item[First Class Double] (I Double)
            \item[Second Class Double] (II Double)
            \item[Greater Double]
            \item[Double] Feasts of Double and higher have I and II Evensong, beginning with I Evensong and ending with Compline of the following day.
            \item[Semidouble] Sundays and days within Common Octaves are Semidouble.\par
            The Semidouble Office begins with I Evensong and ends with Compline of the following day.
            \item[Memorial] No Office is to be said for a Memorial, but on the day on which they are noted in the Kalendar, Commemoration only is made of them at I Evensong and Mattins; except on II Doubles, on which Commemoration of the Memorial is omitted at Evensong; and on I Doubles, on which no Commemoration is made of Memorials.
            \item[Simple] The Office is Simple on weekdays when the Office is of the Feria, and on the Octave Day of a Simple Octave.\par
                The Simple Office begins with I Evensong, and ends with None of the following day.
        \end{description}
    \item Ferial Days
        \begin{description}
            \item[First Class Feria] No Feast Day may be celebrated. Begins at Mattins.
            %Not sure if this is true of the English Office:
                %The Office of the three Greater Ferias of Holy Week is said as directed in the Proper of the Season.
                \begin{itemize}
                    \item Ash Wednesday
                    \item Holy Monday
                    \item Holy Tuesday
                    \item Holy Wednesday
                \end{itemize}
            \item[Second Class Feria] Only I and II Doubles may be celebrated, lesser Feast Days being commemorated. Begins at Mattins.
            	\begin{itemize}
            		\item Ember Days
            		\item Rogation Monday
            		\item Weekdays in Lent \& Passion Week
            	\end{itemize}
            \item[Greater Feria] Begins at Mattins. 
                \begin{itemize}
                    \item Weekdays in Advent
                \end{itemize}
            \item[Feria] Begins when the Office of the preceding day ceases. The Ferial Office ends at None if a Double or Simple follow.
        \end{description}
    \item Octaves
        \begin{itemize}
            \item Privileged Octaves
                \begin{description}
                    \item[First Order Privileged Octaves] The Feast Day, and second \& third days of the Octave, are I Double. The rest of the Octave until the Octave Day is Semidouble. The Octave Day is I Double. Feast Days cannot outrank the Days within the Octave, nor the Octave Day.\par
                    A I or II Double is transferred to the first unhindered day after the Octave; other Feasts are commemorated, except on the Feast Day and the two following days.
                        \begin{itemize}
                            \item Easter Octave
                            \item Whitsun Octave
                        \end{itemize}
                    \item[Second Order Privileged Octaves] The Feast Day is I Double. The Days within the Octave are Semidouble, being only outranked by I Double, with the Octave commemorated. The Octave Day is Greater Double, and gives place only to a I Double of universal observance.
                        \begin{itemize}
                            \item Epiphany Octave
                            \item Corpus Christi Octave
                        \end{itemize}
                    \item[Third Order Privileged Octaves] Days within the Octave are trumped by any feast over Simple. Any Double Feast within this Octave has its Office with Commemoration of the Octave, but on the Octave Day of the Ascension only a I or II Double is kept, with Commemoration of the Octave.
                        \begin{itemize}
                            \item Christmas Octave
                            \item Ascension Octave
                        \end{itemize}
                \end{description}
            \item Common Octaves\par
            The Days within the Octave are Semidouble. A Double Feast occurring within a Common Octave has its Office with Commemoration of the Octave, unless otherwise noted.\par
            \textsc{Note,} Common Octaves are not commemorated during Lent.
                \begin{itemize}
                    \item Conception of the BVM Octave
                    \item Assumption Octave
                    \item Nativity of St. John Baptist Octave
                    \item Solemnity of St. Joseph Octave
                    \item Sts. Peter \& Paul Octave
                    \item All Hallows Octave
                    \item The Octave of the principal patron saint of a church, cathedral, order, town, diocese, province, or nation. 
                \end{itemize}
            \item Simple Octaves\par
            The Feast Day is II Double. The Octave Day is Simple. The Days within the Octave are not commemorated.
                \begin{itemize}
                    \item St. Stephen Octave
                    \item St. John the Evangelist Octave
                    \item Holy Innocents Octave
                    \item St. Lawrence Octave
                    \item Nativity of the B.V.M. Octave
                    \item Divine Compassion Octave
                    \item An Octave of Secondary Patrons
                \end{itemize}
        \end{itemize}
    \item Sundays
        \begin{description}
            \item[First Class] Sundays of the First Class give place to no Feast. A I or II Double Feast Day falling on a Sunday of the First Class is transferred to the first unhindered day.\par
            \textsc{Note,} When I or II Double Feast Day falls on (or is transferred to) Monday, I Evensong is of the Feast Day with a commemoration of the Sunday. Otherwise, it is of the Sunday with a commemoration of the Feast Day.
                \begin{itemize}
                    \item First Sunday of Advent
                    \item Sundays in Lent \& Passiontide
                    \item Easter Sunday
                    \item Low Sunday
                    \item Whitsunday
                \end{itemize}
            \item[Second Class] Sundays of the Second Class give place only to I Doubles, and the Sunday is commemorated.
                \begin{itemize}
                    \item Second Sunday of Advent
                    \item Third Sunday of Advent
                    \item Fourth Sunday of Advent
                    \item Septuagesima Sunday
                    \item Sexagesima Sunday
                    \item Quinquagesima Sunday
                \end{itemize}
            \item[All Other Sundays] On lesser Sundays, the Office is of the Sunday unless a Double I or II Class occur thereon, when the Office is of the Feast with Commemoration of the Sunday.
        \end{description}
    \item Privileged Vigils\par
        \begin{description}
            \item[First Class] Vigils of the First Class are preferred over any other Feast Day.
                \begin{itemize}
                    \item Vigil of the Nativity of Our Lord
                    \item Vigil of Whitsunday
                \end{itemize}
            \item[Second Class] Vigils of the Second Class are preferred over any Feast Day, except I \& II Doubles and Feasts of Our Lord.
                \begin{itemize}
                    \item Vigil of Epiphany
                \end{itemize}
        \end{description}
    \item Common Vigils\par
        Non-privileged Vigils are preferred over any Feast Day, except Doubles \& above or an Octave.
\end{itemize}

%\subsection{First \& Second Evensong}
%When I Evensong of a Feast Day overlaps with II Evensong of another Feast Day, one will always take priority.
%
%  	\begin{description}
 % 		\item[First Class Double with a Common Octave] I Evensong is always said.
 % 		\item[First Class Double] I Evensong is said, unless the II Evensong is of a First Class Double with a Common Octave. If the II Evensong Feast Day be of Greater Double or higher, Commemoration is made of it.
%  		\item[Second Class Double] I Evensong is said, unless the II Evensong be of a Second Class Double.\par
%  	\end{description}

\include{pages/Front/Kalendar}
%Lectionary
\cleardoublepage
\fancyhead[RE,LO]{}\fancyhead[RO,LE]{}
\fancyhead[C]{}
\thispagestyle{empty}
\phantomsection
\addcontentsline{toc}{chapter}{Daily Office Lectionary}
%  \begin{center}
%   \includegraphics[width=98mm, height=140mm]{Coronation.eps}
%   \par
%   \vspace{2ex}
%   	\textsc{\Huge{Daily Office Lectionary}}
%   \end{center}
   
\begin{figure}[H]
    \centering
    \includegraphics[width=98mm, height=140mm]{Lectionary/Coronation.eps}
    \caption{\textsc{\Huge{Daily Office Lectionary}}}
\end{figure}
\include{pages/Lectionary/Lectionary}
\include{pages/Lectionary/AlternativePsalter}
%Daily Office
\cleardoublepage
\fancyhead[RE,LO]{}\fancyhead[RO,LE]{}
\fancyhead[C]{}\thispagestyle{empty}
\mainmatter
\phantomsection
\addcontentsline{toc}{chapter}{Daily Office}

\begin{figure}[H]
    \centering
    \includegraphics[width=98mm, height=140mm]{Office/Resurrection.eps}
    \caption{\textsc{\Huge{Daily Office}}}
\end{figure}
\include{pages/Office/ForeOffice}
\include{pages/Office/Mattins}
\include{pages/Office/OTCant}
\fancyhead[RO,LE]{\textit{Prime}}
\fancyhead[RE,LO]{}
\section{Hour of Prime}
\begin{secrubric}
	This service may be said at any convenient time before or after the saying of Morning Prayer, but not in substitution therefor.
\end{secrubric}

\begin{multicols}{2}

℣. O God, {\ding{64}} make speed to save us.

℟. O Lord, make haste to help us.

℣. Glory be to the Father, and to the Son, and to the Holy Ghost.

℟. As it was in the beginning, is now, and ever shall be, world without end. Amen.

℣. Praise ye the Lord.

℟. The Lord's Name be praised.

\subsection{Jam Lucis Ordo Sidere}
\begin{hangparas}{1.25em}{1}
\lett{N}{ow} that the daylight fills the sky,\par
\secondline{We lift our hearts to God on high,}\par\noindent
\thirdline{That he, in all we do or say,}\par\noindent
Would keep us free from harm to-day:\\
\par\noindent
Would guard our hearts and tongues from strife;\par\noindent
From anger's din would hide our life;\par\noindent
From all ill sights would turn our eyes,\par\noindent
Would close our ears from vanities:\\
\par\noindent
Would keep our inmost conscience pure;\par\noindent
Our souls from folly would secure;\par\noindent
Would bid us check the pride of sense\par\noindent
With due and holy abstinence.\\
\par\noindent
So we, when this new day is gone,\par\noindent
And night in turn is drawing on,\par\noindent
With conscience by the world unstained\par\noindent
Shall praise his name for victory gained.\\
\par\noindent
All laud to God the Father be;\par\noindent
All praise, eternal Son, to thee;\par\noindent
All glory, as is ever meet,\par\noindent
To God the holy Paraclete. Amen.\\
\end{hangparas}

\subsection{Psalm 54. \textit{Deus, in nomine}}
\lett{S}{ave} me, O God, for thy Name's sake : and avenge me in thy strength.\par
\secondline{\psanum{2}Hear my prayer, O God : and hearken unto the words of my mouth.}
\thirdline{\psanum{3}For strangers are risen up against me : and tyrants, which have not God before their eyes, seek after my soul.}
\psanum{4}Behold, God is my helper : the Lord is with them that uphold my soul.\par
\psanum{5}He shall reward evil unto mine enemies : destroy thou them in thy truth.\par
\psanum{6}An offering of a free heart will I give thee, and praise thy Name, O Lord : because it is so comfortable.\par
\psanum{7}For he hath delivered me out of all my trouble : and mine eye hath seen his desire upon mine enemies.\par
℣. Glory be to the Father, and to the Son, and to the Holy Ghost.\par
℟. As it was in the beginning, is now, and ever shall be, world without end. Amen.

\subsection{Psalm 119. \textit{Beati immaculati}}
\lett{B}{lessed} are those that are undefiled in the way : and walk in the law of the Lord.\par
\secondline{\psanum{2}Blessed are they that keep his testimonies : and seek him with their whole heart.}
\thirdline{\psanum{3}For they who do no wickedness : walk in his ways.}
\psanum{4}Thou hast charged : that we shall diligently keep thy commandments.\par
\psanum{5}O that my ways were made so direct : that I might keep thy statutes!\par
\psanum{6}So shall I not be confounded : while I have respect unto all thy commandments.\par
\psanum{7}I will thank thee with an unfeigned heart : when I shall have learned the judgements of thy righteousness.\par
\psanum{8}I will keep thy ceremonies : O forsake me not utterly.\par
℣. Glory be to the Father, and to the Son, and to the Holy Ghost.\par
℟. As it was in the beginning, is now, and ever shall be, world without end. Amen.

\subsection{In quo corriget?}
\lett{W}{herewithal} shall a young man cleanse his way : even by ruling himself after thy word.\par
\secondline{\psanum{10}With my whole heart have I sought thee : O let me not go wrong out of thy commandments.}
\thirdline{\psanum{11}Thy words have I hid within my heart : that I should not sin against thee.}
\psanum{12}Blessed art thou, O Lord : O teach me thy statutes.\par
\psanum{13}With my lips have I been telling : of all the judgements of thy mouth.\par
\psanum{14}I have had as great delight in the way of thy testimonies : as in all manner of riches.\par
\psanum{15}I will talk of thy commandments : and have respect unto thy ways.\par
\psanum{16}My delight shall be in thy statutes : and I will not forget thy word.\par
℣. Glory be to the Father, and to the Son, and to the Holy Ghost.\par
℟. As it was in the beginning, is now, and ever shall be, world without end. Amen.

\subsection{Retribue servo tuo}
\lett{O}{do} well unto thy servant : that I may live, and keep thy word.\par
\secondline{\psanum{18}Open thou mine eyes : that I may see the wondrous things of thy law.}
\thirdline{\psanum{19}I am a stranger upon earth : O hide not thy commandments from me.}
\psanum{20}My soul breaketh out for the very fervent desire : that it hath alway unto thy judgements.\par
\psanum{21}Thou hast rebuked the proud : and cursed are they that do err from thy commandments.\par
\psanum{22}O turn from me shame and rebuke : for I have kept thy testimonies.\par
\psanum{23}Princes also did sit and speak against me : but thy servant is occupied in thy statutes.\par
\psanum{24}For thy testimonies are my delight : and my counsellors.\par
℣. Glory be to the Father, and to the Son, and to the Holy Ghost.\par
℟. As it was in the beginning, is now, and ever shall be, world without end. Amen.

\subsection{Adhaesit pavimento}
\lett{M}{y} soul cleaveth to the dust : O quicken thou me, according to thy word.\par
\secondline{\psanum{26}I have acknowledged my ways, and thou heardest me : O teach me thy statutes.}
\thirdline{\psanum{27}Make me to understand the way of thy commandments : and so shall I talk of thy wondrous works.}
\psanum{28}My soul melteth away for very heaviness : comfort thou me according unto thy word.\par
\psanum{29}Take from me the way of lying : and cause thou me to make much of thy law.\par
\psanum{30}I have chosen the way of truth : and thy judgements have I laid before me.\par
\psanum{31}I have stuck unto thy testimonies : O Lord, confound me not.\par
\psanum{32}I will run the way of thy commandments : when thou hast set my heart at liberty.\par
℣. Glory be to the Father, and to the Son, and to the Holy Ghost.\par
℟. As it was in the beginning, is now, and ever shall be, world without end. Amen.

\begin{rubric}
    As in the ancient use of Prime, the Athanasian Creed (p. \pageref{Ath}) may here be said, concluding with the \blackrubric{Gloria Patri.}
\end{rubric}

\subsection{Chapter}

%From Sarum:
\begin{inhead}
	Ferias Throughout the Year
\end{inhead}\par\noindent
Love the truth and peace, saith the Lord of hosts.

\begin{inhead}
	Sundays and Feast Days
\end{inhead}\par\noindent
Now unto the King eternal, immortal, invisible, the only wise God, be honour and glory for ever and ever.

\begin{inhead}
	Eastertide
\end{inhead}\par\noindent
O Lord, be gracious unto us; we have waited for thee: be thou our arm every morning, our salvation also in the time of trouble.

\begin{rubric}
	After the Chapter is said,
\end{rubric}
℟. Thanks be to God.

\subsection{Short Respond}
℣. Jesu Christ, Son of the living God, have mercy upon us;

℟. Jesu Christ, Son of the living God, have mercy upon us.

℣. Thou that sittest at the right hand of the Father;

℟. Have mercy upon us.

℣. Glory be to the Father, and to the Son, and to the Holy Ghost.

℟. Jesu Christ, Son of the living God, have mercy upon us.\\

℣. O Lord, arise, help us;

℟. And deliver us for thy name’s sake.

%MANUAL ADJUSTMENT:
\newcolumn

\subsection{Suffrages}
\begin{rubric}
	The Suffrages are said, kneeling, on all Ferias outside of Eastertide.
\end{rubric}
℣. Lord, have mercy upon us.

℟. Christ, have mercy upon us.

℣. Lord, have mercy upon us.

\subsubsection{Lord's Prayer}
\lett{O}{ur} Father, who art in heaven, Hallowed be thy Name. Thy kingdom come. Thy will be done on earth, As it is in heaven. Give us this day our daily bread. And forgive us our trespasses, As we forgive those who trespass against us. And lead us not into temptation; But deliver us from evil. Amen.

\begin{rubric}
	Then shall follow these Versicles and Responses.
\end{rubric}

℣. O let my mouth be filled with thy praise;

℟. That I may sing of thy glory and honour all the day long.

℣. Turn thy face from my sins, O Lord;

℟. And put out all my misdeeds.

℣. Make me a clean heart, O God;

℟. And renew a right spirit within me.

℣. Cast me not away from thy presence;

℟. And take not thy Holy Spirit from me.

℣. O give me the comfort of thy help again;

℟. And stablish me with thy free spirit.

\subsubsection{Confession}
\begin{rubric}
	Then shall the Minister and People say together the Confession following,
\end{rubric}
%Replaced with Sarum. Text from the Monastic Diurnal.
\lett{I}{confess} to God, to blessed Mary, to all the saints, and to you: that I have sinned exceedingly in thought, word, and deed, by my own fault. I beg holy Mary, all the Saints of God, and you, to pray for me.

\begin{rubric}
	The Minister then stands and says,
\end{rubric}

%1928 English Translation
\lett{A}{lmighty} God, have mercy upon us, forgive us all our sins and deliver us from all evil, confirm and strengthen us in all goodness, and bring us to life everlasting. \textit{Amen.}

\begin{rubric}
	If a Priest be present, he shall stand and pronounce the following Absolution,
\end{rubric}

%1928 English Translation
\lett{M}{ay} the Almighty and Merciful Lord grant unto you pardon {\ding{64}} and remission of all your sins, time for amendment of life, and the grace and comfort of the Holy Spirit. \textit{Amen.}

℣. Wilt thou not turn again and quicken us, O God?

℟. That thy people may rejoice in thee.

℣. O Lord, show thy mercy upon us.

℟. And grant us thy salvation.

℣. Vouchsafe, O Lord, to keep us this day without sin;

℟. O Lord, have mercy upon us, have mercy upon us.

\subsection{Collect}
℣. O Lord, hear our prayer.

℟. And let our cry come unto thee.

\letuspray
%From Sarum:
\lett{I}{n} this hour of this day, fill us, O Lord, with thy mercy; that going forth in thy strength, we may make our boast of thee all the day long. Through.

℣. The Lord be with you.\par
\textit{or,} O Lord, hear our prayer.

℟. And with thy spirit.\par
\textit{or,} And let our cry come unto thee.

℣. Let us bless the Lord.

℟. Thanks be to God.

\lett{T}{he} Lord bless us, and preserve us from all evil, and bring us to everlasting life; and may the souls {\ding{64}} of the faithful departed, through the mercy of God, rest in peace. \textit{Amen.}
\end{multicols}

\begin{rubric}
	The Office of Prime may end here, or the following devotions may be added.
\end{rubric}
\section{Pretiosa}
\fancyhead[RO,LE]{\textit{Pretiosa}}
\fancyhead[RE,LO]{}
\begin{secrubric}
    The \blackrubric{Pretiosa} is a laudable custom and may be added to \blackrubric{Prime.} The Martyrology reading for the Day may here be read. The following is then said.
\end{secrubric}
\begin{multicols}{2}
℣. Right dear in the sight of the Lord.

℟. Is the death of his Saints.\\

\lett{M}{ay} holy Mary and all the Saints intercede for us to the Lord: that we may be worthy to obtain from him help and salvation, who liveth and reigneth for ever and ever. \textit{Amen.}\\

℣. O God, {\ding{64}} make speed to save me.

℟. O Lord, make haste to help me.

℣. O God, {\ding{64}} make speed to save me.

℟. O Lord, make haste to help me.

℣. O God, {\ding{64}} make speed to save me.

℟. O Lord, make haste to help me.\\

℣. Glory be to the Father, and to the Son, and to the Holy Ghost.

℟. As it was in the beginning, is now, and ever shall be, world without end. Amen.

℣. Lord, have mercy upon us.

℟. Christ, have mercy upon us.

℣. Lord, have mercy upon us.\\

\lett{O}{ur} Father, who art in heaven, Hallowed be thy Name. Thy kingdom come. Thy will be done on earth, As it is in heaven. Give us this day our daily bread. And forgive us our trespasses, As we forgive those who trespass against us. And lead us not into temptation; But deliver us from evil. Amen.\\

℣. Let thy loving mercy come also unto us, O Lord.

℟. Even thy salvation, according unto thy word.

℣. And show thy servants thy work.

℟. And their children thy glory.

℣. And the glorious Majesty of the Lord our God be upon us.

℟. And prosper thou the work of our hands upon us: O prosper thou our handy-work.\\

\letuspray

\lett{O}{Almighty} Lord and Everlasting God, vouchsafe, we beseech thee, to direct, sanctify, and govern, both our hearts and bodies in the ways of thy laws and in the works of thy commandments: that through thy most mighty protection, both here and ever, we may be preserved in body and soul; through our Lord and Saviour Jesus Christ. \textit{Amen.}\\

℣. May the Divine Assistance remain always with us.

℟. Amen.
\end{multicols}
\include{pages/Office/CommemorationDeparted}
\include{pages/Office/Athanasian}
\include{pages/Office/Terce}
\fancyhead[RO,LE]{\textit{Sext}}
\fancyhead[RE,LO]{}
\section{Hour of Sext}

\begin{multicols}{2}

℣. O God, {\ding{64}} make speed to save us.

℟. O Lord, make haste to help us.

℣. Glory be to the Father, and to the Son, and to the Holy Ghost.

℟. As it was in the beginning, is now, and ever shall be, world without end. Amen.

℣. Praise ye the Lord.

℟. The Lord's Name be praised.

\subsection{Rector Potens}
\begin{hangparas}{1.25em}{1}
\lett{O}{God} of truth, O Lord of might,\par
\secondline{Who ord'rest time and change aright,}
\par\noindent
\thirdline{And send'st the early morning ray,}\par\noindent
And light'st the glow of perfect day.\\
\par\noindent
Extinguish thou each sinful fire,\par\noindent
And banish every ill desire:\par\noindent
And while thou keep'st the body whole,\par\noindent
Shed forth thy peace upon the soul.\\
\par\noindent
O Father, that we ask be done,\par\noindent
Through Jesus Christ, thine only Son;\par\noindent
Who, with the Holy Ghost and thee,\par\noindent
Doth live and reign eternally. Amen.\\
\end{hangparas}

\subsection{Psalm 123. \textit{Ad te levavi oculos meos}}
\lett{U}{nto} thee lift I up mine eyes : O thou that dwellest in the heavens.\par
\secondline{\psanum{2}Behold, even as the eyes of servants look unto the hand of their masters, and as the eyes of a maiden unto the hand of her mistress : even}
\noindent
so our eyes wait upon the Lord our God, until he have mercy upon us.\par
\psanum{3}Have mercy upon us, O Lord, have mercy upon us : for we are utterly despised.\par
\psanum{4}Our soul is filled with the scornful reproof of the wealthy : and with the despitefulness of the proud.\par
℣. Glory be to the Father, and to the Son, and to the Holy Ghost.\par
℟. As it was in the beginning, is now, and ever shall be, world without end. Amen.

\subsection{Psalm 124. \textit{Nisi quia Dominus}}
\lett{I}{f} the Lord himself had not been on our side, now may Israel say : if the Lord himself had not been on our side, when men rose up against us;\par
\secondline{\psanum{2}They had swallowed us up quick : when they were so wrathfully displeased at us.}
\thirdline{\psanum{3}Yea, the waters had drowned us : and the stream had gone over our soul.}
\psanum{4}The deep waters of the proud : had gone even over our soul.\par
\psanum{5}But praised be the Lord : who hath not given us over for a prey unto their teeth.\par
\psanum{6}Our soul is escaped even as a bird out of the snare of the fowler : the snare is broken, and we are delivered.\par
\psanum{7}Our help standeth in the Name of the Lord : who hath made heaven and earth.\par
℣. Glory be to the Father, and to the Son, and to the Holy Ghost.\par
℟. As it was in the beginning, is now, and ever shall be, world without end. Amen.

\subsection{Psalm 125. \textit{Qui confidunt}}
\lett{T}{hey} that put their trust in the Lord shall be even as the mount Sion : which may not be removed, but standeth fast for ever.\par
\secondline{\psanum{2}The hills stand about Jerusalem : even so standeth the Lord round about his people, from this time forth for evermore.}
\thirdline{\psanum{3}For the rod of the ungodly cometh not into the lot of the righteous : lest the righteous put their hand unto wickedness.}
\psanum{4}Do well, O Lord : unto those that are good and true of heart.\par
\psanum{5}As for such as turn back unto their own wickedness : the Lord shall lead them forth with the evil-doers; but peace shall be upon Israel.\par
℣. Glory be to the Father, and to the Son, and to the Holy Ghost.\par
℟. As it was in the beginning, is now, and ever shall be, world without end. Amen.

\newcolumn

\subsection{Chapter}
\begin{inhead}
	Ferias Throughout the Year
\end{inhead}\par\noindent
Prove all things, hold fast to that which is good: abstain from all appearance of evil.

\begin{inhead}
	Sundays \& Feast Days
\end{inhead}\par\noindent
There are three that bear record in heaven: The Father, the Word, and the Holy Ghost: and these three are one.

\begin{inhead}
	Eastertide
\end{inhead}\par\noindent
Who is he that overcometh the world, but he that believeth that Jesus is the Son of God? This is he that came by water and blood, even Jesus Christ.

℟. Thanks be to God.

\subsection{Short Respond}
\begin{inhead}
	Ferias Throughout the Year
\end{inhead}
℣. I will bless the Lord at all times.

℟. I will bless the Lord at all times.

℣. His praise shall ever be in my mouth.

℟. At all times.

℣. Glory be to the Father, and to the Son, and to the Holy Ghost.

℟. I will bless the Lord at all times.\\

℣. Behold, God is my helper.

℟. The Lord is with them that uphold my soul.

\begin{inhead}
	Sundays, Feasts, \& in Eastertide
\end{inhead}

℣. O Lord, thy Word endureth for ever in heaven (alleluia, alleluia).

℟. O Lord, thy Word endureth for ever in heaven (alleluia, alleluia).

℣. Thy truth also remaineth from one generation to another.

℟. For ever in heaven (alleluia, alleluia).

℣. Glory be to the Father, and to the Son, and to the Holy Ghost.

℟. O Lord, thy Word endureth for ever in heaven (alleluia, alleluia).

℣. The Lord is my shepherd, therefore can I lack nothing (alleluia).

℟. He shall feed me in a green pasture (alleluia).

\subsection{Suffrages}
\begin{rubric}
	The Suffrages are said, kneeling, on all Ferias outside of Eastertide.
\end{rubric}

℣. Lord, have mercy upon us.

℟. Christ, have mercy upon us.

℣. Lord, have mercy upon us.

\lett{O}{ur} Father, who art in heaven, Hallowed be thy Name. Thy kingdom come. Thy will be done on earth, As it is in heaven. Give us this day our daily bread. And forgive us our trespasses, As we forgive those who trespass against us. And lead us not into temptation; But deliver us from evil. Amen.

℣. Hide not thy face from thy servant, for I am in trouble.

℟. O haste thee and hear me.

℣. Draw nigh unto my soul and save it.

℟. O deliver me, because of mine enemies.

℣. Arise, O Christ, and help us.

℟. And deliver us for thy Name's sake.

\subsection{Collects}
℣. The Lord be with you.\par
\textit{or,} O Lord, hear our prayer.

℟. And with thy spirit.\par
\textit{or,} And let our cry come unto thee.

℣. Let us pray.

\begin{rubric}
	The Collect(s) of the Day is here said, followed by this memorial Collect.
\end{rubric}

\lett{L}{ord} Jesus, who didst stretch out thine arms of love on the hard wood of the Cross, that all men might come within the reach of thy saving embrace; clothe us in thy Spirit, that we, stretching forth our hands in loving labour for others, may bring those who know thee not to the knowledge and love of thee, who with the Father and the Holy Ghost livest and reignest ever one God, world without end. \textit{Amen.}\\

℣. The Lord be with you.\par
\textit{or,} O Lord, hear our prayer.

℟. And with thy spirit.\par
\textit{or,} And let our cry come unto thee.

℣. Let us bless the Lord.

℟. Thanks be to God.

℣. May the souls {\ding{64}} of the faithful departed, through the mercy of God, rest in peace.

\lett{N}{ow} unto the King eternal, immortal, invisible, the only wise God, be honour and glory, for ever and ever. \textit{Amen.}

\end{multicols}

%\begin{center}
%	{\textsc{Here endeth the Order of Sext.}}
%\end{center}
\fancyhead[RO,LE]{\textit{None}}
\fancyhead[RE,LO]{}
\section{Hour of None}

\begin{multicols}{2}

℣. O God, {\ding{64}} make speed to save us.

℟. O Lord, make haste to help us.

℣. Glory be to the Father, and to the Son, and to the Holy Ghost.

℟. As it was in the beginning, is now, and ever shall be, world without end. Amen.

℣. Praise ye the Lord.

℟. The Lord's Name be praised.

\subsection{Rerum Deus}
%MANUAL ADJUSTMENT:
\lett{O}{God,} creation’s secret force,\par
\leftskip=3.25em
\par\noindent
Thyself unmoved, all motion's source,

\leftskip=0em

\begin{hangparas}{1.25em}{1}
\par\noindent
Who from the morn till evening ray,\par\noindent
Through all its changes guid'st the day:\\
\par\noindent
Grant us, when this short life is past,\par\noindent
The glorious evening that shall last;\par\noindent
That, by a holy death attained,\par\noindent
Eternal glory may be gained.\\
\par\noindent
O Father, that we ask be done,\par\noindent
Through Jesus Christ, thine only Son;\par\noindent
Who, with the Holy Ghost and thee,\par\noindent
Doth live and reign eternally. Amen.
\end{hangparas}

\subsection{Psalm 126. \textit{In convertendo}}

\lett{W}{hen} the Lord turned again the captivity of Sion : then were we like unto them that dream.\par
\secondline{\psanum{2}Then was our mouth filled with laughter : and our tongue with joy.}
\thirdline{\psanum{3}Then said they among the heathen : The Lord hath done great things for them.}
\psanum{4}Yea, the Lord hath done great things for us already : whereof we rejoice.\par
\psanum{5}Turn our captivity, O Lord : as the rivers in the south.\par
\psanum{6}They that sow in tears : shall reap in joy.\par
\psanum{7}He that now goeth on his way weeping, and beareth forth good seed : shall doubtless come again with joy, and bring his sheaves with him.\par
℣. Glory be to the Father, and to the Son, and to the Holy Ghost.\par
℟. As it was in the beginning, is now, and ever shall be, world without end. Amen.

\subsection{Psalm 127. \textit{Nisi Dominus}}

\lett{E}{xcept} the Lord build the house : their labour is but lost that build it.\par
\secondline{\psanum{2}Except the Lord keep the city : the watchman waketh but in vain.}
\thirdline{\psanum{3}It is but lost labour that ye haste to rise up early, and so late take rest, and eat the bread of carefulness : for so he giveth his beloved sleep.}
\psanum{4}Lo, children and the fruit of the womb : are an heritage and gift that cometh of the Lord.\par
\psanum{5}Like as the arrows in the hand of the giant : even so are the young children.\par
\psanum{6}Happy is the man that hath his quiver full of them : they shall not be ashamed when they speak with their enemies in the gate.\par
℣. Glory be to the Father, and to the Son, and to the Holy Ghost.\par
℟. As it was in the beginning, is now, and ever shall be, world without end. Amen.

\subsection{Psalm 128. \textit{Beati omnes}}

\lett{B}{lessed} are all they that fear the Lord : and walk in his ways.\par
\secondline{\psanum{2}For thou shalt eat the labours of thine hands : O well is thee, and happy shalt thou be.}
\thirdline{\psanum{3}Thy wife shall be as the fruitful vine : upon the walls of thine house.}
\psanum{4}Thy children like the olive-branches : round about thy table.\par
\psanum{5}Lo, thus shall the man be blessed : that feareth the Lord.\par
\psanum{6}The Lord from out of Sion shall so bless thee : that thou shalt see Jerusalem in prosperity all thy life long.\par
\psanum{7}Yea, that thou shalt see thy children's children : and peace upon Israel.\par
℣. Glory be to the Father, and to the Son, and to the Holy Ghost.\par
℟. As it was in the beginning, is now, and ever shall be, world without end. Amen.

\subsection{Chapter}
\begin{inhead}
	Ferias Throughout the Year
\end{inhead}\par\noindent
Bear ye one another's burdens, and so fulfil the law of Christ.

\begin{inhead}
	Sundays \& Feast Days
\end{inhead}\par\noindent
One Lord, one Faith, one Baptism, one God and Father of all, who is above all, and through all, and in you all.

\begin{inhead}
	Eastertide
\end{inhead}\par\noindent
There are three that bear witness in earth: the Spirit, and the Water, and the Blood: and these three agree in one.

℟. Thanks be to God.

\subsection{Short Respond}
\begin{inhead}
	Ferias Throughout the Year
\end{inhead}
℣. Deliver me, O Lord, and be merciful unto me.

℟. Deliver me, O Lord, and be merciful unto me.

℣. My foot standeth right; I will praise the Lord in the congregations.

℟. And be merciful unto me.

℣. Glory be to the Father, and to the Son, and to the Holy Ghost.
    
℟. Deliver me, O Lord, and be merciful unto me.

℣. I will praise thy Name, O Lord.

℟. For thou hast delivered me out of all my trouble.

\begin{inhead}
	Sundays, Feasts, \& in Eastertide
\end{inhead}

℣. I call with my whole heart; hear me, O Lord (alleluia, alleluia).

℟. I call with my whole heart; hear me, O Lord (alleluia, alleluia).

℣. I will keep thy statutes.

℟. Hear me, O Lord (alleluia, alleluia).

℣. Glory be to the Father, and to the Son, and to the Holy Ghost.

℟. I will call with my whole heart; hear me, O Lord (alleluia, alleluia).

℣. Cleanse me, O Lord, from my secret faults (alleluia).

℟. Keep thy servant also from presumptuous sins (alleluia).

\subsection{Suffrages}
\begin{rubric}
	The Suffrages are said, kneeling, on all Ferias outside of Eastertide.
\end{rubric}

℣. Lord, have mercy upon us.

℟. Christ, have mercy upon us.

℣. Lord, have mercy upon us.

\lett{O}{ur} Father, who art in heaven, Hallowed be thy Name. Thy kingdom come. Thy will be done on earth, As it is in heaven. Give us this day our daily bread. And forgive us our trespasses, As we forgive those who trespass against us. And lead us not into temptation; But deliver us from evil. Amen.

℣. Cast me not away in the time of age.

℟. Forsake me not when my strength faileth me.

℣. Hide not thy face from me.

℟. Lest I be like unto them that go down into the pit.

℣. Quicken me, O Lord, for thy Name's sake.

℟. And for thy righteousness' sake bring my soul out of trouble.

\newcolumn

\subsection{Collects}
℣. The Lord be with you.\par
\textit{or,} O Lord, hear our prayer.

℟. And with thy spirit.\par
\textit{or,} And let our cry come unto thee.

℣. Let us pray.

\begin{rubric}
	The Collect(s) of the Day is here said, followed by this memorial Collect.
\end{rubric}

\lett{L}{ord} Jesus Christ, who for our sakes didst tread the paths of death; make known to us the way of life; that as thou wast reckoned with the transgressors in thy death, and with the rich in thy burial, so we, who are dead in trespasses and sins, may be raised up by thee to the land of true riches. Who livest and reignest with the Father and the Holy Ghost, ever one God, world without end. \textit{Amen.}\\

℣. The Lord be with you.\par
\textit{or,} O Lord, hear our prayer.

℟. And with thy spirit.\par
\textit{or,} And let our cry come unto thee.

℣. Let us bless the Lord.

℟. Thanks be to God.

℣. May the souls {\ding{64}} of the faithful departed, through the mercy of God, rest in peace.

\lett{N}{ow} unto him that is able to keep us from falling and to present us faultless before the presence of his glory with exceeding joy; to the only wise God our Saviour be glory and majesty, dominion, and power, both now and ever. \textit{Amen.}

\end{multicols}

%\begin{center}
%	{\textsc{Here endeth the Order of None.}}
%\end{center}
\include{pages/Office/Evensong}
\include{pages/Office/Compline}
\include{pages/Office/Mary}
%Prayers \& Thanksgivings
\cleardoublepage
\fancyhead[RE,LO]{}\fancyhead[RO,LE]{}
\fancyhead[C]{}\thispagestyle{empty}
\phantomsection
\addcontentsline{toc}{chapter}{Prayers \& Thanksgivings}

\begin{figure}[H]
    \centering
    \includegraphics[width=98mm, height=140mm]{Prayers/Prayer.eps}
    \caption{\textsc{\Huge{Prayers \& Thanksgivings}}}
\end{figure}
\include{pages/Prayers/PrayersThanksgivings}
%Litanies \& Offices
\cleardoublepage
\fancyhead[RE,LO]{}\fancyhead[RO,LE]{}
\fancyhead[C]{}
\thispagestyle{empty}
\phantomsection
\addcontentsline{toc}{chapter}{Litanies \& Offices}

\begin{figure}[H]
    \centering
    \includegraphics[width=98mm, height=140mm]{Litanies/DavidRepentance.eps}
    \caption{\textsc{\Huge{Litanies \& Offices of the Church}}}
\end{figure}
\include{pages/Litanies/GreatLitany}
\include{pages/Litanies/Penitential}
\include{pages/Litanies/BaptismalBlessing}
\section{Litany of Saints}\label{LitanySaints}
\fancyhead[RO,LE]{\textit{Litany of Saints}}
\fancyhead[RE,LO]{}

%MANUAL ADJUSTMENT:
\vspace{-1\baselineskip}

\begin{multicols}{2}
\par\noindent
Lord, have mercy upon us.\par
℟. Lord, have mercy upon us.\par\noindent
Christ, have mercy upon us.\par
℟. Christ, have mercy upon us.\par\noindent
Lord, have mercy upon us.\par
℟. Lord, have mercy upon us.\par\noindent
O Christ hear us.\par
℟. O Christ graciously hear us.\par\noindent
O God the Father of heaven,\par
℟. Have mercy upon us.\par\noindent
O God the Son, Redeemer of the world,\par
℟. Have mercy upon us.\par\noindent
O God the Holy Ghost,\par
℟. Have mercy upon us.\par\noindent
Holy Trinity one God,\par
℟. Have mercy upon us.\par\noindent
Holy Mary,\par
℟. Pray for us.\par\noindent
Holy Mother of God,\par
℟. Pray for us.\par\noindent
Holy Virgin of virgins,\par
℟. Pray for us.\par\noindent
Holy Michael,\par
℟. Pray for us.\par\noindent
Holy Gabriel,\par
℟. Pray for us.\par\noindent
Holy Raphael,\par
℟. Pray for us.\par\noindent
All ye holy Angels and Archangels,\par
℟. Pray for us.\par\noindent
All ye holy orders of blessed Spirits,\par
℟. Pray for us.\par\noindent
St. John Baptist,\par
℟. Pray for us.\par\noindent
St. Joseph,\par
℟. Pray for us.\par\noindent
All ye holy Patriarchs and Prophets,\par
℟. Pray for us.\par\noindent
St. Peter,\par
℟. Pray for us.\par\noindent
St. Paul,\par
℟. Pray for us.\par\noindent
St. Andrew,\par
℟. Pray for us.\par\noindent
St. James,\par
℟. Pray for us.\par\noindent
St. John,\par
℟. Pray for us.\par\noindent
St. Thomas,\par
℟. Pray for us.\par\noindent
St. James,\par
℟. Pray for us.\par\noindent
St. Philip,\par
℟. Pray for us.\par\noindent
St. Bartholomew,\par
℟. Pray for us.\par\noindent
St. Matthew,\par
℟. Pray for us.\par\noindent
St. Simon,\par
℟. Pray for us.\par\noindent
St. Jude,\par
℟. Pray for us.\par\noindent
St. Matthias,\par
℟. Pray for us.\par\noindent
St. Barnabas,\par
℟. Pray for us.\par\noindent
St. Luke,\par
℟. Pray for us.\par\noindent
St. Mark,\par
℟. Pray for us.\par\noindent
All ye holy Apostles and Evangelists,\par
℟. Pray for us.\par\noindent
All ye holy Disciples of the Lord,\par
℟. Pray for us.\par\noindent
All ye holy Innocents,\par
℟. Pray for us.\par\noindent
St. Stephen,\par
℟. Pray for us.\par\noindent
St. Lawrence,\par
℟. Pray for us.\par\noindent
St. Vincent,\par
℟. Pray for us.\par\noindent
St. Fabian and St. Sebastian,\par
℟. Pray for us.\par\noindent
St. John and St. Paul,\par
℟. Pray for us.\par\noindent
St. Cosmas and Damien,\par
℟. Pray for us.\par\noindent
St. Gervase and St. Protase,\par
℟. Pray for us.\par\noindent
All ye holy Martyrs,\par
℟. Pray for us.\par\noindent
St. Sylvester,\par
℟. Pray for us.\par\noindent
St. Gregory,\par
℟. Pray for us.\par\noindent
St. Ambrose,\par
℟. Pray for us.\par\noindent
St. Augustine,\par
℟. Pray for us.\par\noindent
St. Jerome,\par
℟. Pray for us.\par\noindent
St. Martin,\par
℟. Pray for us.\par\noindent
St. Nicholas,\par
℟. Pray for us.\par\noindent
All ye holy Bishops and Confessors,\par
℟. Pray for us.\par\noindent
All ye holy Doctors,\par
℟. Pray for us.\par\noindent
Holy St. Anthony,\par
℟. Pray for us.\par\noindent
Holy St. Benedict,\par
℟. Pray for us.\par\noindent
All ye holy Priests and Levites,\par
℟. Pray for us.\par\noindent
All ye holy Monks and Hermits,\par
℟. Pray for us.\par\noindent
St. Mary Magdalene,\par
℟. Pray for us.\par\noindent
St. Agatha,\par
℟. Pray for us.\par\noindent
St. Lucy,\par
℟. Pray for us.\par\noindent
St. Agnes,\par
℟. Pray for us.\par\noindent
St. Cecilia,\par
℟. Pray for us.\par\noindent
St. Catherine,\par
℟. Pray for us.\par\noindent
St. Anastasia,\par
℟. Pray for us.\par\noindent
All ye holy Virgins and Widows,\par
℟. Pray for us.\par\noindent
All ye Saints of God,\par
℟. Intercede for us.\par\noindent
Be thou merciful,\par
℟. Spare us, O Lord.\par\noindent
Be thou merciful,\par
℟. Hear us, O Lord.\par\noindent
From all evil,\par
℟. Deliver us, O Lord.\par\noindent
From all sin,\par
℟. Deliver us, O Lord.\par\noindent
From thy wrath,\par
℟. Deliver us, O Lord.\par\noindent
From sudden and unprepared death,\par
℟. Deliver us, O Lord.\par\noindent
From the crafts of the devil,\par
℟. Deliver us, O Lord.\par\noindent
From anger, hatred, and all uncharitableness,\par
℟. Deliver us, O Lord.\par\noindent
From the spirit of fornication,\par
℟. Deliver us, O Lord.\par\noindent
From lightning and tempest,\par
℟. Deliver us, O Lord.\par\noindent
From the scourge of earthquake,\par
℟. Deliver us, O Lord.\par\noindent
From pestilence, famine, and war,\par
℟. Deliver us, O Lord.\par\noindent
From everlasting death,\par
℟. Deliver us, O Lord.\par\noindent
By the mystery of thy holy Incarnation,\par
℟. Deliver us, O Lord.\par\noindent
By thine advent,\par
℟. Deliver us, O Lord.\par\noindent
By thy nativity,\par
℟. Deliver us, O Lord.\par\noindent
By thy baptism and holy fasting,\par
℟. Deliver us, O Lord.\par\noindent
By thy cross and passion,\par
℟. Deliver us, O Lord.\par\noindent
By thy death and burial,\par
℟. Deliver us, O Lord.\par\noindent
By thy holy resurrection,\par
℟. Deliver us, O Lord.\par\noindent
By thy wonderful ascension,\par
℟. Deliver us, O Lord.\par\noindent
By the coming of the Holy Ghost the Paraclete,\par
℟. Deliver us, O Lord.\par\noindent
In the day of judgment,\par
℟. Deliver us, O Lord.\par\noindent
We sinners,\par
℟. Beseech thee, hear us.\par\noindent
That thou wouldest spare us,\par
℟. We beseech thee, hear us.\par\noindent
That thou wouldest vouchsafe to bring us to true repentance,\par
℟. We beseech thee, hear us.\par\noindent
That thou wouldest vouchsafe to govern and preserve thy holy Church,\par
℟. We beseech thee, hear us.\par\noindent
That thou wouldest vouchsafe to preserve our Apostolic lord and all orders of the Church in holy religion,\par
℟. We beseech thee, hear us.\par\noindent
That thou wouldest vouchsafe to humble the enemies of holy Church,\par
℟. We beseech thee, hear us.\par\noindent
That thou wouldest vouchsafe to give unto Christian kings and rulers peace and true concord,\par
℟. We beseech thee, hear us.\par\noindent
That thou wouldest vouchsafe to grant to all the Christian people peace and unity,\par
℟. We beseech thee, hear us.\par\noindent
That thou wouldest vouchsafe to recall such as do err into the unity of the Church, and to bring all heathen unto the light of the gospel,\par
℟. We beseech thee, hear us.\par\noindent
That thou wouldest vouchsafe to strengthen, and preserve us in thy holy service,\par
℟. We beseech thee, hear us.\par\noindent
That thou wouldest vouchsafe to lift up our minds unto heavenly desires,\par
℟. We beseech thee, hear us.\par\noindent
That thou wouldest vouchsafe to reward all our benefactors with everlasting blessings,\par
℟. We beseech thee, hear us.\par\noindent
That thou wouldest vouchsafe to deliver our souls, and the souls of our brethren, kinsfolk, and benefactors, from eternal damnation,\par
℟. We beseech thee, hear us.\par\noindent
That thou wouldest vouchsafe to give and preserve the fruits of the earth,\par
℟. We beseech thee, hear us.\par\noindent
That thou wouldest vouchsafe to grant unto all the faithful departed {\ding{64}} rest eternal,\par
℟. We beseech thee, hear us.\par\noindent
That thou wouldest vouchsafe graciously to hear us,\par
℟. We beseech thee, hear us.\par\noindent
Son of God,\par
℟. We beseech thee, hear us.\par\noindent
O Lamb of God, who takest away the sins of the world,\par
℟. Spare us, O Lord, hear us.\par\noindent
O Lamb of God, who takest away the sins of the world,\par
℟. Graciously hear us, O Lord.\par\noindent
O Lamb of God, who takest away the sins of the world,\par
℟. Have mercy upon us.\par\noindent
O Christ, hear us.\par
℟. O Christ, graciously hear us.\par\noindent
Lord, have mercy upon us.\par
℟. Christ, have mercy upon us.\par\noindent
Lord, have mercy upon us.
\end{multicols}

\begin{rubric}
    {The Lord's Prayer is here said until,}
\end{rubric}

℣. And lead us not into temptation,

℟. But deliver us from evil.

\subsection{Psalm 70. \textit{Deus, in adjutorium}}
\lett{H}{aste} thee, O God, to deliver me : make haste to help me, O Lord.\par
\secondline{\psanum{2}Let them be ashamed and confounded that seek after my soul : let them be turned backward and put to confusion that wish me evil.}
\thirdline{\psanum{3}Let them for their reward be soon brought to shame : that cry over me, There, there.}
\psanum{4}But let all those that seek thee be joyful and glad in thee : and let all such as delight in thy salvation say alway, The Lord be praised.\par
\psanum{5}As for me, I am poor and in misery : haste thee unto me, O God.\par
\psanum{6}Thou art my helper and my redeemer : O Lord, make no long tarrying.\par
℣. Glory be to the Father, and to the Son, and to the Holy Ghost.\par
℟. As it was in the beginning, is now, and ever shall be, world without end. Amen.

\begin{multicols}{2}
\sloppy

℣. Save thy servants.\par
℟. My God, who put their trust in thee.\par
℣. Be unto us, O Lord, a strong tower.\par
℟. From the face of the enemy.\par
℣. Let the enemy have no advantage over us.\par
℟. Nor the son of wickedness approach to hurt us.\par
℣. O Lord, deal not with us after our sins.\par
℟. Neither reward us after our iniquities.\par
℣. Let us pray for our \textit{Bishop}, \textit{N.}\par
℟. The Lord preserve him and keep him alive, and make him blessed upon earth, and deliver him not into the hands of his enemies.\par
℣. Let us pray for our benefactors.\par
℟. Vouchsafe, O Lord, to reward with eternal life all them that do us good for thy name's sake. Amen.\par
℣. Let us pray for the faithful departed.\par
℟. Rest eternal grant unto them, O Lord, and let light perpetual shine upon them.\par
℣. May they rest in peace.\par
℟. Amen.\par
℣. For our absent brethren.\par
℟. My God, save thy servants who put their trust in thee.\par
℣. O Lord, send them help from thy holy place.\par
℟. And strengthen them out of Sion.\par
℣. O Lord, hear my prayer.\par
℟. And let my prayer come unto thee.\par
℣. The Lord be with you.\par
℟. And with thy spirit.\par
℣. Let us pray.

\fussy
\end{multicols}

\subsection{Collects}
\lett{O}{God,} whose nature and property is ever to have mercy and to forgive: receive our humble petitions; and though we and all thy servants be bound with the chain of our sins, yet let the pitifulness of thy great mercy loose us.

\lett{W}{e} beseech thee, O Lord, mercifully to hear the prayers of thy humble servants, and to forgive the sins of them that confess the same unto thee: that they may obtain of thy loving-kindness pardon and peace.

\lett{G}{raciously} shew forth upon us, O Lord, thy unspeakable mercy: that thou wouldest both loose us from all our sins, and likewise deliver us from the punishment which for the same we deserve.

\lett{O}{God,} who art wroth with them that sin against thee, and sparest them that are penitent: mercifully look upon the prayers of thy people who call upon thee; and turn away the scourges of thy wrath, which for our sins we justly deserve.

\par\noindent
	\centerline{\rule{0.5\textwidth}{0.4pt}}
\par\noindent
\begin{rubric}
	If the Bishoprick be vacant, the following prayer is omitted.
\end{rubric}
\lett{A}{lmighty} and everlasting God, have mercy upon thy servant \textit{N.}, our Bishop, and, according to thy great goodness, direct him into the way of everlasting salvation: that, by thy grace, he may desire that which is well pleasing unto thee, and with all his strength perform the same.
\par\noindent
	\centerline{\rule{0.5\textwidth}{0.4pt}}
\par\noindent

%MANUAL ADJUSTMENT:
%\vspace{0.5ex}
\clearpage
\lett{O}{God,} from whom all holy desires, all good counsels, and all just works do proceed: give unto thy servants that peace which the world cannot give; that both our hearts may be set to obey thy commandments, and also, that by thee we being defended from the fear of our enemies may pass our time in rest and quietness.

\lett{K}{indle,} O Lord, with the fire of the Holy Spirit our reins and our hearts: that we may serve thee with a chaste body, and please thee with a clean heart.

\lett{O}{God,} the Creator and Redeemer of all the faithful, grant unto the souls of thy servants and handmaids the remission of all their sins: that through devout supplications they may obtain the pardon which they have alway desired.

\lett{P}{revent} us, O Lord, we beseech thee, in all our doings with thy most gracious favour, and further us with thy continual help: that all our prayer and work may be begun, continued, and ended in thee.

\lett{A}{lmighty} and everlasting God, who hast dominion both of the living and of the dead, and hast mercy upon all whom thou foreknowest will be thine in faith and works: we humbly beseech thee; that all those for whom we are minded to pour forth our prayers, whether in this present world they still be held in the flesh, or being delivered from the body have passed into that which is to come, may at the intercession of all thy Saints obtain of thy bountiful goodness the remission of all their sins. Through Jesus Christ thy Son, our Lord, who liveth and reigneth with thee, in the unity of the Holy Ghost, God, throughout all ages, world without end.

℟. Amen.

\subsection{Conclusion}
℣. The Lord be with you.

℟. And with thy spirit.

℣. May the almighty and merciful Lord graciously hear us.

℟. Amen.

%English Missal corrected here for conformity with the rest of the BCP.
℣. And may the souls {\ding{64}} of the faithful departed, through the mercy of God, rest in peace.\par
℟. Amen.
%\begin{center}
%    \textsc{Here endeth the Litany of the Saints.}
%\end{center}
%Office of the Dead
\cleardoublepage
\fancyhead[RE,LO]{}
\fancyhead[RO,LE]{}
\fancyhead[C]{}
\thispagestyle{empty}
\phantomsection
\addcontentsline{toc}{chapter}{Office of the Dead}

\begin{figure}[H]
    \centering
    \includegraphics[width=98mm, height=140mm]{OfficeDead/Harrowing.eps}
    \caption{\textsc{\Huge{Office of the Dead}}}
\end{figure}
\include{pages/OfficeDead/VesperDead}
\include{pages/OfficeDead/MatinsDead}
\section{Lauds of the Dead}\label{laudsdead}
\fancyhead[RO,LE]{\textit{Lauds}}
\begin{secrubric}
    %The Monastic Diurnal is ambiguous on this rubric, but the Diurnale is clear.
    Before Lauds---if it be said separately from Matins---the Lord's Prayer and Angelic Salutation are said secretly. Otherwise, the Office begins at once with the Antiphon.
\end{secrubric}
\subsection{Psalm 51. \textit{Miserere mei, Deus}}\par\noindent
\textit{Ant.} The bones which thou hast broken {\dag} shall rejoice in the Lord.\par
\lett{H}{ave} mercy upon me, O God, after thy great goodness : according to the multitude of thy mercies do away mine offences.\par
\secondline{\psanum{2}Wash me throughly from my wickedness : and cleanse me from my sin.}
\thirdline{\psanum{3}For I acknowledge my faults : and my sin is ever before me.}
\psanum{4}Against thee only have I sinned, and done this evil in thy sight : that thou mightest be justified in thy saying, and clear when thou art judged.\par
\psanum{5}Behold, I was shapen in wickedness : and in sin hath my mother conceived me.\par
\psanum{6}But lo, thou requirest truth in the inward parts: and shalt make me to understand wisdom secretly.\par
\psanum{7}Thou shalt purge me with hyssop, and I shall be clean : thou shalt wash me, and I shall be whiter than snow.\par
\psanum{8}Thou shalt make me hear of joy and gladness : that the bones which thou hast broken may rejoice.\par
\psanum{9}Turn thy face from my sins : and put out all my misdeeds.\par
\psanum{10}Make me a clean heart, O God : and renew a right spirit within me.\par
\psanum{11}Cast me not away from thy presence : and take not thy holy Spirit from me.\par
\psanum{12}O give me the comfort of thy help again : and stablish me with thy free Spirit.\par
\psanum{13}Then shall I teach thy ways unto the wicked : and sinners shall be converted unto thee.\par
\psanum{14}Deliver me from blood-guiltiness, O God, thou that art the God of my health : and my tongue shall sing of thy righteousness.\par
\psanum{15}Thou shalt open my lips, O Lord : and my mouth shall shew thy praise.\par
\psanum{16}For thou desirest no sacrifice, else would I give it thee : but thou delightest not in burnt-offerings.\par
\psanum{17}The sacrifice of God is a troubled spirit : a broken and contrite heart, O God, shalt thou not despise.\par
\psanum{18}O be favourable and gracious unto Sion : build thou the walls of Jerusalem.\par
\psanum{19}Then shalt thou be pleased with the sacrifice of righteousness, with the burnt-offerings and oblations : then shall they offer young bullocks upon thine altar.\par
℣. Rest eternal * grant unto them, O Lord.\par
℟. And let light perpetual * shine upon them.\par\noindent
\textit{Ant.} The bones which thou hast broken shall rejoice in the Lord.

\subsection{Psalm 65. \textit{Te decet hymnus}}\par\noindent
\textit{Ant.} Thou, O Lord, {\dag} that hearest the prayer, unto thee shall all flesh come.\par
\lett{T}{hou}, O God, art praised in Sion : and unto thee shall the vow be performed in Jerusalem.\par
\secondline{\psanum{2}Thou that hearest the prayer : unto thee shall all flesh come.}
\thirdline{\psanum{3}My misdeeds prevail against me : O be thou merciful unto our sins.}
\psanum{4}Blessed is the man whom thou choosest, and receivest unto thee : he shall dwell in thy court, and shall be satisfied with the pleasures of thy house, even of thy holy temple.\par
\psanum{5}Thou shalt shew us wonderful things in thy righteousness, O God of our salvation : thou that art the hope of all the ends of the earth, and of them that remain in the broad sea.\par
\psanum{6}Who in his strength setteth fast the mountains : and is girded about with power.\par
\psanum{7}Who stilleth the raging of the sea : and the noise of his waves, and the madness of the people.\par
\psanum{8}They also that dwell in the uttermost parts of the earth shall be afraid at thy tokens : thou that makest the outgoings of the morning and evening to praise thee.\par
\psanum{9}Thou visitest the earth, and blessest it : thou makest it very plenteous.\par
\psanum{10}The river of God is full of water : thou preparest their corn, for so thou providest for the earth.\par
\psanum{11}Thou waterest her furrows, thou sendest rain into the little valleys thereof : thou makest it soft with the drops of rain, and blessest the increase of it.\par
\psanum{12}Thou crownest the year with thy goodness : and thy clouds drop fatness.\par
\psanum{13}They shall drop upon the dwellings of the wilderness : and the little hills shall rejoice on every side.\par
\psanum{14}The folds shall be full of sheep : the valleys also shall stand so thick with corn, that they shall laugh and sing.\par
℣. Rest eternal * grant unto them, O Lord.\par
℟. And let light perpetual * shine upon them.\par\noindent
\textit{Ant.} Thou, O Lord, that hearest the prayer, unto thee shall all flesh come.

\subsection{Psalm 63. \textit{Deus, Deus meus}}\par\noindent
\textit{Ant.} Thy right hand {\dag} hath upholden me, O Lord.\par
\lett{O}{God,} thou art my God : early will I seek thee.\par
\secondline{\psanum{2}My soul thirsteth for thee, my flesh also longeth after thee : in a barren and dry land where no water is.}
\thirdline{\psanum{3}Thus have I looked for thee in holiness : that I might behold thy power and glory.}
\psanum{4}For thy loving-kindness is better than the life itself : my lips shall praise thee.\par
\psanum{5}As long as I live will I magnify thee on this manner : and lift up my hands in thy Name.\par
\psanum{6}My soul shall be satisfied, even as it were with marrow and fatness : when my mouth praiseth thee with joyful lips.\par
\psanum{7}Have I not remembered thee in my bed : and thought upon thee when I was waking?\par
\psanum{8}Because thou hast been my helper : therefore under the shadow of thy wings will I rejoice.\par
\psanum{9}My soul hangeth upon thee : thy right hand hath upholden me.\par
\psanum{10}These also that seek the hurt of my soul : they shall go under the earth.\par
\psanum{11}Let them fall upon the edge of the sword : that they may be a portion for foxes.\par
\psanum{12}But the King shall rejoice in God; all they also that swear by him shall be commended : for the mouth of them that speak lies shall be stopped.\par
℣. Rest eternal * grant unto them, O Lord.\par
℟. And let light perpetual * shine upon them.\par\noindent
\textit{Ant.} Thy right hand hath upholden me, O Lord.

\subsection{Ego dixi}\par\noindent
\textit{Ant.} From the gate of hell, {\dag} deliver my soul, O Lord.\par
\lett{I}{said} in the cutting off of my days, {\dag} I shall go to the gates of the grave: * I am deprived of the residue of my years.\par
\secondline{I said, I shall not see the Lord, {\dag} even the Lord, in the land of the living: * I shall behold man no more with the inhabitants of the world.}
\thirdline{Mine age is departed, * and is removed from me as a shepherd's tent;}
I have cut off like a weaver my life: {\dag} he will cut me off with pining sickness: * from day even to night wilt thou make an end of me.\par
I reckoned till morning, {\dag} that, as a lion, so will he break all my bones: * from day even to night wilt thou make an end of me.\par
Like a crane or a swallow, so did I chatter: * I did mourn as a dove.\par
Mine eyes fail with looking upward: * O Lord, I am oppressed; undertake for me.\par
What shall I say? {\dag} he hath both spoken unto me, and himself hath done it: * I shall go softly all my years in the bitterness of my soul.\par
O Lord, by these things men live, {\dag} and in all these things is the life of my spirit: * so wilt thou recover me, and make me to live.\par
Behold, for peace I had great bitterness: {\dag} but thou hast in love to my soul delivered it from the pit of corruption: * for thou hast cast all my sins behind thy back.\par
For the grave cannot praise thee, {\dag} death can not celebrate thee: * they that go down into the pit cannot hope for thy truth.\par
The living, the living, he shall praise thee, as I do this day: * the father to the children shall make known thy truth.\par
The Lord was ready to save me: * therefore we will sing my songs to the stringed instruments all the days of our life in the house of the Lord.\par
℣. Rest eternal * grant unto them, O Lord.\par
℟. And let light perpetual * shine upon them.\par\noindent
\textit{Ant.} From the gate of hell, deliver my soul, O Lord.\par

\subsection{Psalm 150. \textit{Laudate Dominum}}\par\noindent
\textit{Ant.} Let everything that hath breath {\dag} praise the Lord.\par
\lett{O}{praise} God in his holiness : praise him in the firmament of his power.\par
\secondline{\psanum{2}Praise him in his noble acts : praise him according to his excellent greatness.}
\thirdline{\psanum{3}Praise him in the sound of the trumpet : praise him upon the lute and harp.}
\psanum{4}Praise him in the cymbals and dances : praise him upon the strings and pipe.\par
\psanum{5}Praise him upon the well-tuned cymbals : praise him upon the loud cymbals.\par
\psanum{6}Let every thing that hath breath : praise the Lord.\par
℣. Rest eternal * grant unto them, O Lord.\par
℟. And let light perpetual * shine upon them.\par\noindent
\textit{Ant.} Let everything that hath breath praise the Lord.\\

℣. I heard a voice from heaven, saying unto me.\par
℟. Blessed are the dead which die in the Lord.

%MANUAL ADJUSTMENT:
\clearpage
\subsection{Benedictus}\par\noindent
\textit{Ant.} I am the resurrection {\dag} and the life: he that believeth in me, though he were dead, yet shall he live: and whosoever liveth and believeth in me, shall never die.
\lett{B}{lessed} {\ding{64}} be the Lord God of Israel; * for he hath visited and redeemed his people;\par
    \secondline{And hath raised up a mighty salvation for us, * in the house of his servant David;}
    \thirdline{As he spake by the mouth of his holy Prophets, * which have been since the world began;}
    That we should be saved from our enemies, * and from the hand of all that hate us.\par
    To perform the mercy promised to our forefathers, * and to remember his holy covenant;\par
    To perform the oath which he sware to our forefather Abraham, * that he would give us;\par
    That we being delivered out of the hand of our enemies * might serve him without fear;\par
    In holiness and righteousness before him, * all the days of our life.\par
    And thou, child, shalt be called the prophet of the Highest: * for thou shalt go before the face of the Lord to prepare his ways;\par
    To give knowledge of salvation unto his people * for the remission of their sins,\par
    Through the tender mercy of our God; * whereby the day-spring from on high hath visited us;\par
    To give light to them that sit in darkness, and in the shadow of death, * and to guide our feet into the way of peace.\par
    ℣. Rest eternal * grant unto them, O Lord.\par
    ℟. And let light perpetual * shine upon them.\par\noindent
\textit{Ant.} I am the resurrection and the life: he that believeth in me, though he were dead, yet shall he live; and whosoever liveth and believeth in me, shall never die.

\subsection{Lord's Prayer}
\begin{rubric}
The Lord's Prayer is said, in secret, kneeling, ending with,
\end{rubric}
℣. And lead us not into temptation.

℟. But deliver us from evil.

\subsection{Psalm 130. \textit{De profundis}}
\begin{rubric}
{Psalm 130 is not said on All Souls Day, on the day of death or burial, nor at any time when the Office is recited with Double rite.}
\end{rubric}
\lett{O}{ut} of the deep have I called unto thee, O Lord : Lord, hear my voice.\par
\secondline{\psanum{2}O let thine ears consider well : the voice of my complaint.}
\thirdline{\psanum{3}If thou, Lord, wilt be extreme to mark what is done amiss : O Lord, who may abide it?}
\psanum{4}For there is mercy with thee : therefore shalt thou be feared.\par
\psanum{5}I look for the Lord; my soul doth wait for him : in his word is my trust.\par
\psanum{6}My soul fleeth unto the Lord : before the morning watch, I say, before the morning watch.\par
\psanum{7}O Israel, trust in the Lord, for with the Lord there is mercy : and with him is plenteous redemption.\par
\psanum{8}And he shall redeem Israel : from all his sins.\par
℣. Rest eternal * grant unto them, O Lord.\par
℟. And let light perpetual * shine upon them.

\subsection{Responsory}
℣. From the gate of hell.

℟. Deliver \textit{his soul}, O Lord.

℣. May \textit{he} rest in peace.

℟. Amen.

\subsection{Conclusion}
℣. O Lord, hear my prayer.

℟. And let my cry come unto thee.

℣. The Lord be with you.

℟. And with thy spirit.

\begin{rubric}
{The appropriate collect is here said (p. \pageref{dead}).}
\end{rubric}

℣. Rest eternal * grant unto them, O Lord.

℟. And let light perpetual * shine upon them.

℣. May they rest in peace.

℟. Amen.

\begin{center}
    \textsc{Here endeth the Order of Lauds of the Dead.}
\end{center}
%Eucharistic Devotions
\cleardoublepage
\fancyhead[RE,LO]{}\fancyhead[RO,LE]{}
\fancyhead[C]{}
\thispagestyle{empty}
\phantomsection
\addcontentsline{toc}{chapter}{Eucharistic Devotions}

\begin{figure}[H]
    \centering
    \includegraphics[width=98mm, height=140mm]{Eucharistic/Supper.eps}
    \caption{\textsc{\Huge{Eucharistic Devotions}}}
\end{figure}
\include{pages/Eucharistic/Benediction}
\include{pages/Eucharistic/Preparation}
\include{pages/Eucharistic/Thanksgiving}
\include{pages/Eucharistic/Asperges}
\section{Prayers at the Foot of the Altar}
\fancyhead[RO,LE]{\textit{Foot of the Altar}}
\begin{secrubric}
    A Hymn may be sung during the \blackrubric{Prayers at the Foot of the Altar,} except in a Low Mass.
\end{secrubric}
\begin{secrubric}
    After the Rite of Aspersion (or immediately upon processing to the Foot of the Altar if the Rite not occur), the Priest---having vested---approaches the Altar, makes the due reverence, signs himself with the sign of the Cross from forehead to breast, and says in the clear voice:
\end{secrubric}
\elcol{℣. In the Name of the Father, {\ding{64}} and of the Son, and of the Holy Ghost.}{℣. In nómine Patris, {\ding{64}} et Fílii, et Spíritus Sancti.}
\elcol{℟. Amen.}{℟. Amen.}
\begin{rubric}
    Then, with hands joined before his breast, the Priest begins the Antiphon:
\end{rubric}
\elcol{℣. I will go unto the altar of God.}{℣. Introíbo ad altáre Dei.}
\elcol{℟. Even unto the God of my joy and gladness.}{℟. Ad Deum, qui l{\ae}tíficat juventútem meam.}
\begin{rubric}
    In Masses for the Dead, and in Masses of the Season from Passion Sunday to Holy Saturday exclusive, the following Psalm \blackrubric{Give sentence} is omitted together with \blackrubric{Glory be} and the repetition of the Antiphon. And the Priest then moves immediately to \blackrubric{Our help.}
\end{rubric}
    \par\noindent
	\centerline{\rule{0.5\textwidth}{0.4pt}}\\
	\par\noindent
\elcol{℣. Give sentence with me, O God, and defend my cause against the ungodly people: O deliver me from the deceitful and wicked man.

℟. For thou art the God of my strength; why hast thou put me from thee: and why go I so heavily, while the enemy oppresseth me?

℣. O send out thy light and thy truth, that they may lead me: and bring me unto thy holy hill, and to thy dwelling.

℟. And that I may go unto the altar of God, even the God of my joy and gladness: and upon the harp will I give thanks unto thee, O God, my God.

℣. Why art thou so heavy, O my soul: and why art thou so disquieted within me?

℟. O put thy trust in God: for I will yet give him thanks, which is the help of my countenance, and my God.

℣. Glory be to the Father, and to the Son: and to the Holy Ghost.

℟. As it was in the beginning, is now, and ever shall be: world without end. Amen.

℣. I will go unto the altar of God:

℟. Even unto the God of my joy and gladness.}
{℣. Júdica me, Deus, et discérne causam meam de gente non sancta: ab hómine iníquo et dolóso érue me.

℟. Quia tu es, Deus, fortitúdo mea: quare me repulísti, et quare tristis incédo, dum afflígit me inimícus?

℣. Emítte lucem tuam et veritátem tuam: ipsa me deduxérunt, et adduxérunt in montem sanctum tuum et in tabernácula tua.

℟. Et introíbo ad altáre Dei: ad Deum, qui l{\ae}tíficat juventútem meam.

℣. Confitébor tibi in cíthara, Deus, Deus meus: quare tristis es, ánima mea, et quare contúrbas me?

℟. Spera in Deo, quóniam adhuc confitébor illi: salutáre vultus mei, et Deus meus.

℣. Glória Patri, et Fílio, et Spirítui Sancto.

℟. Sicut erat in princípio, et nunc, et semper: et in sǽcula s{\ae}culórum. Amen.

℣. Introíbo ad altáre Dei.

℟. Ad Deum, qui l{\ae}tíficat juventútem meam.}
    \par\noindent
	\centerline{\rule{0.5\textwidth}{0.4pt}}\\
	\par\noindent
\elcol{℣. Our help is in the Name {\ding{64}} of the Lord:

℟. Who hath made heaven and earth.}{℣. Adjutórium nostrum {\ding{64}} in nómine Dómini.

℟. Qui fecit c{\ae}lum et terram.}
\begin{rubric}
    Then, with hands joined, bowing profoundly, the Priest says,
\end{rubric}
\elcol{\lett{I}{confess} to Almighty God, to Blessed Mary ever Virgin, to Blessed Michael the Archangel, to Blessed John Baptist, to the Holy Apostles Peter and Paul, to all the Saints, and to you, brethren: that I have sinned exceedingly in thought, word, and deed: \inrub{Strike breast thrice.} by my fault, my own fault, my own most grievous fault. Wherefore I beg Blessed Mary ever Virgin, Blessed Michael the Archangel, Blessed John Baptist, the Holy Apostles Peter and Paul, all the Saints, and you, brethren, to pray for me to the Lord our God.}{\lett{C}{onf\smash{í}teor} Deo omnipoténti, beát{\ae} Marí{\ae} semper Vírgini, beáto Michaéli Archángelo, beáto Joánni Baptíst{\ae}, sanctis Apóstolis Petro et Paulo, ómnibus Sanctis, et vobis, fratres: quia peccávi nimis cogitatióne, verbo et ópere: \inrub{Strike breast thrice.} mea culpa, mea culpa, mea máxima culpa. Ideo precor beátam Maríam semper Vírginem, beátum Michaélem Archángelum, beátum Joánnem Baptístam, sanctos Apóstolos Petrum et Paulum, omnes Sanctos, et vos, fratres, oráre pro me ad Dóminum, Deum nostrum.}
\elcol{℟. Almighty God have mercy upon thee, forgive thee thy sins, and bring thee to everlasting life.}{℟. Misereátur tui omnípotens Deus, et, dimíssis peccátis tuis, perdúcat te ad vitam {\ae}térnam.}
\elcol{℟. Amen.\margcolrub{1}{The Priest rises. The Ministers bow.}}{℟. Amen.\marglatincolrub{1}{The Priest rises. The Ministers bow.}}
\clearpage
\begin{rubric}
	The Ministers, bowing, say,
\end{rubric}
\elcol{\lett{I}{confess} to Almighty God, to Blessed Mary ever Virgin, to Blessed Michael the Archangel, to Blessed John Baptist, to the Holy Apostles Peter and Paul, to all the Saints, and to thee, Father: that I have sinned exceedingly in thought, word, and deed: \inrub{Strike breast thrice.} by my fault, my own fault, my own most grievous fault. Wherefore I beg Blessed Mary ever Virgin, Blessed Michael the Archangel, Blessed John Baptist, the Holy Apostles Peter and Paul, all the Saints, and thee, Father, to pray for me to the Lord our God.}{\lett{C}{onf\smash{í}teor} Deo omnipoténti, beát{\ae} Marí{\ae} semper Vírgini, beáto Michaéli Archángelo, beáto Joánni Baptíst{\ae}, sanctis Apóstolis Petro et Paulo, ómnibus Sanctis, et tibi, pater: quia peccávi nimis cogitatióne, verbo et ópere: \inrub{Strike breast thrice.} mea culpa, mea culpa, mea máxima culpa. Ideo precor beátam Maríam semper Vírginem, beátum Michaélem Archángelum, beátum Joánnem Baptístam, sanctos Apóstolos Petrum et Paulum, omnes Sanctos, et te, pater, oráre pro me ad Dóminum, Deum nostrum.}
\elcol{℣. Almighty God have mercy upon you, forgive you your sins, and bring you to everlasting life.}{℣. Misereátur vestri omnípotens Deus, et, dimíssis peccátis vestris, perdúcat vos ad vitam {\ae}térnam.}
\elcol{℟. Amen.}{℟. Amen.}
\elcol{℣. The Almighty and Merciful Lord grant unto us pardon, {\ding{64}} absolution, and remission of our sins.

℟. Amen.\margcolrub{2}{Ministers rise.}

℣. Wilt thou not turn again and quicken us, O God?

℟. That thy people may rejoice in thee.

℣. O Lord, show thy mercy upon us.

℟. And grant us thy salvation.

℣. O Lord, hear my prayer.

℟. And let my cry come unto thee.

℣. The Lord be with you.

℟. And with thy spirit.}
{℣. Indulgéntiam, {\ding{64}} absolutiónem et remissiónem peccatórum nostrórum tríbuat nobis omnípotens et miséricors Dóminus.

℟. Amen.\marglatincolrub{2}{Ministers rise.}

℣. Deus, tu convérsus vivificábis nos?

℟. Et plebs tua l{\ae}tábitur in te.

℣. Osténde nobis, Dómine, misericórdiam tuam.

℟. Et salutáre tuum da nobis.

℣. Dómine, exáudi oratiónem meam.

℟. Et clamor meus ad te véniat.

℣. Dóminus vobíscum.

℟. Et cum spíritu tuo.}

\begin{rubric}
    The Priest, extending and joining his hands, says in a loud voice:
\end{rubric}

\elcol{℣. Let us pray.}{℣. Orémus.}
%Holy Mass
\cleardoublepage
\fancyhead[RE,LO]{}\fancyhead[RO,LE]{}
\fancyhead[C]{}\thispagestyle{empty}
\phantomsection
\addcontentsline{toc}{chapter}{Holy Mass}

\begin{figure}[H]
    \centering
    \includegraphics[width=98mm, height=140mm]{HolyMass/Crucifixion.eps}
    \caption{\textsc{\Huge{Holy Mass}}}
\end{figure}
\fancyhead[C]{{\LARGE Holy Mass}}
\fancyhead[RO,LE]{}
\fancyhead[RE,LO]{}
%\fancyhead[RO,LE]{\textit{Liturgy of St. Tikhon}}
%\section{Liturgy of St. Tikhon}\label{tikhon}
\phantomsection
\addcontentsline{toc}{section}{Liturgy of St. Tikhon}
%\begin{rubric}
%If any be an open and notorious evil liver, or have done any wrong to his neighbours by word or deed, so that the Congregation be thereby offended; the Curate, having knowledge thereof, shall call him and advertise him, that in any wise he presume not to come to the Lord’s Table, until he have openly declared himself to have truly repented and amended his former naughty life, that the Congregation may thereby be satisfied, which before were offended; and that he have recompensed the parties, to whom he hath done wrong; or at least declare himself to be in full purpose so to do, as soon as he conveniently may.
%\end{rubric}
%\begin{rubric}
%The same order shall the Curate use with those betwixt whom he perceiveth malice and hatred to reign; not suffering them to be partakers of the Lord's Table, until he know them to be reconciled. And if one of the parties so at variance be content to forgive from the bottom of his heart all that the other hath trespassed against him, and to make amends for that he himself hath offended; and the other party will not be persuaded to a godly unity, but remain still in his frowardness and malice the Minister in that case ought to admit the penitent person to the Holy Communion, and not him that is obstinate.
%\end{rubric}
%\begin{rubric}
%Provided that every Minister so advertising or repelling any, as is specified in the two precedent paragraphs, shall he obliged forthwith to give an account of the same to the Bishop, and therein to obey his order and direction.
%\end{rubric}
%
%
%There is a dispute between the Missals as to where to place the Collect for Purity and the Summary of the Law. The St. Tikhon's Liturgy (as it exists) and the English Missal place it at the Altar, after the Ascent. While this is a fitting place for the Collect, it creates a liturgical novelty for the Summary of the Law. At the Epistle side is offered the prayers to God: the various propers and the sacrifice of the Epistle. In contrast, the Summary of the Law is not directed towards God (who needs no preaching of His own Law) but rather towards the people. Therefore, it is absurd to have the place of sacrifice to God (the Epistle side) be the place of exhortation to the People of God, especially if does not even turn towards the Congregation. There is an example where the Priest does fully turn away from God towards the Congregation: the \emph{Orate Frates}, which occurs in the middle of the Altar. Therefore, to avoid novelties and liturgical inconsistencies, I have placed the Collect and the Summary/Decalogue at the Foot of the Altar, as other Missals have so done (it is appropriate for the Collect, for that is where other collects are read, and it is supremely expedient for an exhortation of the People).
\subsection{Collect for Purity}
\begin{secrubric}
    The Priest, at the Foot of the Altar, saith the Collect for Purity, as followeth, in an audible voice.
\end{secrubric}
\lett{A}{lmighty} God, unto whom all hearts are open, all desires known, and from whom no secrets are hid; Cleanse the thoughts of our hearts by the inspiration of thy Holy Spirit, that we may perfectly love thee, and worthily magnify thy holy Name; through Christ our Lord. \textit{Amen.}
\begin{rubric}
Then shall the Priest, turning to the People, rehearse distinctly---in an audible voice---The Ten Commandments; and the People, still kneeling, shall, after every Commandment, ask God mercy for their transgressions for the time past, and grace to keep the law for the time to come.
\end{rubric}
\begin{rubric}
And \textsc{Note,} that in rehearsing the Ten Commandments, the Priest may omit that part of the Commandment which is inset.
\end{rubric}
\begin{rubric}
	The Decalogue may be omitted. But \textsc{Note,} That when ever it is omitted, the Minister shall say the Summary of the Law, beginning, \blackrubric{Hear what our Lord Jesus Christ saith.}
\end{rubric}

%MANUAL ADJUSTMENT:
\vspace{-2ex}

\subsection{Ten Commandments}

%MANUAL ADJUSTMENT:
\vspace{-0.5ex}

\begin{center}
	\textsc{God spake these words, and said:}
\end{center}

%MANUAL ADJUSTMENT:
\vspace{-0.5ex}

%\begin{multicols}{2}
\par\noindent
    I am the \textsc{Lord} thy God; Thou shalt have none other gods but me.\par
    \textit{Lord, have mercy upon us, and incline our hearts to keep this law.}
    \par\noindent
    Thou shalt not make to thyself any graven image, nor the likeness of any thing that is in heaven above, or in the earth beneath, or in the water under the earth; thou shalt not bow down to them, nor worship them:
    \par\noindent
    \leftskip=2em
	{\small{for I the Lord thy God am a jealous God, and visit the sins of the fathers upon the children, unto the third and fourth generation of them that hate me; and show mercy unto thousands in them that love me and keep my commandments.}}
	\par
	\leftskip=0cm
    \textit{Lord, have mercy upon us, and incline our hearts to keep this law.}
    \par\noindent
    Thou shalt not take the Name of the Lord thy God in vain;
    \par\noindent
    \leftskip=2em
	{\small{for the Lord will not hold him guiltless, that taketh his Name in vain.}}
	\par
	\leftskip=0em
    \textit{Lord, have mercy upon us, and incline our hearts to keep this law.}
    \par\noindent
    Remember that thou keep holy the Sabbath-day.
    \par\noindent
    \leftskip=2em
	{\small{Six days shalt thou labour, and do all that thou hast to do; but the seventh day is the Sabbath of the Lord thy God. In it thou shalt do no manner of work; thou, and thy son, and thy daughter, thy man-servant, and thy maid-servant, thy cattle, and the stranger that is within thy gates. For in six days the Lord made heaven and earth, the sea, and all that in them is, and rested the seventh day: wherefore the Lord blessed the seventh day, and hallowed it.}}
	\par
	\leftskip=0em
    \textit{Lord, have mercy upon us, and incline our hearts to keep this law.}
    \par\noindent
    Honour thy father and thy mother;
    \par\noindent
    \leftskip=2em
	{\small{that thy days may be long in the land which the Lord thy God giveth thee.}}
	\par
	\leftskip=0em
    \textit{Lord, have mercy upon us, and incline our hearts to keep this law.}
    \par\noindent
    Thou shalt do no murder.\par
    \textit{Lord, have mercy upon us, and incline our hearts to keep this law.}
    \par\noindent
    Thou shalt not commit adultery.\par
    \textit{Lord, have mercy upon us, and incline our hearts to keep this law.}
\par\noindent
    Thou shalt not steal.\par
    \textit{Lord, have mercy upon us, and incline our hearts to keep this law.}
\par\noindent
    Thou shalt not bear false witness against thy neighbour.\par
    \textit{Lord, have mercy upon us, and incline our hearts to keep this law.}
    \par\noindent
    Thou shalt not covet.
    \par\noindent
    \leftskip=2em
	{\small{thy neighbour's house, thou shalt not covet thy neighbour's wife, nor his servant, nor his maid, nor his ox, nor his ass, nor any thing that is his.}}
	\par
	\leftskip=0em
	\textit{Lord, have mercy upon us, and write all these thy laws in our hearts, we beseech thee.}
	
%	\end{multicols}
\begin{rubric}
	Then may the Priest say the Summary of the Law, as followeth, in an audible voice.
\end{rubric}
\subsection{Summary of the Law}
\begin{center}
	{\textsc{Hear what our Lord Jesus Christ saith.}}
\end{center}

\lett{T}{hou} shalt love the Lord thy God with all thy heart, and with all thy soul, and with all thy mind. This is the first and great commandment. And the second is like unto it; Thou shalt love thy neighbour as thyself. On these two commandments hang all the Law and the Prophets.

\subsection{Ascending the Altar}

\begin{rubric}
    Then shall the Priest turn towards the Altar and ascend, saying secretly,%\blackrubric{Take away}.
\end{rubric}
\lett{T}{ake} away from us, we beseech thee, O Lord, our iniquities: that we may be worthy to enter with pure minds into the Holy of holies. Through Christ, our Lord. Amen.
\begin{rubric}
    Then, with hands joined upon the Altar, the Priest saith, bowing (kissing the Altar at the middle),% \blackrubric{We pray thee}.
\end{rubric}
\lett{W}{e} pray thee, O Lord, by the merits of thy Saints, * whose relics are here, and of all the Saints: that thou wouldest vouchsafe to pardon all my sins. Amen.
\begin{rubric}
    At a solemn Mass, the Celebrant, before reading the Introit, shall bless the incense saying: \blackrubric{Be thou bles {\ding{64}} sed by him in whose honour thou shalt be burnt. Amen.}
\end{rubric}
\begin{rubric}
	Receiving the thurible from the Deacon, he censeth the Altar, saying nothing. Then the Deacon taketh the thurible from the Celebrant and censeth him only.
\end{rubric}
\begin{rubric}
%Added clarity that the Priest should be in the middle of the Altar for the Kyrie.
    Then shall the Priest move to the Epistle side and, signing himself with the sign of the Cross, begin the Introit: which finished, at the middle of the Altar, with joined hands, he saith alternately with the Ministers the \blackrubric{Kyrie,}
\end{rubric}

%MANUAL ADJUSTMENT:
\vspace{-0.75\baselineskip}

\subsection{Kyrie}
\elcoldent{℣. Lord, have mercy upon us.}{℣. Kyrie, eléison.}
\elcoldent{℟. Lord, have mercy upon us.}{℟. Kyrie, eléison.}
\elcoldent{℣. Lord, have mercy upon us.}{℣. Kyrie, eléison.}
\elcoldent{℟. Christ, have mercy upon us.}{℟. Christe, eléison.}
\elcoldent{℣. Christ, have mercy upon us.}{℣. Christe, eléison.}
\elcoldent{℟. Christ, have mercy upon us.}{℟. Christe, eléison.}
\elcoldent{℣. Lord, have mercy upon us.}{℣. Kyrie, eléison.}
\elcoldent{℟. Lord, have mercy upon us.}{℟. Kyrie, eléison.}
\elcoldent{℣. Lord, have mercy upon us.}{℣. Kyrie, eléison.}

%MANUAL ADJUSTMENT:
\vspace{-0.15\baselineskip}

\subsection{Gloria in excelsis}

%MANUAL ADJUSTMENT:
\vspace{-0.1\baselineskip}

\begin{rubric}
    Then the Priest shall, at the middle of the Altar, extend and join his hands and---bowing his head slightly---say, if it is to be said,
\end{rubric}
\elcol{\lett{G}{lory} be to God on high, and on earth peace, good will towards men. We praise thee, we bless thee, we worship thee,\margcolrub{1}{Bow head.} we glorify thee, we give thanks\margcolrub{2}{Bow head.} to thee for thy great glory, O Lord God, heavenly King, God the Father Almighty.\par
	O Lord, the only-begotten Son, Jesus Christ; O Lord God, Lamb of God, Son of the Father, that takest away the sins of the world, have mercy upon us. Thou that takest away the sins of the world, receive our prayer.\margcolrub{3}{Bow head.} Thou that sittest at the right hand of God the Father, have mercy upon us.\par
	For thou only art holy; thou only art the Lord; thou only, O Christ, with the Holy Ghost, {\ding{64}} art most high in the glory of God the Father. Amen.}
    {\lett{G}{l\smash{ó}ria} in excélsis Deo. Et in terra pax homínibus bon{\ae} voluntátis. Laudámus te. Benedícimus te. Adorámus te.\marglatincolrub{1}{Bow head.} Glorificámus te. Grátias ágimus tibi\marglatincolrub{2}{Bow head.} propter magnam glóriam tuam. Dómine Deus, Rex c{\ae}léstis, Deus Pater omnípotens.\par
    Dómine Fili unigénite, Jesu Christe. Dómine Deus, Agnus Dei, Fílius Patris. Qui tollis peccáta mundi, miserére nobis. Qui tollis peccáta mundi, súscipe deprecatiónem nostram.\marglatincolrub{3}{Bow head.} Qui sedes ad déxteram Patris, miserére nobis.\par
    Quóniam tu solus Sanctus. Tu solus Dóminus. Tu solus Altíssimus, Jesu Christe. Cum Sancto Spíritu {\ding{64}} in glória Dei Patris. Amen.}

\begin{rubric}
    Then shall the Priest kiss the Altar in the middle, turn to the People, extend his hands, and say,
\end{rubric}

\elcol{℣. The Lord be with you.}{℣. Dóminus vobíscum.}
\elcol{℟. And with thy spirit.}{℟. Et cum spíritu tuo.}
\elcol{℣. Let us pray.}{℣. Orémus.}

%MANUAL ADJUSTMENT:
\vspace{-2ex}

\subsection{Collect}
\begin{rubric}
	Then shall the Priest move to the Epistle corner of the Altar and pray the Collect(s) of the Day.
\end{rubric}

%MANUAL ADJUSTMENT:
\vspace{-4ex}

\subsection{Epistle}
\begin{rubric}
	The Epistle for the Day shall then be read by the Subdeacon, first saying, \blackrubric{The Epistle is written in the \emph{--} Chapter of \emph{--}, beginning at the \emph{--} Verse.}
\end{rubric}
\begin{rubric}
	The Epistle ended, he shall say, \blackrubric{Here endeth the Epistle,} the people responding, \blackrubric{Thanks be to God.}
\end{rubric}
\begin{rubric}
    The Gradual and Alleluia (or Tract) and (if provided) Sequence is here chanted by the Choir.
\end{rubric}

%MANUAL ADJUSTMENT:
\vspace{-4ex}

\subsection{Gospel}
%\begin{multicols}{2}
\begin{rubric}
    These being ended, if it be a Solemn Mass, the Deacon shall then place the book of the Gospels on the middle of the Altar, and the Celebrant shall bless the incense with the \blackrubric{Be thou blessed,} as above.
\end{rubric}
\begin{rubric}
	Then shall the Deacon kneeling before the Altar (or the Priest only bowing in the midst of the Altar) say with joined hands,% the \blackrubric{Cleanse my heart} continuing to the \blackrubric{Bid, sir}, with the priest saying \blackrubric{The Lord be in thy}.
\end{rubric}

\lett{C}{leanse} my heart and my lips, O almighty God, who didst cleanse the lips of the prophet Isaiah with a live coal: so of thy gracious mercy vouchsafe to cleanse me, that I may worthily proclaim thy holy Gospel. Through Christ, our Lord. Amen.

\begin{paracol}{2}

\sloppy

\begin{inhead}
	Deacon
\end{inhead}
\par\noindent
	\emph{Deacon.} Bid, sir, a blessing.\par\noindent
	\emph{Priest.} The Lord be in thy heart and on thy lips: that thou mayest worthily and fitly proclaim his Gospel: In the name of the Father, and of the Son, {\ding{64}} and of the Holy Ghost. Amen.

\switchcolumn

\begin{inhead}
	Priest
\end{inhead}
	\par\noindent
	\emph{Priest.} Bid, Lord, a blessing.\par\noindent
	\emph{Priest.} The Lord be in my heart and on my lips: that I may worthily and fitly proclaim his Gospel. 
	
	\fussy
	
\end{paracol}


\begin{rubric}
    Having received the blessing, the Deacon kisseth the hand of the Celebrant. And going with the other Ministers, with the incense and the lights, to the place of the Gospel, he standeth with joined hands.%, saying what is below.
\end{rubric}
%Unnecessary rubric:
%\begin{rubric}
%    If, however, the Priest celebrate without Deacon and Subdeacon, when the book has been carried to the other corner of the Altar, he boweth in the middle, and with joined hands, saith the \blackrubric{Cleanse my heart} and the \blackrubric{Bid, Lord} and \blackrubric{The Lord be in my}.
%\end{rubric}

%\end{multicols}
\begin{rubric}
	Then, the People standing, the Deacon, turning to the Book, saith with joined hands,
\end{rubric}
\elcol{℣. The Lord be with you.}{℣. Dóminus vobíscum.}
\elcol{℟. And with thy spirit.}{℟. Et cum spíritu tuo.}
\elcol{℣. The Beginning (or, Continuation) {\ding{66}} of the Holy Gospel according to \textit{N.}

℟. Glory be to thee, O Lord.}{℣. Inítium (vel, Sequéntia) {\ding{66}} sancti Evangélii secúndum \textit{N.}

℟. Glória tibi, Dómine.}

\begin{rubric}
    Then shall the Deacon sign the book with the thumb of his right hand at the beginning of the Gospel text which he is to read, then himself on the forehead, the mouth, and the breast: and while the Ministers respond, he censeth the book thrice, then readeth the Gospel with joined hands.
\end{rubric}
\begin{rubric}
    At the end of the Gospel, the Ministers respond,
\end{rubric}
\elcol{℟. Praise be to thee, O Christ.}{℟. Laus tibi, Christe.}
\begin{rubric}
    Then shall the Subdeacon carry the book to the Priest, who kisseth the Gospel text, saying: \blackrubric{Through the words of the Gospel may our sins be blotted out.} Then the Priest is censed by the Deacon.\par
    \textsc{Note,} In Masses of the Dead, \blackrubric{Cleanse} is said, but a blessing is not asked, lights are not carried, and the Celebrant doth not kiss the book.
\end{rubric}
\begin{rubric}
    Then, in the middle of the Altar, extending, raising, and joining his hands, the Priest shall say, if it is to be said, \blackrubric{I believe in one God,} proceeding with joined hands.
\end{rubric}
\subsection{Nicene Creed}\label{NiceneCreed}
\elcol{\lett{I}{believe} in one God\margcolrub{4}{Bow head to Cross.} the Father Almighty, Maker of heaven and earth, And of all things visible and invisible:\par
    And in one Lord Jesus Christ,\margcolrub{5}{Bow head to Cross.} the only-begotten Son of God; Begotten of his Father before all worlds, God of God, Light of Light, Very God of very God; Begotten, not made; Being of one substance with the Father; By whom all things were made: Who for us men and for our salvation came down from heaven, \inrub{Everyone genuflects.} And was incarnate by the Holy Ghost of the Virgin Mary, And was made man: \inrub{Everyone rises.} And was crucified also for us under Pontius Pilate; He suffered and was buried: And the third day he rose again according to the Scriptures: And ascended into heaven, And sitteth on the right hand of the Father: And he shall come again, with glory, to judge both the quick and the dead; Whose kingdom shall have no end.\par
    And I believe in the Holy Ghost, The Lord, and Giver of Life, Who proceedeth from the Father; Who with the Father and the Son together is worshiped\margcolrub{6}{Bow head to Cross.} and glorified; Who spake by the Prophets: And I believe one, holy, catholic, and apostolic Church: I acknowledge one Baptism for the remission of sins: And I look for the Resurrection of the dead: {\ding{64}} And the Life of the world to come. Amen.}
    {\lett{C}{redo} in unum Deum,\marglatincolrub{4}{Bow head to Cross.} Patrem omnipoténtem, factórem c{\ae}li et terr{\ae}, visibílium ómnium et invisibílium.\par
    Et in unum Dóminum Jesum Christum,\marglatincolrub{5}{Bow head to Cross.} Fílium Dei unigénitum. Et ex Patre natum ante ómnia sǽcula. Deum de Deo, lumen de lúmine, Deum verum de Deo vero. Génitum, non factum, consubstantiálem Patri: per quem ómnia facta sunt. Qui propter nos hómines et propter nostram salútem descéndit de c{\ae}lis. \inrub{Everyone genuflects.} Et incarnátus est de Spíritu Sancto ex María Vírgine: Et homo factus est. \inrub{Everyone rises.} Crucifíxus étiam pro nobis: sub Póntio Piláto passus, et sepúltus est. Et resurréxit tértia die, secúndum Scriptúras. Et ascéndit in c{\ae}lum: sedet ad déxteram Patris. Et íterum ventúrus est cum glória judicáre vivos et mórtuos: cujus regni non erit finis.\par
    Et in Spíritum Sanctum, Dóminum et vivificántem: qui ex Patre procédit. Qui cum Patre et Fílio simul adorátur\marglatincolrub{6}{Bow head to Cross.} et conglorificátur: qui locútus est per Prophétas. Et unam sanctam cathólicam et apostólicam Ecclésiam. Confíteor unum baptísma in remissiónem peccatórum. Et exspécto resurrectiónem mortuórum. {\ding{64}} Et vitam ventúri s{\ae}culi. Amen.}
\begin{rubric}
Then shall be declared unto the People what Holy-days, or Fasting-days, are in the week following to be observed; and (if occasion be) shall Notice be given of the Communion, and of the Banns of Matrimony, and other matters to be published.
\end{rubric}

\begin{rubric}
Here, or immediately after the Creed, may be said the Bidding Prayer, or other authorized Prayers and intercessions.
\end{rubric}

\begin{rubric}
	Then followeth the Sermon.
\end{rubric}

\begin{rubric}
    Then shall the Priest ascend the Altar, kiss it, turn towards the People, and say,
\end{rubric}
\elcol{℣. The Lord be with you.}{℣. Dóminus vobíscum.}
\elcol{℟. And with thy spirit.}{℟. Et cum spíritu tuo.}
\elcol{℣. Let us pray.}{℣. Orémus.}

%MANUAL ADJUSTMENT:
\clearpage
\subsection{Offertory}

%\begin{multicols}{2}    
\begin{rubric}
    The Priest then saith the Offertory Verse.\par
    \textsc{Note,} At High Mass, a hymn may be sung while the Priest prepareth the Oblations in the Offertory, and while the Collection is taken, if there be one.
\end{rubric}

\subsubsection{Offering the Bread}
\begin{secrubric}
    Then shall the Deacon present the Paten with the Bread to the Celebrant, which he then offereth, saying,%: \blackrubric{Receive, O holy Father}.
\end{secrubric}
\lett{R}{eceive,} O holy Father, almighty everlasting God, this spotless host, which I, thine unworthy servant, offer unto thee, my living and true God, for my numberless sins, offences, and negligences, and for all who stand around, as also for all faithful Christians, both living and departed: that to me and to them it may avail for salvation unto life eternal. Amen.
\begin{rubric} 
    Then, making a cross with the same Paten, he shall place the Bread upon the Corporal (or the Paten with the Bread upon the Corporal).
\end{rubric}

\subsubsection{Offering the Wine}
\begin{secrubric}
	Then the Deacon shall minister the Wine, and the Subdeacon the Water into the Cup. The Priest (blessing with the sign of the Cross the Water to be mixed in the Cup, unless it be a Requiem Mass) shall say,%: \blackrubric{O God, who didst wonderfully}.
\end{secrubric}
\lett{O}{God,} who didst wonderfully create, and yet more wonderfully renew the dignity of man's nature: grant that by the mystery of this water and wine we may be made partakers of his divinity, who vouchsafed to share our humanity, Jesus Christ thy Son our Lord: Who liveth and reigneth with thee in the unity of the Holy Ghost, God: throughout all ages world without end. Amen.
\begin{rubric}
    Then shall the Priest receive the Cup and offer it, saying,% \blackrubric{We offer unto thee}.
\end{rubric}
\lett{W}{e} offer unto thee, O Lord, the cup of salvation, humbly beseeching thy mercy: that in the sight of thy divine majesty it may ascend as a sweet-smelling savour for our salvation, and for that of the whole world. Amen.

%MANUAL ADJUSTMENT:
\clearpage
\subsubsection{Offering the Gifts}
\begin{secrubric}
    Then shall he make the sign of the Cross with the Cup, place it upon the Corporal, and cover it with the Pall. Then, with hands joined upon the Altar, he saith, bowing slightly,% \blackrubric{In a humble spirit}.
\end{secrubric}
\lett{I}{n} a humble spirit, and with a contrite heart, may we be accepted of thee, O Lord: and so let our sacrifice be offered in thy sight this day, that it may be pleasing unto thee, O Lord God.
\begin{rubric}
    Then shall the Priest---standing upright---extend his hands, raise them, and join them; and lifting his eyes to heaven and lowering them immediately, he saith,% \blackrubric{Come, O Sanctifier}.
\end{rubric}
\lett{C}{ome,} O Sanctifier almighty, everlasting God, and bl {\ding{64}} ess this sacrifice, prepared for thy holy name.

\lettspace

\subsubsection{Censing the Altar}
\begin{secrubric}
    If he be celebrating solemnly, the Priest blesseth incense, saying,% \blackrubric{Through the intercession of blessed Michael the Archangel standing at the right hand of the altar of incense, and of all his elect, may the Lord vouchsafe to bl {\ding{64}} ess this incense, and receive it for a sweet-smelling savour. Through Christ, our Lord. Amen.}.
\end{secrubric}
\lett{T}{hrough} the intercession of blessed Michael the Archangel standing at the right hand of the altar of incense, and of all his elect, may the Lord vouchsafe to bl {\ding{64}} ess this incense, and receive it for a sweet-smelling savour. Through Christ, our Lord. Amen.
\begin{rubric}
    Receiving the Thurible from the Deacon, he censeth the Oblations, in the manner prescribed in the traditional manner, saying,%: \blackrubric{May this incense which thou hast blessed, ascend unto thee, O Lord: and let thy mercy descend upon us.}.\par
\end{rubric}
\lett{M}{ay} this incense which thou hast blessed, ascend unto thee, O Lord: and let thy mercy descend upon us.

\lettspace

\begin{rubric}
    Then the Priest censeth the Altar, saying,%: \blackrubric{Let my prayer}.\par
\end{rubric}
\lett{L}{et} my prayer, O Lord, be set forth in thy sight as the incense: and let the lifting up of my hands be an evening sacrifice. Set a watch, O Lord, before my mouth, and keep the door of my lips: O let not mine heart be inclined to any evil thing, let me not be occupied in ungodly works.
\begin{rubric}
    While he returneth the Thurible to the Deacon, he saith,% \emph{The Lord kindle}.\par
\end{rubric}
\lett{T}{he} Lord kindle in us the fire of his love, and the flame of everlasting charity. Amen.

\lettspace

\begin{rubric}
    Then the Priest is censed by the Deacon, and afterwards the others in order.
\end{rubric}

\subsubsection{Washing the Hands}
\begin{secrubric}
    Meanwhile, the Priest washeth his hands, saying,% \blackrubric{I will wash \&c}.
\end{secrubric}
\lett{I}{will} wash my hands in innocency, O Lord: and so will I go to thine altar: That I may shew the voice of thanksgiving, and tell of all thy wondrous works. Lord, I have loved the habitation of thy house and the place where thine honour dwelleth. O shut not up my soul with the sinners, nor my life with the blood-thirsty: In whose hands is wickedness: and their right hand is full of gifts. But as for me, I will walk innocently: O deliver me, and be merciful unto me. My foot standeth right: I will praise the Lord in the congregations. Glory be to the Father, and to the Son, and to the Holy Ghost. As it was in the beginning, is now, and ever shall be, world without end. Amen.

%\begin{rubric}
%	\textsc{Note,} In Requiem Masses and during Passiontide, \blackrubric{Glory be} is omitted.
%\end{rubric}

%MANUAL ADJUSTMENT:
\vspace{-1.5ex}

\subsubsection{Offering the Oblation}

\begin{secrubric}
    Then bowing slightly, in the middle of the Altar with hands joined upon it, he saith,% \blackrubric{Receive, O holy Trinity}.
\end{secrubric}
\lett{R}{eceive,} O holy Trinity, this oblation which we offer unto thee in memory of the Passion, Resurrection, and Ascension of Jesus Christ, our Lord: and in honour of Blessed Mary ever Virgin, of Blessed John Baptist, of the Holy Apostles Peter and Paul, of these, and of all the Saints: that it may avail to their honour, and for our salvation: and may they, whose memory we celebrate on earth, vouchsafe to intercede for us in heaven. Through the same Christ, our Lord. Amen.
%\end{multicols}

\begin{rubric}
    Then shall the Priest kiss the Altar and, turning to the People, extend and join his hands, and say, raising his voice a little,
\end{rubric}

\elcol{℣. Pray, brethren: that my sacrifice and yours may be acceptable to God the Father Almighty.}{℣. Oráte, fratres: ut meum ac vestrum sacrifícium acceptábile fiat apud Deum Patrem omnipoténtem.}

%Anglican Missal (more familiar):
\elcol{℟. The Lord receive this sacrifice at thy hands, to the praise and glory of his Name, both to our benefit, and that of all his holy Church}{℟. Suscípiat Dóminus sacrifícium de mánibus tuis ad laudem et glóriam nominis sui, ad utilitátem quoque nostram, totiúsque Ecclési{\ae} su{\ae} sanct{\ae}.}
\begin{rubric}
    The Priest saith in a low voice: \blackrubric{Amen.}
\end{rubric}

%MANUAL ADJUSTMENT:
\vspace{-1.5ex}

\subsubsection{Secret Prayers}

\begin{secrubric}
	Then shall the Priest, with hands extended, immediately (without \blackrubric{Let us pray}), add the Secret Prayers.\par
	\textsc{Note,} When these are ended, he saith in a loud voice,
\end{secrubric}\par\noindent
℣. Throughout all ages, world without end.\par
℟. Amen.

\clearpage
\begin{rubric}
    Then shall the Deacon say,
\end{rubric}
\centerline{Let us pray for the whole state of Christ's Church.}
\begin{rubric}
    The Priest then saith, with hands extended, the following.
\end{rubric}
\lett{A}{lmighty} and everliving God, who by thy holy Apostle hast taught us to make prayers, and supplications, and to give thanks for all men; We humbly beseech thee most mercifully to accept our (alms and) oblations, and to receive these our prayers, which we offer unto thy Divine Majesty; beseeching thee to inspire continually the Universal Church with the spirit of truth, unity, and concord: And grant that all those who do confess thy holy Name may agree in the truth of thy holy Word, and live in unity and godly love.

\lett{W}{e} beseech thee also, so to direct and dispose the hearts of all Christian Rulers, that they may truly and impartially administer justice, to the punishment of wickedness and vice, and to the maintenance of thy true religion, and virtue.

%Mention of bishops imported from 1929 Scottish BCP:
\lett{G}{ive} grace, O heavenly Father, to all Bishops and other Ministers, and especially to thy servant \emph{N.} our Metropolitan (and \emph{N.} our Bishop), that \textit{he} may, both by \textit{his} life and doctrine, set forth thy true and lively Word, and rightly and duly administer thy holy Sacraments.

\lett{A}{nd} to all thy People give thy heavenly grace; and especially to this congregation here present; that, with meek heart and due reverence, they may hear, and receive thy holy Word; truly serving thee in holiness and righteousness all the days of their life.

\lett{A}{nd} we most humbly beseech thee, of thy goodness, O Lord, to comfort and succour all those who, in this transitory life, are in trouble, sorrow, need, sickness, or any other adversity.

%Mention of the Saints imported from 1929 Scottish BCP:
\lett{A}{nd} we also bless thy holy Name for all thy servants departed this life in thy faith and fear; beseeching thee to grant them continual growth in thy love and service, and to give us grace so to follow their good examples, chiefly the Blessed Virgin Mary, Mother of thy Son Jesus Christ our Lord and God, and the Holy Patriarchs, Prophets, Apostles, and Martyrs, that with them we may be partakers of thy heavenly kingdom.
\begin{rubric}
	 The Priest placeth both hands upon the Altar, outside the Corporal, and continueth,
\end{rubric}

%\lett{G}{rant} 
\par\noindent
Grant this, O Father, for Jesus Christ's sake,% our only Mediator and Advocate.\par
%\secondline{℟. Amen.\\}

\gregorioscore{resources/gabc/Mediator.gabc}


%This is a difficult situation to discern. The Prayer Book Traditions reflects the uncertainty as to where to place the Confession. The current American placement is reflective of Cranmer's Protestant theology and liturgical sense, since the origins come from placing the Confession close to the reception of Communion (which was immediately after the Words of Institution). The 1928 reflects a more catholic understanding of the liturgy, including the epiklesis and the other integral parts of the Eucharistic Sacrifice before the Reception of Holy Communion. This understanding was more present in the 1549 which placed the Confession in a congruous place as the Roman tradition: after the Sacrifice and before the Communion of the Faithful. While this is obscured here due to the Byzantine devotion imposed, this is only a temporary imposition and the overall shape of the liturgy (due to the Orthodox theology present here) is better preserved in its Catholic and even Prayer Book form by placing the Confession immediately before the `Ecce Agnus Dei', which is where the Second Confiteor would go in the Roman Rite.
%MANUAL ADJUSTMENT:
\clearpage
\subsection{Preface}
\begin{rubric}
    Then shall the Priest begin the Preface, with both hands placed apart on the Altar.
\end{rubric}
\elcol{℣. The Lord be with you.}{℣. Dóminus vobíscum.}
\elcol{℟. And with thy spirit.}{℟. Et cum spíritu tuo.}
\elcol{℣. Lift up your hearts.\margcolrub{7}{He raises his hands a little.}}{℣. Sursum corda.\marglatincolrub{7}{He raises his hands a little.}}
\elcol{℟. We lift them up unto the Lord.}{℟. Habémus ad Dóminum.}
\elcol{℣. Let us give thanks unto our Lord God.\margcolrub{8}{He joins his hands before his breast, \& bows his head.}}{℣. Grátias agámus Dómino, Deo nostro.\marglatincolrub{8}{He joins his hands before his breast, \& bows his head.}}
\elcol{℟. It is meet and right so to do.}{℟. Dignum et justum est.}
\begin{rubric}
	The Priest separateth his hands and continueth with the following.
\end{rubric}
\lett{I}{t} is very meet, right, and our bounden duty, that we should at all times, and in all places, give thanks unto thee, O Lord, Holy Father, Almighty, Everlasting God.
\begin{rubric}
    The Proper Preface is then here said (p. \pageref{prefaces})---if there be one---concluding with the following.
\end{rubric}
\label{PrefaceEnding}
\lett{T}{herefore} with Angels and Archangels, and with all the company of heaven, we laud and magnify thy glorious Name; evermore praising thee, and saying:\margrubleft{9}{He joins his hands \& bows.}


\subsection{Sanctus}
\elcol{\lett{H}{oly, Holy, Holy,} Lord God of hosts, Heaven and earth are full of thy glory: Glory be to thee, O Lord Most High. Blessed\margcolrub{10}{He stands upright \& signs himself.} {\ding{64}} is he that cometh in the Name of the Lord. Hosanna in the Highest.}
{\lett{S}{anctus, Sanctus, Sanctus,} Dóminus, Deus Sábaoth. Pleni sunt c{\ae}li et terra glória tua. Glória tibi, Dómine altíssime. Benedíctus,\marglatincolrub{10}{He stands upright \& signs himself.} {\ding{64}} qui venit in nómine Dómini. Hosánna in excélsis.}

\vfill

	\begin{figure}[H]
		\centering
		\includegraphics[scale=.25]{tailpiece/Crown.eps}
	\end{figure}

\clearpage

\begin{figure}
    \centering
    \includegraphics[width=98mm, height=135mm]{HolyMass/Canon.eps}
\end{figure}

\thispagestyle{empty}
\clearpage
%\reversemarginpar
\subsection{Canon of the Mass}
\begin{rubric}
    The Priest, extending, raising somewhat, and joining his hands, raising his eyes towards heaven, and immediately lowering them, shall begin the Canon, as followeth.
\end{rubric}
\lett{A}{ll} glory be to thee, Almighty God, our heavenly Father, for that thou, of thy tender mercy, didst give thine only Son Jesus Christ to suffer death upon the Cross for our redemption; who (by his own oblation of himself once offered) made a full, perfect, and sufficient sacrifice, oblation, and satisfaction, for the sins of the whole world; and did institute, and in his holy Gospel command us to continue, a perpetual memory of that his precious death and sacrifice, until his coming again:
\subsubsection{Words of Institution}

\par\noindent
  \begin{minipage}[t]{0.7\textwidth}% Match width to wrapfigure
	\lett{F}{or} in the night in which he was betrayed, \textit{(a)} he took Bread; \textit{(b)} and when he had given {\ding{64}} thanks, he brake it, and gave it to his disciples, saying, Take, eat, \textit{(c)}
	
	\begin{center}
		\large{This is my Body, which is given for you;\\Do this in remembrance of me.}
	\end{center}
  \end{minipage}
\hfill
  \begin{minipage}[t]{0.25\textwidth}\setlength{\parindent}{10pt}% Match width to wrapfigure
  	{\small\indent
    \textit{(a) He takes \& holds the Bread.}\par\indent
    \textit{(b) He looks up to heaven \& bows his head.}\par\indent
    \textit{(c) He holds the Bread between both of his thumbs \& forefingers.}
    }
  \end{minipage}

\vspace{1ex}

%\begin{center}
%\large{This is my Body, which is given for you;\\Do this in remembrance of me.}
%\end{center}
\begin{secrubric}
    Then shall the Priest immediately genuflect, briefly elevate the Bread for the People to see, replace it upon the Corporal (or Paten), genuflect again, and then immediately continue,\par
    \textsc{Note,} From henceforth, the Priest doth not disjoin his forefingers \& thumbs until the ablutions.
\end{secrubric}
\par\noindent
  \begin{minipage}[t]{0.7\textwidth}
  \noindent% Match width to wrapfigure
	Likewise, after supper, \textit{(d)} he took the Cup; \textit{(e)} and when he had given {\ding{64}} thanks, he gave it to them, saying, Drink ye all of this; for \textit{(f)}
	
	\begin{center}
		\large{This is my Blood of the New Testament, which is shed for you, and for many, for the remission of sins;}
	\end{center}
  \end{minipage}
\hfill
  \begin{minipage}[t]{0.25\textwidth}\setlength{\parindent}{10pt}% Match width to wrapfigure
  	{\small\indent
    \textit{(d) He uncovers the Cup, holds it with both hands, \& bows his head.}\par\indent
    \textit{(e) He holds the Cup in his left hand.}\par\indent
    \textit{(f) He holds the Cup slightly raised.}
    }
  \end{minipage}
  
  \vspace{1ex}
  
  
  \begin{secrubric}
    Then shall the Priest place the Cup back upon the Corporal and immediately say,
\end{secrubric}\par\noindent
    \centerline{Do this, as oft as ye shall drink it, in remembrance of me.}
\begin{rubric}
    Then shall the Priest immediately genuflect, briefly elevate the Cup, replace the Cup upon the Corporal, cover it, genuflect again, and continue with hands extended.
\end{rubric}
\subsubsection{Recollection}
%Added signs of the Cross from American \& Anglican Missal:
\lett{W}{herefore} O Lord and heavenly Father, according to the institution of thy dearly beloved Son our Saviour Jesus Christ, we, thy humble servants, do celebrate and make here before thy Divine Majesty, with these thy holy {\ding{64}} gifts, which we now offer unto thee, the memorial thy Son hath commanded us to make; having in remembrance his blessed passion and precious death, his mighty resurrection and glorious ascension; rendering unto thee most hearty thanks for the innumerable benefits procured unto us by the same.

\subsubsection{Invocation}
\begin{rubric}
	Then shall the Priest uncover the Cup, bow profoundly, join his hands upon the Altar, and continue,
\end{rubric}
%Epiklesis of St. James through the 1929 Scottish BCP:
%Crosses supplied from the Anglican Missal.
%Moving the Pall from precedent in the LAP Missal:
	\par\noindent
  \begin{minipage}[t]{0.71\textwidth}\noindent
	\lett{A}{nd} we thine unworthy servants beseech thee, most merciful Father, to hear us, \textit{(g)} and to send thy Holy Spirit upon us and upon the~{\ding{64}}~se thy
  \end{minipage}
  \hfill
  \begin{minipage}[t]{0.25\textwidth}
  	{\small\itshape
  	(g) He stands upright and imposes his hands over the Gifts.}
  \end{minipage}
  
  \vspace{0.5ex}
\par\noindent
gifts and creatures of Br~{\ding{64}}~ead and W~{\ding{64}}~ine, that, being blessed and hallowed by his life-giving power, they may become the Bo~{\ding{64}}~dy and Bl~{\ding{64}}~ood of thy most dearly beloved Son, to the end that all who shall receive the same may be sanctified both in body and soul, and preserved unto everlasting life.


\begin{rubric}
	Then shall the Priest cover the Cup, genuflect, and continue, with hands extended,
\end{rubric}
%From Scottish 16-something.
%\lett{H}{ear} us, O merciful Father, we most humbly beseech thee, and of thy almighty goodness, vouchsafe\margrub{He imposes his hands over the Bread \& Wine, right hand over left.} so to bl {\ding{64}} ess and sanct {\ding{64}} ify, with thy Word and Holy Spirit, these thy gifts and creatures of Bread and Wine, that they may %be unto us
%become the Bo {\ding{64}} dy and Bl {\ding{64}} ood of thy most dearly beloved Son; so that we receiving them according to thy Son our Saviour Jesus Christ's holy institution, in remembrance of his death and passion, may be partakers of the same his most precious Body and Blood.
%Interpolation of the Liturgy of St. Tikhon \& 1928 BCP.
%Grant that we, receiving them according to thy Son our Saviour Jesus Christ's holy institution, in remembrance of his Death and Passion, may be partakers of his most blessed Body and Blood.

%\lett{A}{nd} we most humbly beseech thee, O merciful Father, to hear us; and, of thy almighty goodness, vouchsafe to bless and sanctify, with thy Word and Holy Spirit, these thy gifts and creatures of bread and wine; that we, receiving them according to thy Son our Saviour Jesus Christ’s holy institution, in remembrance of his death and passion, may be partakers of his most blessed Body and Blood.

%\lett{A}{nd} we most humbly beseech thee, O merciful Father, to hear us: and, of thy almighty goodness, vouchsafe to send down thy Holy Spirit upon these thy gifts and creatures of Bread and Wine, that they may be changed into the Bo {\ding{64}} dy and Bl {\ding{64}} ood of thy most dearly beloved Son.\margrub{He extends his hands.} Grant that we, receiving them according to thy Son our Saviour Jesus Christ's holy institution, in remembrance of his Death and Passion, may be partakers of his most blessed Body and Blood.
\par

\lett{A}{nd} we earnestly desire thy fatherly goodness, mercifully to accept this our Sacrifice of Praise and Thanksgiving; most humbly beseeching thee to grant that, by the merits and death of thy Son Jesus Christ, and through faith in his Blood, we, and all thy whole Church, may obtain remission of our sins, and all other benefits of his Passion.
\subsubsection{Commemoration for the Departed}
\lett{R}{emember} also, O Lord, thy servants and handmaids \emph{N.} and \emph{N.}, who have gone before us with the sign of faith, and rest in the sleep of peace.

\lettspace

\begin{rubric}
    Then shall the Priest join his hands, pray for the dead he hath in mind, then extend his hands and say,
\end{rubric}\par\noindent
To them, O Lord, and to all that rest in Christ, we beseech thee to grant a place of refreshing, light, and peace.
%The English & American Missals lacks this addition, being only in the Anglican Missal tradition:
%\lett{A}{nd} vouchsafe to grant some part and fellowship with thy holy Apostles and Martyrs: with John, Stephen, Matthias, Barnabas, Ignatius, Alexander, Marcellinus, Peter, Felicitas, Perpetua, Agatha, Lucy, Agnes, Cecilia, Anastasia, and with all thy Saints: within whose fellowship we beseech thee admit us.
%\lett{R}{emember} Lord, also the souls of thy servants and handmaidens, which are gone before us with the mark of faith, and rest in the sleep of peace. \inrub{The Priest now prays for the souls of the dead.} We beseech thee, O Lord, that unto them, and unto all such as rest in Christ, thou wilt grant a place of refreshing, of light, and of peace. And vouchsafe to give unto us some portion and fellowship with thy holy Apostles and Martyrs; with John, Stephen, Matthias, Barnabas, Ignatius, Alexander, Marcellinus, Peter, Felicitas, Perpetua, Agatha, Lucia, Agnes, Cecilia, Anastasia, and with all thy Saints; within whose fellowship we beseech thee to admit us.





\subsubsection{Oblation}
	\lett{A}{nd} here we offer and present unto thee, O Lord, our selves, our souls and bodies, to be a reasonable, holy, and living sacrifice unto thee; humbly beseeching thee, \inrub{He kisses the Altar} that we, and all others who shall be partakers of this Holy Communion, may worthily receive the most precious
	\par\noindent
  \begin{minipage}[t]{0.7\textwidth}\setlength{\parindent}{15pt}\noindent
  Bo~{\ding{64}}~dy and Bl~{\ding{64}}~ood \textit{(h)} of thy Son Jesus Christ, be filled with thy grace and heavenly benediction, and made one body with him, that he may dwell in us, and we in him.
  \end{minipage}
  \hfill
  \begin{minipage}[t]{0.25\textwidth}%\setlength{\parindent}{10pt}
  \vspace{0.25ex}
  	{\small\itshape%\indent
    (h) After signing the Gifts, he signs himself.}
  \end{minipage}
  
  \vspace{0.5ex}
  \par\noindent
  
	
	\lett{A}{nd} although we are unworthy, through our manifold sins, to offer unto
	
		\par\noindent
  \begin{minipage}[t]{0.7\textwidth}%\setlength{\parindent}{15pt}\noindent
  
  \leftskip=3.25em
  
  thee any sacrifice; yet we beseech thee to accept this our bounden duty and service; not weigh-
  
	\leftskip=0em
  
  ing our merits, but pardoning our offences. \inrub{He joins his hands \& bows his head profoundly.}
  
  \lett{T}{hrough} Jesus Christ our Lord; \textit{(i)} by~{\ding{64}}~whom, and with wh~{\ding{64}}~om, in the unity of the Holy~{\ding{64}}~Ghost, \textit{(j)} all ho~{\ding{64}}~nour and gl~{\ding{64}}~ory be unto thee,
  \end{minipage}
  \hfill
  \begin{minipage}[t]{0.25\textwidth}\setlength{\parindent}{10pt}
    \vspace{1ex}
  	{\small\itshape
  	(i) He genuflects \& makes the sign the Cross with the Host over the Cup thrice, from lip to lip.\par\indent
  	(j) He makes the sign of the Cross twice with the Host between the Cup \& his breast.}
  \end{minipage}

\begin{rubric}
	The Priest then uncovereth the Cup, elevateth the Host and Cup to the height of his breast then replacing them on the Altar, covereth the Cup, genuflecteth, then joineth his hands.
\end{rubric}

\gregorioscore{resources/gabc/CanonEnd.gabc}
%\par\noindent
%O Father Almighty, world without end.

%℟. Amen.


\subsection{Lord's Prayer}\par\noindent
	Let us pray. And now, as our Saviour Christ hath taught us, we are bold to say, 

\begin{rubric}
	The Priest alone extendeth his hands while the Priest and Congregation say together,
\end{rubric}

\lett{O}{ur} Father, who art in heaven, Hallowed be thy Name. Thy kingdom come. Thy will be done, on earth, As it is in heaven. Give us this day our daily bread. And forgive us our trespasses, As we forgive those who trespass against us. And lead us not into temptation, But deliver us from evil. For thine is the kingdom, and the power, and the glory, for ever and ever. Amen.

\begin{rubric}
    Unless the Priest consecrated on the Paten, he takes the Paten between the fore and middle fingers of his right hand, and holds it upright upon the Altar.
\end{rubric}
%
\lett{D}{eliver} us, O Lord, we beseech thee, from all evils, past, present, and to come: and at the intercession of the blessed and glorious ever Virgin Mary, Mother of God, with thy blessed Apostles Peter and Paul, and with Andrew, and all the Saints, \inrub{The Priest signs himself---with the Paten, if he consecrated on the Corporal} favourably grant peace in our days: that by the help of thine availing mercy we may ever both be free from sin and safe from all distress.
\begin{rubric}
    If the Host be not already on the Paten, the Priest is to slide the Paten underneath the Host. He then shall uncover the Cup, genuflect, and break the Host in half over the Cup, saying,
\end{rubric}\par\noindent
Through the same Jesus Christ, thy Son our Lord.
\begin{rubric}
    The Priest placeth the half in his right hand upon the Paten. He then breaketh a particle off from the half in his left hand, saying,
\end{rubric}\par\noindent
Who liveth and reigneth with thee in the unity of the Holy Ghost, God.
\begin{rubric}
    Then shall the Priest join the other half, which he holdeth in his left hand, to the half laid upon the Paten, and retaining the small particle in his right hand over the Cup, which he holdeth with his left by the knop below the Cup, say in an audible voice, or sing,
\end{rubric}\par\noindent
Throughout all ages world without end.\par
℟. Amen.

\begin{rubric}
	With the same particle, the Priest signs thrice over the Cup, saying,
\end{rubric}

\elcol{℣. The peace {\ding{64}} of the Lord be {\ding{64}} alway with {\ding{64}} you.}{℣. Pax {\ding{64}} Dómini sit {\ding{64}} semper vobís {\ding{64}} cum.}
\elcol{℟. And with thy spirit.}{℟. Et cum spíritu tuo.}

\begin{rubric}
    The Priest, putting the same particle into the Cup, saith, in the secret voice:
\end{rubric}
\lett{M}{ay} this commixture and consecration of the Body and Blood of our Lord Jesus Christ be to us who receive it unto everlasting life. Amen.

%MANUAL ADJUSTMENT:
\vspace{0.5\baselineskip}

\subsection{Agnus Dei}
\begin{rubric}
    Then shall the Priest cover the Cup, genuflect, and---bowing to the Sacrament---join his hands and strike his breast thrice, saying in an audible voice:
\end{rubric}
\elcol{\lett{O}{Lamb} of God, that takest away the sins of the world, have mercy upon us.}{\lett{A}{gnus} Dei, qui tollis peccáta mundi: miserére nobis.}
\elcol{O Lamb of God, that takest away the sins of the world, have mercy upon us.}{Agnus Dei, qui tollis peccáta mundi: miserére nobis.}
\elcol{O Lamb of God, that takest away the sins of the world, grant us thy peace.}{Agnus Dei, qui tollis peccáta mundi: dona nobis pacem.}
\begin{rubric}
	In Requiem Masses, \blackrubric{Have mercy upon us} is replaced by \blackrubric{grant them rest} and \blackrubric{grant us thy peace} by \blackrubric{grant them rest everlasting.}
\end{rubric}
%This is the original location of the Prayer Book Confession, mirroring the placement of the Roman Confession. However, in order to emphasise Reformed theology, the later English Prayer Books placed reception of Holy Communion to immediately after the Words of Institution, requiring the Confession to be placed before the Canon. Even though the American tradition developed a more Catholic Eucharistic rite, the Reformed placement of the Confession remained. Here, this is corrected in order to better express our Orthodox Catholic theology.
\subsection{Confession}
\begin{rubric}
    Then the Deacon turneth to the People and saith to those who come to receive the Holy Communion,
\end{rubric}
\lett{Y}{e} who do truly and earnestly repent you of your sins, and are in love and charity with your neighbours, and intend to lead a new life, following the commandments of God, and walking from henceforth in his holy ways; Draw near with faith, and take this holy Sacrament to your comfort; and make your humble confession to Almighty God, devoutly kneeling.
\begin{rubric}
Then shall this General Confession be made, in an audible voice, by the Priest, bowing, and all those who are minded to receive the Holy Communion, humbly kneeling.
\end{rubric}

\lett{A}{lmighty} God, Father of our Lord Jesus Christ, Maker of all things, Judge of all men; We acknowledge and bewail our manifold sins and wickedness, Which we, from time to time, most grievously have committed, By thought, word, and deed, Against thy Divine Majesty, Provoking most justly thy wrath and indignation against us. We do earnestly repent, And are heartily sorry for these our misdoings; The remembrance of them is grievous unto us; The burden of them is intolerable. Have mercy upon us, Have mercy upon us, most merciful Father; For thy Son our Lord Jesus Christ's sake, Forgive us all that is past; And grant that we may ever hereafter Serve and please thee in newness of life, To the honour and glory of thy Name; Through Jesus Christ our Lord. Amen.
\begin{rubric}
	Then shall the Priest (Bishop if he be present) stand up and, turning to the People, say,
\end{rubric}

\lett{A}{lmighty} God, our heavenly Father, who of his great mercy hath promised forgiveness of sins to all those who with hearty repentance and true faith turn unto him; Have mercy upon you; pardon {\ding{64}} and deliver you from all your sins; confirm and strengthen you in all goodness; and bring you to everlasting life; through Jesus Christ our Lord.

℟. Amen.

%MANUAL ADJUSTMENT:
\clearpage
\subsection{Comfortable Words}
\begin{rubric}
Then shall the Priest, facing the People, say,
\end{rubric}\noindent
\begin{center}
	\textsc{Hear what comfortable words our Saviour Christ saith unto all who truly turn to him.}
\end{center}
\par\noindent
Come unto me, all ye that travail and are heavy laden, and I will refresh you. \vr{Matt. 11:28}
\par\noindent
    So God loved the world, that he gave his only-begotten Son, to the end that all that believe in him should not perish, but have everlasting life. \vr{John 3:16}
    \par\noindent
    \begin{center}
		\textsc{Hear also what Saint Paul saith.}
	\end{center}
\par\noindent
    This is a true saying, and worthy of all men to be received, That Christ Jesus came into the world to save sinners. \vr{1 Tim. 1:15}

    \begin{center}
		\textsc{Hear also what Saint John saith.}
	\end{center}
    \par\noindent
    If any man sin, we have an Advocate with the Father, Jesus Christ the righteous; and he is the Propitiation for our sins. \vr{1 John 2:1-2}
\begin{rubric}
    Then the Priest, turning to the Altar, bowing with hands joined upon the Altar, saith in the secret voice,
\end{rubric}
    \lett{O}{Lord} Jesu Christ, who saidst to thine Apostles: Peace I leave with you, my peace I give unto you: regard not my sins but the faith of thy Church; and vouchsafe to grant her peace and unity according to thy will: Who livest and reignest God, throughout all ages, world without end. Amen.
\begin{rubric}
    If the \blackrubric{Pax} be given, the Priest is to kiss the Altar and---giving the \blackrubric{Pax}---say: \blackrubric{Peace be with thee.} with the response \blackrubric{And with thy spirit.}
\end{rubric}

\subsection{Prayer of Humble Access}
\begin{rubric}
    The Priest alone then saith, in an audible voice,
\end{rubric}
\lett{W}{e} do not presume to come to this thy Table, O merciful Lord, trusting in our own righteousness, but in thy manifold and great mercies. We are not worthy so much as to gather up the crumbs under thy Table. But thou art the same Lord, whose property is always to have mercy: Grant us therefore, gracious Lord, so to eat the flesh of thy dear Son Jesus Christ, and to drink his blood, that our sinful bodies may be made clean by his Body, and our souls washed through his most precious Blood, and that we may evermore dwell in him, and he in us.\par
℟. Amen.


\subsection{Communion of the Priest}
%\begin{multicols}{2}
\begin{rubric}
    Then the Priest shall say that which followeth,%\blackrubric{O Lord Jesu Christ} and \blackrubric{Let the partaking}.
\end{rubric}
\lett{O}{Lord} Jesu Christ, Son of the living God, who by the will of the Father, and the co-operation of the Holy Ghost, hast through thy death given life unto the world: deliver me by this thy most sacred Body and Blood from all mine iniquities, and from every evil: and make me ever to cleave unto thy commandments, and suffer me never to be separated from thee: Who with the same God the Father and the Holy Ghost livest and reignest God, world without end. Amen.
\lett{L}{et} the partaking of thy Body, O Lord Jesu Christ, which I, unworthy, presume to receive, turn not to my judgement and condemnation: but of thy goodness let it avail unto me for protection of mind and of body, that I may receive thy healing: Who livest and reignest with God the Father in the unity of the Holy Ghost, God, throughout all ages, world without end. Amen.
\begin{rubric}
    He then genuflecteth and saith: \blackrubric{I will receive the bread of heaven, and call upon the name of the Lord.}
\end{rubric}
\begin{rubric}
    Then, bowing slightly, the Priest shall take both parts of the Host between the thumb and forefinger of his left hand, and place the Paten between the same forefinger and middle finger, and striking his breast three times with his right hand, say thrice, devoutly and humbly, raising his voice slightly:
\end{rubric}\par\noindent
    Lord, I am not worthy, \inrub{Proceeding in the secret voice:} that thou shouldest come under my roof: but speak the word only, and my soul shall be healed.
\begin{rubric}
After signing himself with his right hand with the Host over the Paten, he saith,
\end{rubric}\par\noindent
The Body of our Lord Jesus Christ preserve my soul unto everlasting life. Amen.
\begin{rubric}
And bowing, he reverently taketh both parts of the Host. After consuming the Host, he is to place the Paten down upon the Corporal, and raising himself, join his hands, and remain still for a short time in meditation on the Most Holy Sacrament.
\end{rubric}
\begin{rubric}
	Then he shall uncover the Cup, genuflect, collect the fragments (if there be any), and cleanse the Paten over the Cup, saying meanwhile:
\end{rubric}
\par\noindent
What reward shall I give unto the Lord for all the benefits that he hath done unto me? I will receive the cup of salvation, and call upon the Name of the Lord. I will call upon the Lord which is worthy to be praised, so shall I be safe from mine enemies.
\begin{rubric}
    Taking the Cup in his right hand and signing himself with it, he saith,
\end{rubric}\par\noindent
The Blood of our Lord Jesus Christ preserve my soul unto everlasting life. Amen.
\begin{rubric}
    Holding the Paten under the Cup with his left hand, the Priest is to reverently receive the Blood with the particle.
\end{rubric}
\begin{rubric}
	Having received, if there be any to be communicated, let him communicate them before he purifieth himself.
\end{rubric}
%\end{multicols}

\subsection{Communion of the People}
\begin{rubric}
    If there be any to be communicated, the Priest shall genuflect and place the consecrated particles in a Ciborium (or if there are few to be communicated, on the Paten), unless from the beginning they had been placed in a Ciborium.
\end{rubric}
\begin{rubric}
    If the Priest is to administer Communion from the reserved Sacrament, he openeth the Tabernacle, genuflecteth, taketh out the Ciborium, \& placeth it upon the Corporal.
\end{rubric}
\begin{rubric}
    Then shall the Priest genuflect, take the Ciborium with his left hand (or the Cup, if he consecrated on the Paten), and take one particle with his right hand takes, which he holdeth between his thumb and forefinger slightly raised above the Ciborium (Cup); and turning to the people in the midst of the Altar, he saith in the clear voice:
\end{rubric}
\elcol{℣. Behold the Lamb of God. Behold him who taketh away the sins of the world.}{℣. Ecce Agnus Dei, ecce, qui tollit peccáta mundi.}
\elcol{℟. Lord, I am not worthy that thou shouldest come under my roof, but speak the word only, and my soul shall be healed.}{℟. Dómine, non sum dignus, ut intres sub tectum meum, sed tantum dic verbo, et sanábitur ánima mea.}
\elcol{℟. Lord, I am not worthy that thou shouldest come under my roof, but speak the word only, and my soul shall be healed.}{℟. Dómine, non sum dignus, ut intres sub tectum meum, sed tantum dic verbo, et sanábitur ánima mea.}
\elcol{℟. Lord, I am not worthy that thou shouldest come under my roof, but speak the word only, and my soul shall be healed.}{℟. Dómine, non sum dignus, ut intres sub tectum meum, sed tantum dic verbo, et sanábitur ánima mea.}
\begin{rubric}
%This rubric is adapted from the Antiochian Service Book (Eastern Rite).
	After the Priest returneth to the Altar, the Pre-Communion Hymn (p. \pageref{byzantine}) is to then be sung as communicants approach the Altar Rail.\par
	\textsc{Note,} The Pre-Communion Hymn may be said or sung any time after the Prayer of Humble Access.
\end{rubric}
\begin{rubric}\label{communionrubrics}
    Only Orthodox Christians in good standing and properly disposed may receive the Eucharist.\par
    \textsc{Note,} Each communicant shall receive the Communion kneeling.
\end{rubric}

%MANUAL ADJUSTMENT:
\clearpage

\begin{rubric}
    If they are to communicate, the Priest first communicateth the Sacred Ministers, and then the other Priests and Clerics in choir. (Priest and Deacons shall wear a stole either of white colour or of the same colour as the administering Priest.) And last of all, he proceedeth to communicate the others, beginning with those on the Epistle side.
\end{rubric}
\subsubsection{Distribution of the Body \& Blood Separately}
\begin{rubric}
    If both kinds are to be administered separately, the Priest, when giving the Body, shall make with the Host the sign of the Cross over the Ciborium and place the Host on the tongue of each communicant, while saying:
\end{rubric}

 \lett{T}{he} Body of our Lord Jesus Christ, which was given for thee, preserve thy body and soul unto everlasting life. Take and eat this in remembrance that Christ died for thee, and feed on him in thy heart by faith, with thanksgiving.
\begin{rubric}
    And then, while administering the Cup to each communicant, the Priest, keeping hold of the Cup, shall say:
\end{rubric}

\lett{T}{he} Blood of our Lord Jesus Christ, which was shed for thee, preserve thy body and soul unto everlasting life. Drink this in remembrance that Christ's Blood was shed for thee, and be thankful.
\subsubsection{Distribution of the Body \& Blood Together}
\begin{rubric}
    But if the Body and Blood of Christ are to be administered together by intinction, then the Priest, when giving the Sacrament to each one, intincteth the Host into the Cup; then, making with the Host the sign of the Cross over the Cup, he placeth the Host on the tongue of each communicant, while saying:
\end{rubric}

%Adapted the Proposed English 1928 BCP
\lett{T}{he} Body of our Lord Jesus Christ which was given for thee, and his Blood which was shed for thee. Take them in remembrance that Christ died for thee, and feed on him in thy heart by faith with thanksgiving.

\begin{rubric}
    The Communion ended, the Priest shall return to the Altar, saying nothing, nor doth he give the Blessing since he will give it at the end of Mass.    
\end{rubric}

%MANUAL ADJUSTMENT:
\clearpage
\subsection{Ablutions}
\begin{rubric}
    When the faithful have received Holy Communion, the Deacon and Priest return the Blessed Sacrament to the Altar. The Priest shall then consume everything not being reserved.\par
    \textsc{Note,} When prepared for reservation, the Host shall be touched with the Blood, the Blood then being completely consumed.
\end{rubric}

\begin{rubric}
    Then shall the Priest say,
\end{rubric}
\lett{G}{rant,} O Lord, that what we have taken with our mouths we may receive with a pure heart: and from a temporal gift may it become to us an everlasting remedy.

\begin{rubric}
    Meanwhile, he presenteth the Cup to the Minister, who shall pour in a little wine and water, wherewith he purifieth himself. Then he continueth,    
\end{rubric}
\lett{L}{et} thy Body, O Lord, which I have taken, and thy Blood which I have drunk, cleave to my members: and grant; that no stain of sin may remain in me, whom thou hast refreshed with pure and holy sacraments: Who livest and reignest, world without end. Amen.
\begin{rubric}
    Then shall the Priest wash and wipe his fingers and take the ablution. Then he wipeth his mouth and the Cup.
\end{rubric}
\begin{rubric}
	After folding the Corporal, he is to cover the Cup and place it upon the Altar as before. Then he proceedeth with the Mass.
\end{rubric}

%MANUAL ADJUSTMENT:
\vspace{-1ex}

\subsection{Thanksgiving}
\begin{rubric}
    Standing with his hands joined, the Priest shall read the Communion Antiphon. Then, at the Epistle corner, he alone saith the Thanksgiving, in an audible voice.
\end{rubric}
\letuspray
\lett{A}{lmighty} and everliving God, we most heartily thank thee, for that thou dost vouchsafe to feed us who have duly received these holy mysteries with the spiritual food of the most precious Body and Blood of thy Son our Saviour Jesus Christ; and dost assure us thereby of thy favour and goodness towards us; and that we are very members incorporate in the mystical body of thy Son, which is the blessed company of all faithful people; and are also heirs through hope of thy everlasting kingdom, by the merits of his most precious death and passion. And we humbly beseech thee, O heavenly Father, so to assist us with thy grace, that we may continue in that holy fellowship, and do all such good works as thou hast prepared for us to walk in; through Jesus Christ our Lord, to whom, with thee and the Holy Ghost, be all honour and glory, world without end.\par
℟. Amen.

%MANUAL ADJUSTMENT:
\clearpage
\subsection{Dismissal}

\begin{rubric}
    Then, with hands joined before his breast, the Priest shall proceed to the midst of the Altar, kiss the Altar, and turn to the People, with hands extended, saying:
\end{rubric}
\elcol{℣. The Lord be with you.}{℣. Dóminus vobíscum.}
\elcol{℟. And with thy spirit.}{℟. Et cum spíritu tuo.}
\begin{rubric}
    Then, turning back to the Book, he saith:
\end{rubric}
\elcol{℣. Let us pray.}{℣. Orémus.}

\subsubsection{Postcommunion Prayers}
\begin{rubric}
The Priest then saith the Postcommunion Prayers in the same manner, number, and order as the Collects at the beginning of Mass.
\end{rubric}

\subsubsection{Final Blessing}
\begin{rubric}
    After saying the last Prayer, the Priest is to return to the midst of the Altar, kissing it, and turning toward the people, saying:
\end{rubric}
\elcol{℣. The Lord be with you.}{℣. Dóminus vobíscum.}
\elcol{℟. And with thy spirit.}{℟. Et cum spíritu tuo.}
    \begin{rubric}
    	{On days when the \blackrubric{Gloria in exclesis} is said,}
    \end{rubric}
    \elcol{℣. Depart in peace.}{℣. Ite, Missa est.}
    \elcol{℟. Thanks be to God.}{℟. Deo grátias.}
    \begin{rubric}
    	{On days when the \blackrubric{Gloria in excelsis} is not said,}
    \end{rubric}
    \elcol{℣. Let us bless the Lord.}{℣. Benedicámus Dómino.}
    \elcol{℟. Thanks be to God.}{℟. Deo grátias.}
\begin{rubric}
In Requiem Masses,
\end{rubric}
    \elcol{℣. May they rest in peace.}{℣. Requiéscant in pace.}
    \elcol{℟. Amen.}{℟. Amen.}
\begin{rubric}
    In Eastertide, in Masses of the Season,
\end{rubric}
    \elcol{℣. Depart in peace, alleluia, alleluia.}{℣. Ite, Missa est, allelúja, allelúja.}
    \elcol{℟. Thanks be to God, alleluia, alleluia.}{℟. Deo grátias, allelúja, allelúja.}
\begin{rubric}
    Having said the Dismissal, the Priest boweth before the midst of the Altar, and with hands joined thereon, saith in the secret voice,% \blackrubric{Let this my bounden duty}.
\end{rubric}
\lett{L}{et} this my bounden duty and service be pleasing to thee, O holy Trinity: and grant; that the sacrifice, which I unworthy have offered before the eyes of thy majesty, may be acceptable to thee, and through thy mercy obtain thy gracious favour for me and all for whom I have offered it. Through Christ, our Lord. Amen.
\begin{rubric}
    Then---unless it be a Requiem Mass---the Congregation kneeling, the Priest shall kiss the Altar; raise his eyes; extend, raise, and join his hands; bow his head to the Cross; and give the Blessing, saying,
\end{rubric}
\lett{T}{he} Peace of God, which passeth all understanding, keep your hearts and minds in the knowledge and love of God, and of his Son Jesus Christ our Lord: And the Blessing of God Almighty,\margrub{1}{He turns to the people, blessing them once only, even in solemn Masses.} the Father, {\ding{64}} the Son, and the Holy Ghost, be amongst you, and remain with you always. \textit{Amen.}
\begin{rubric}
    In Pontifical Masses, the Blessing is given by the Bishop and is threefold, as ordered in the Pontifical.
\end{rubric}
\begin{rubric}
    In Masses of the Dead, the Blessing is not given, but having said \blackrubric{May they rest in peace} and \blackrubric{Let this my bounden duty,} he kisseth the Altar and readeth the Last Gospel.
\end{rubric}
\subsection{Last Gospel}
\begin{rubric}
    Unless a proper Last Gospel be provided, the Priest shall conclude the Mass with this Last Gospel, at the Gospel corner and with hands joined:
\end{rubric}
\elcol{℣. The Lord be with you.}{℣. Dóminus vobíscum.}
\elcol{℟. And with thy spirit.}{℟. Et cum spíritu tuo.}
\elcol{℣. The\margcolrub{2}{He signs with the sign of the Cross first the Altar or the Book, then himself on the forehead, mouth, \& breast.} {\ding{66}} beginning of the Holy Gospel according to John.}{℣. Inítium\marglatincolrub{2}{He signs with the sign of the Cross first the Altar or the Book, then himself on the forehead, mouth, \& breast.} {\ding{66}} sancti Evangélii secúndum Joánnem}
\elcol{℟. Glory be to thee, O Lord.}{℟. Glória tibi, Dómine.}
\elcol{In the beginning was the Word, and the Word was with God, and the Word was God. The same was in the beginning with God. All things were made by him; and without him was not any thing made that was made. In him was life; and the life was the light of men. And the light shineth in darkness; and the darkness comprehended it not. There was a man sent from God, whose name was John. The same came for a witness, to bear witness of the Light, that all men through him might believe. He was not that Light, but was sent to bear witness of that Light. That was the true Light, which lighteth every man that cometh into the world. He was in the world, and the world was made by him, and the world knew him not. He came unto his own, and his own received him not. But as many as received him, to them gave he power to become the sons of God, even to them that believe on his name: Which were born, not of blood, nor of the will of the flesh, nor of the will of man, but of God. \inrub{Everyone genuflects.} And the Word was made flesh, and dwelt among us, \inrub{Everyone rises.} (and we beheld his glory, the glory as of the only begotten of the Father,) full of grace and truth.\par
℟. Thanks be to God.}{In princípio erat Verbum, et Verbum erat apud Deum, et Deus erat Verbum. Hoc erat in princípio apud Deum. Omnia per ipsum facta sunt: et sine ipso factum est nihil, quod factum est: in ipso vita erat, et vita erat lux hóminum: et lux in ténebris lucet, et ténebræ eam non comprehendérunt.
Fuit homo missus a Deo, cui nomen erat Joánnes. Hic venit in testimónium, ut testimónium perhibéret de lúmine, ut omnes créderent per illum. Non erat ille lux, sed ut testimónium perhibéret de lúmine.
Erat lux vera, quæ illúminat omnem hóminem veniéntem in hunc mundum. In mundo erat, et mundus per ipsum factus est, et mundus eum non cognóvit. In própria venit, et sui eum non recepérunt. Quotquot autem recepérunt eum, dedit eis potestátem fílios Dei fíeri, his, qui credunt in nómine ejus: qui non ex sanguínibus, neque ex voluntáte carnis, neque ex voluntáte viri, sed ex Deo nati sunt. \inrub{Everyone genuflects.} Et Verbum caro factum est, \inrub{Everyone rises.} et habitávit in nobis: et vídimus glóriam ejus, glóriam quasi Unigéniti a Patre, plenum gráti{\ae} et veritátis.\par
℟. Deo grátias.}
\begin{rubric}
    At High Mass, a hymn may be sung while the Priest and servers leave the sanctuary.
\end{rubric}
\begin{rubric}
    As he departeth from the Altar, the Priest saith for thanksgiving the Antiphon \blackrubric{Let us sing}, with the rest (p. \pageref{CommunionThanksgiving}).
\end{rubric}
\subsection{General Rubrics}\label{GeneralRubrics}
%The Book of Common Prayer has always contained rubrics at the beginning and end of the liturgical text in order for the good order of the Church. Due to the unique situation of the Vicariate, it is fitting that these be accommodated to our use.
%\begin{rubric}
%Notwithstanding anything that is elsewhere enjoined in any Rubrick or Canon, the Priest, in celebrating the Holy Communion, shall wear the vestments proper to the Priest as of 1950 in dissident Rome. That is, Amice, Alb, Cincture, Maniple, Stole, and Chasuble (the latter three of the colour of the day).
%\end{rubric}
\begin{rubric}
    It is expedient for the Priest to have, set up on the Altar, cards with the most common prayers (or at least a convenient pamphlet).\par
\end{rubric}
%\begin{rubric}
%The Priest, the presence of an Ordinary notwithstanding, when he processes or reposes, wears a Canterbury Cap or Biretta.
%\end{rubric}
%\begin{rubric}
%There shall be no celebration of the Lord's Supper without at least one person, other than the Priest, physically present.
%\end{rubric}
\begin{rubric}
The Altar at the Communion-time---having a fair white linen cloth, a cross, and at least two candles upon it---shall stand in the Sanctuary.
\end{rubric}
\begin{rubric}
The Bread for the Eucharist ought to be provided by the Priest or the members of the Congregation. It must be leavened. That is, before and after being initially baked, it must be made of---and only made of---wheat flour, water, yeast, and salt. The Wine should be made of grapes, without any additives or supplements. The Cup for the Eucharist should be made of precious metal.
\end{rubric}
\begin{rubric}
It is an ancient and laudable rule of the Church to receive this Holy Sacrament fasting. That is, one who desireth to receive the Eucharist must consume neither food nor drink (water excepted) from the midnight before the liturgy until the time of reception. In the case of Masses which occur after noon, it is sufficient that the fasting begin at noon or immediately after one's noon-time meal.
\end{rubric}

%\begin{center}
%	{\textsc{Here Endeth the Liturgy of St. Tikhon.}}
%\end{center}
%\clearpage
%\subsection{Exhortations for Holy Communion}
%\begin{rubric}
%	At the time of the celebration of the Communion, the Priest may say this Exhortation. And Note, that the Exhortation shall be said on the First Sunday in Advent, the First Sunday in Lent, and Trinity Sunday.
%\end{rubric}
%\lett{D}{early} beloved in the Lord, ye who mind to come to the holy Communion of the Body and Blood of our Saviour Christ, must consider how Saint Paul exhorteth all persons diligently to try and examine themselves, before they presume to eat of that Bread, and drink of that Cup. For as the benefit is great, if with a true penitent heart and lively faith we receive that holy Sacrament; 
%Addition:
%that is, the very and true Body and Blood of Our Lord Jesus Christ; 
%
%so is the danger great, if we receive the same unworthily. Judge therefore yourselves, brethren, that ye be not judged of the Lord; repent you truly for your sins past; 
%Addition:
%confess them to your spiritual father; 
%
%have a lively and steadfast faith in Christ our Saviour; amend your lives, and be in perfect charity with all men; so shall ye be meet partakers of those holy mysteries. And above all things ye must give most humble and hearty thanks to God, the Father, the Son, and the Holy Ghost, for the redemption of the world by the death and passion of our Saviour Christ, both God and man; who did humble himself, even to the death upon the Cross, for us, miserable sinners, who lay in darkness and the shadow of death; that he might make us the children of God, and exalt us to everlasting life. And to the end that we should always remember the exceeding great love of our Master, and only Saviour, Jesus Christ, thus dying for us, and the innumerable benefits which by his precious blood-shedding he hath obtained for us; he hath instituted and ordained holy mysteries, as pledges of his love, and for a continual remembrance of his death, to our great and endless comfort. To him therefore, with the Father and the Holy Ghost, let us give (as we are most bounden) continual thanks; submitting ourselves wholly to his holy will and pleasure, and studying to serve him in true holiness and righteousness all the days of our life. Amen.
%\begin{rubric}
%	When the Minister giveth warning for the Celebration of the Holy Communion, (which he shall always do upon the Sunday, or some Holy-day, immediately preceding,) he shall read this Exhortation following, or so much thereof as, in his discretion, he may think convenient.
%\end{rubric}
%\lett{D}{early} beloved, on \textit{N.} day next I purpose, through God's assistance, to administer to all such as shall be religiously and devoutly disposed the most comfortable Sacrament of the Body and Blood of Christ; to be by them received in remembrance of his meritorious Cross and Pas-sion; whereby alone we obtain remission of our sins, and are made partakers of the Kingdom of heaven. Wherefore it is our duty to render most humble and hearty thanks to Almighty God, our heavenly Father, for that he hath given his Son our Saviour Jesus Christ, not only to die for us, but also to be our spiritual food and sustenance in that holy Sacrament. Which being so divine and comfortable a thing to them who receive it worthily, and so dangerous to those who will presume to receive it unworthily; my duty is to exhort you, in the mean season to consider the dignity of that holy mystery, and the great peril of the unworthy receiving thereof; and so to search and examine your own consciences, and that not lightly, and after the manner of dissemblers with God; but so that ye may come holy and clean to such a heavenly Feast, in the marriage-garment required by God in holy Scripture, and be received as worthy partakers of that holy Table.
%
%    The way and means thereto is: First, to examine your lives and conversations by the rule of God's commandments; and whereinsoever ye shall perceive yourselves to have offended, either by will, word, or deed, there to bewail your own sinfulness, and to confess yourselves to Almighty God, with full purpose of amendment of life. And if ye shall perceive your offences to be such as are not only against God, but also against your neighbours; then ye shall reconcile yourselves unto them; being ready to make restitution and satisfaction, according to the uttermost of your powers, for all injuries and wrongs done by you to any other; and being likewise ready to forgive others who have offended you, as ye would have forgiveness of your offences at God's hand: for otherwise the receiving of the holy Communion doth nothing else but increase your condemnation. Therefore, if any of you be a blasphemer of God, an hinderer or slanderer of his Word, an adulterer, or be in malice, or envy, or in any other grievous crime; repent you of your sins, or else come not to that holy Table.
%    
%    And because it is requisite that no man should come to the holy Communion, but with a full trust in God's mercy, and with a quiet conscience; 
    %This phrase is omitted so to indicate that all such sinners ought to confess:
    %therefore, if there be any of you, who by this means cannot quiet his own conscience herein, but requireth further comfort or counsel, 
%    let him come to me, or to some other Priest of 
    %Change:
%    Christ's Church, confessing his sins and opening his grief; that he may receive such godly counsel, advice, and absolution of his sins as may tend to the quieting of his conscience, the removing of all scruple and doubtfulness, and the sure benefit of Jesus Christ's promise: the remission of sins and bestowing of the benefits of his Passion.
%\begin{rubric}
%	Or, in case he shall see the People negligent to come to the Holy Communion, instead of the former, he may use this Exhortation.
%\end{rubric}
%\lett{D}{early} beloved brethren, on \textit{N.} I intend, by God's grace, to celebrate the 
%Change:
%Holy Sacrifice of the Mass: 
%
%unto which, in God's behalf, I bid you all who are here present; and beseech you, for the Lord Jesus Christ's sake, that ye will not refuse to come thereto, being so lovingly called and bidden by God himself. Ye know how grievous and unkind a thing it is, when a man hath prepared a rich feast, decked his table with all kind of provision, so that there lacketh nothing but the guests to sit down; and yet they who are called, without any cause, most unthankfully refuse to come. Which of you in such a case would not be moved? Who would not think a great injury and wrong done unto him? Wherefore, most dearly beloved in Christ, take ye good heed, lest ye, withdrawing yourselves from this holy Supper, provoke God's indignation against you. It is an easy matter for a man to say, I will not communicate, because I am otherwise hindered with worldly business. But such excuses are not so easily accepted and allowed before God. If any man say, I am a grievous sinner, and therefore am afraid to come: wherefore then do ye not repent and amend? 
%Addition:
%Wherefore do ye not approach myself or one of Christ's Priests for Confession and Absolution? 
%
%When God calleth you, are ye not ashamed to say ye will not come? When ye should return to God, will ye excuse yourselves, and say ye are not ready? Consider earnestly with yourselves how little such feigned excuses will avail before God. Those who refused the feast in the Gospel, because they had bought a farm, or would try their yokes of oxen, or because they were married, were not so excused, but counted unworthy of the heavenly feast. Wherefore, according to mine office, I bid you in the Name of God, I call you in Christ's behalf, I exhort you, as ye love your own salvation, that ye will be partakers of this holy Communion. And as the Son of God did vouchsafe to yield up his soul by death upon the Cross for your salvation; so it is your duty to receive the Communion in remembrance of the sacrifice of his death, as he himself hath commanded: which if ye shall neglect to do, consider with yourselves how great is your ingratitude to God, and how sore punishment hangeth over your heads for the same; when ye wilfully abstain from the Lord's Table, and separate from your brethren, who come to feed on the banquet of that most heavenly food. These things if ye earnestly consider, ye will by God’s grace return to a better mind: for the obtaining whereof we shall not cease to make our humble petitions unto Almighty God, our heavenly Father.
\include{pages/HolyMass/PreCommunionHymn}
\fancyhead[RO,LE]{\textit{Prefaces}}
\fancyhead[RE,LO]{}
\section{Prefaces for the Mass}\label{prefaces}

\subsection{Proper Prefaces without Notation}

\begin{multicols}{2}

\raggedcolumns	

\centerline{\textit{Advent}}%From the Ritual of the Altar
\begin{rubric}
	Every day in \blackrubric{Advent,} including \blackrubric{Christmas Eve.}
\end{rubric}
\lett{T}{hrough} Jesus Christ our Lord; whom thou, the merciful and faithful God, didst promise as a Saviour to lost mankind; by whose truth the unlearned should be instructed, by whose holiness the ungodly should be justified, by whose power the weak should be holpen. Now, therefore, when the time in which thou shouldest send him draweth near, and now that the day of our deliverance already shineth; trusting in these thy promises, with holy joy we do rejoice.\\

\centerline{\textit{Christmas}}
%\subsubsection{Christmas}%1928 BCP:
\begin{rubric}
	Upon \blackrubric{Christmas Day} and its Octave, until \blackrubric{Epiphany Eve}, inclusive, and on the Feast of the \blackrubric{Most Holy Body of Christ} and its Octave.
\end{rubric}
\lett{B}{ecause} thou didst give Jesus Christ, thine only Son, to be born (as at this time) for us; who, by the operation of the Holy Ghost, was made very man, of the substance of the Virgin Mary his mother; and that without spot of sin, to make us clean from all sin.%\\

%MANUAL ADJUSTMENT:
\newcolumn

\centerline{\textit{Epiphany}}
%\subsubsection{Epiphany}%1928 BCP:
\begin{rubric}
	Upon the \blackrubric{Epiphany} and its Octave.
\end{rubric}
\lett{T}{hrough} Jesus Christ our Lord, who, in substance of our mortal flesh, manifested forth his glory; that he might bring us out of darkness into his own glorious light.\\

\centerline{\textit{Septuagesima Sunday}}
%\subsubsection{Septuagesima Sunday}
\begin{rubric}
	Optionally upon \blackrubric{Septuagesima Sunday.}
\end{rubric}
\lett{F}{or} by things seen, we are instructed how we should pursue those things unseen. And so are reminded, as the year leadeth on its course, to pass, from things past unto things to come, from the old man into the newness of life; that, being freed from reliance on earthly goods, we may lay hold upon the excelling richness of the heavenly gift; and by that food, which is preserved for us in every age, we may arrive unto that victorious sustenance which remaineth for ever, even Jesus Christ our Lord.\\

%Quia per ea quæ conspiciuntur instruimur quibus modis ad invisibilia tendere debeamus. Denique commonemur, anni ducente successu, de præteritis ad futura, de vetustate in novitatem vitæ transire, Ut terrenis sustentationibus expediti, cælestis doni capiamus desiderabilius ubertatem, Et per eum cibum qui beneficiis prærogatur alternis, perveniamus ad victum sine fine mansurum, Jesum Christum Dominum nostrum.

\centerline{\textit{Sexagesima Sunday}}
%\subsubsection{Sexagesima Sunday}
\begin{rubric}
	Optionally upon \blackrubric{Sexagesima Sunday.}
\end{rubric}
\lett{W}{ho} deigned to instruct the rational creature by such stewardship, lest he be given over to temporal things and strive not for the everlasting reward. That he may neither withdraw from correction nor become haughty in approbation, but rather overcome all adversities, this being his glorious devotion: Jesus Christ our Lord.\\

%Qui rationabilem creaturam, ne temporalibus dedita bonis, ad praemia sempiterna non tendat, ea dispensatione dignaris erudire, ut nec castigatione deficiat, nec prosperitatibus insolescat; sed hoc potius fiat ejus gloriosa devotio, quo nullis adversitatibus obruta superetur. Per Christum Dominum nostrum.

\centerline{\textit{Quinquagesima Sunday}}
%\subsubsection{Quinquagesima Sunday}
\begin{rubric}
	Optionally on \blackrubric{Quinquagesima Sunday.}
\end{rubric}
\lett{A}{nd} to entreat thy Majesty by lowly devotion, that thou, beholding the weakness of our mortal nature, accuse us not of our depravity in thy wrath, but rather cleanse, cultivate, and console us by thine immeasurable gentleness. And since we can do no good thing without thee, bestow upon us thy grace alone that we may live in wholesome conversation.\\

%Et majestatem tuam cernua devotione exorare, ut modulum terrenae fragilitatis aspiciens, non in ira tua pro nostra pravitate nos arguas, sed immensa clementia purifices, erudias, consoleris. Qui quum sine te nihil possumus facere, quod tibi sit placitum, tua nobis gratia sola praestabit, ut salubri conversatione vivamus.

\vspace{-1ex}

\centerline{\textit{Lent}}
%\subsubsection{Lent}%1933 English Missal:
\begin{rubric}
	From \blackrubric{Ash Wednesday} until the Saturday before \blackrubric{Passion Sunday,} inclusive.
\end{rubric}
\lett{W}{ho} by bodily fasting dost overcome vice, dost raise the mind, and dost bestow on us virtue and heavenly rewards: through Christ, our Lord.\\

\vspace{-1ex}

\centerline{\textit{Holy Rood}}
%\subsubsection{Holy Rood}%1933 English Missal:
\begin{rubric}
	From \blackrubric{Passion Sunday} until \blackrubric{Maundy Thursday,} inclusive. And all Masses of the \blackrubric{Holy Rood,} the \blackrubric{Passion of the Lord,} or the \blackrubric{Most Precious Blood,} and within their Octaves, even when celebrated within the Octave of \blackrubric{Christmas.}
\end{rubric}
\lett{W}{ho} by the tree of the Cross didst give salvation unto mankind: that whence death arose, thence life might rise again: and that he, who by a tree overcame, might also by a tree be overcome: through Christ our Lord.

%MANUAL ADJUSTMENT:
\vspace{0.5\baselineskip}

\centerline{\textit{Easter}}
%\subsubsection{Easter} %1928 BCP:
\begin{rubric}
	From \blackrubric{Holy Saturday} until \blackrubric{Ascension Eve,} inclusive.
\end{rubric}
\lett{B}{ut} chiefly are we bound to praise thee for the glorious Resurrection of thy Son Jesus Christ our Lord: for he is the very Paschal Lamb, which was offered for us, and hath taken away the sin of the world; who by his death hath destroyed death, and by his rising to life again hath restored to us everlasting life.

%MANUAL ADJUSTMENT:
\vspace{0.5\baselineskip}

\centerline{\textit{Ascension}}
%\subsubsection{Ascension}%1928 BCP:
\begin{rubric}
	From \blackrubric{Ascension Thursday} until the Friday after its Octave, inclusive. It is also said, even without a Commemoration, on the Friday after the Octave.
\end{rubric}
\lett{T}{hrough} thy most dearly beloved Son Jesus Christ our Lord; who, after his most glorious Resurrection, manifestly appeared to all his Apostles, and in their sight ascended up into heaven, to prepare a place for us; that where he is, thither we might also ascend, and reign with him in glory.\\

\centerline{\textit{Holy Ghost}}
%\subsubsection{Holy Ghost}%1928 BCP Whitsun Preface:
\begin{rubric}
	From \blackrubric{Whitsunday Eve} until \blackrubric{Trinity Sunday,} exclusive, in Votive Masses of the Holy Ghost, even when celebrated within the Octave of \blackrubric{Christmas.}
\end{rubric}
\lett{T}{hrough} Jesus Christ our Lord; according to whose most true promise, the Holy Ghost came down (as at this time) from heaven, lighting upon the Apostles, to teach them, and to lead them into all truth; giving them boldness with fervent zeal constantly to preach the Gospel unto all nations; whereby we have been brought out of darkness and error into the clear light and true knowledge of thee, and of thy Son Jesus Christ.\\

\centerline{\textit{Most Holy Trinity}}
%\subsubsection{Most Holy Trinity}%1928 BCP:
\begin{rubric}
	Upon Sundays, from the \nth{3} Sunday after \blackrubric{Trinity} until the \blackrubric{Sunday Next before Advent,} inclusive. And also in Masses of the \blackrubric{Most Holy Trinity,} even when celebrated within the Octave of \blackrubric{Christmas.} And on Sundays in \blackrubric{Septuagesimatide.}
\end{rubric}
\lett{W}{ho,} with thine only-begotten Son, and the Holy Ghost, art one God, one Lord, in Trinity of Persons and in Unity of Substance. For that which we believe of thy glory, O Father, the same we believe of the Son, and of the Holy Ghost, without any difference of inequality.\\

\centerline{\textit{Compassion of Our Lord}}
%\subsubsection{Compassion of Our Lord}%1933 English Missal:
\begin{rubric}
	Upon the Feast of the of the \blackrubric{Compassion of Our Lord} and its Octave Day.
\end{rubric}
\lett{W}{ho} didst will that thine only-begotten Son, as he hung upon the Cross, should be pierced by the soldier's spear: that his Heart thus opened, the shrine of the divine bounty, might pour forth upon us streams of mercy and of grace: that as it never ceased to burn with love for us, so it might be a rest for the godly, and an open refuge of salvation for the penitent.\\

\centerline{\textit{Christ the King}}
%\subsubsection{Christ the King}%1933 English Missal:
\begin{rubric}
	Masses of \blackrubric{Our Lord Jesus Christ the King,} even when they are celebrated within the Octave of \blackrubric{Christmas.}
\end{rubric}
\lett{W}{ho} didst anoint thine only-begotten Son, our Lord Jesus Christ, with the oil of gladness, to be a Priest for ever and the King of all the world: that, offering himself an unspotted sacrifice of peace upon the altar of the cross, he might accomplish the sacrament of the redemption of mankind: and making all creatures subject to his governance, might deliver up to thine infinite Majesty an eternal and universal kingdom. A kingdom of truth and life: a kingdom of sanctity and grace: a kingdom of justice, love, and peace.\\

\centerline{\textit{All Hallows'}}
%\subsubsection{All Hallows'}%1928 BCP:
\begin{rubric}
	Upon the Feast of \blackrubric{All Hallows} and its Octave.
\end{rubric}
\lett{W}{ho,} in the multitude of thy saints, hast compassed us about with so great a cloud of witnesses that we, rejoicing in their fellowship, may run with patience the race that is set before us, and together with them may receive the crown of glory that fadeth not away.\\

%MANUAL ADJUSTMENT:
\newcolumn

\centerline{\textit{Of Our Lord}}
%\subsubsection{Of Our Lord}
\begin{rubric}
	Upon the Feasts of the \blackrubric{Purification,} \blackrubric{Annunciation,} and \blackrubric{Transfiguration.}
\end{rubric}
\lett{B}{ecause} in the Mystery of the Word made flesh, thou hast caused a new light to shine in our hearts, to give the knowledge of thy glory in the face of thy Son, Jesus Christ our Lord.\\

\centerline{\textit{Blessed Virgin Mary}}
%\subsubsection{Blessed Virgin Mary}
\begin{rubric}
	In Masses of the \blackrubric{Blessed Virgin Mary,} and their Octaves.
\end{rubric}
\lett{A}{nd} that in the \textit{N.} of Blessed Mary ever Virgin we should praise, bless, and magnify thee. In that by the overshadowing of the Holy Ghost she conceived thine only-begotten Son: and, the glory of her maidenhood yet abiding, shed forth upon the world the light eternal, Jesus Christ, our Lord.\\

\centerline{\textit{St. Joseph, Spouse of the Blessed Virgin}}
%\subsubsection{St. Joseph, Spouse of the Blessed Virgin}
\begin{rubric}
	In Masses of \blackrubric{St. Joseph,} and their Octaves.
\end{rubric}
\lett{A}{nd} that in the Festivity (Veneration) of blessed Joseph, we should magnify, bless, and glorify thee with worthy praises. Who being a just man was given by thee for a Spouse to the Virgin Mother of God: and, being a faithful and wise servant, was made ruler over thy Household: that in the office of a father he should guard thine only-begotten Son, conceived by the overshadowing of the Holy Ghost, even Jesus Christ our Lord.\\

\centerline{\textit{Apostles \& Evangelists}}
%\subsubsection{Apostles \& Evangelists}
\begin{rubric}
	In Masses of the \blackrubric{Apostles} or \blackrubric{Evangelists} and in a Mass of the Election and Enthronement of the Patriarch and their anniversaries.
\end{rubric}
\lett{T}{hat} we, O Lord, should humbly entreat thee, that thou, the everlasting Shepherd, wouldst not forsake thy flock: but through thy blessed Apostles keep it by thy continual protection. That it may be governed by those same rulers, whom in thy stead thou hast appointed for thy work, as shepherds of thy people.\\

\centerline{\textit{Requiem}}
%\subsubsection{Requiem}
\begin{rubric}
	In all \blackrubric{Masses for the Dead.}
\end{rubric}
\lett{T}{hrough} Christ, our Lord. In whom hath shone forth unto us the hope of a blessed resurrection, that they who bewail the certain condition of their mortality may be consoled by the promise of immortality to come. For to thy faithful people, O Lord, life is changed not taken away: and at the dissolution of the tabernacle of this earthly sojourning a dwelling-place eternal is made ready in the heavens.

\end{multicols}

%Propers
\cleardoublepage
\fancyhead[RE,LO]{}\fancyhead[RO,LE]{}
\fancyhead[C]{}\thispagestyle{empty}
\phantomsection
\addcontentsline{toc}{chapter}{Propers of the Church Year}

\begin{figure}[H]
    \centering
    \includegraphics[width=98mm, height=140mm]{Propers/Vision.eps}
    \caption{\textsc{\Huge{Propers of the\\
   	Church Year}}}
\end{figure}
\include{pages/Propers/PropSeas}
\include{pages/Propers/PropTime}
   %\section{Table of Prayers}
   \fancyhead[C]{{\LARGE Table of Prayers}}
   \fancyhead[RO,LE]{}
   \fancyhead[RE,LO]{}
   \phantomsection
	\addcontentsline{toc}{section}{Table of Prayers}
\begin{secrubric}
	These are provided for the additional Prayers during Mass.
\end{secrubric}

   \begin{multicols}{2}
   \subsection{Of St. Mary in Advent}\label{SPMaryInAdvent}
   \collect
   \lett{O}{God,} who wast pleased that thy Word should take flesh of the womb of the blessed Virgin Mary at the message of an Angel: grant to thy humble servants; that we, who believe her to be truly the Mother of God, may be aided by her intercession in thy sight. Through the same.
   \secret
   \lett{W}{e} beseech thee, O Lord, stablish in our hearts the mysteries of thy true religion: that, as we believe thy Son, conceived of a Virgin, to be very God and very man; so by the power of his life-giving resurrection we may be found worthy to attain unto everlasting felicity. Through the same.
   \postcommunion
   \lett{W}{e} beseech thee, O Lord, pour thy grace into our hearts: that, as we have known the incarnation of thy Son Christ by the message of an Angel; so by his cross and passion we may be brought unto the glory of his resurrection. (Through the same.)
   
   %MANUAL ADJUSTMENT:
   \newcolumn

    \subsection{Of St. Mary after Christmas}\label{SPMaryPostChristmas}
    \collect
    \lett{O}{God,} who by the virgin child-bearing of blessed Mary hast bestowed upon mankind the rewards of eternal salvation: grant, we beseech thee; that we may perceive her intercession for us, through whom we have been counted worthy to receive the author of life, Jesus Christ thy Son, our Lord. (Who liveth.)
    \secret
    \lett{T}{hrough} thy mercy, O Lord, and the intercession of blessed Mary ever Virgin, may this oblation avail for our prosperity and peace both now and for ever. (Through.)
    \postcommunion
    \lett{M}{ay} this communion, O Lord, cleanse us from guilt: and at the intercession of blessed Mary the Virgin Mother of God, make us partakers of thy heavenly healing. (Through the same.)
    
    %MANUAL ADJUSTMENT:
    \newcolumn
    
    \subsection{Of St. Mary in Eastertide}\label{SPMaryInEaster}
    \collect
    \lett{G}{rant,} we beseech thee, O Lord God, that we thy servants may enjoy perpetual health of mind and of body: and, at the glorious intercession of blessed Mary ever Virgin, be delivered from present sadness, and rejoice in everlasting gladness. (Through.)
    \secret
    \lett{T}{hrough} thy mercy, O Lord, and the intercession of blessed Mary ever Virgin, may this oblation avail for our prosperity and peace, both now and ever. (Through.)
    \postcommunion
    \lett{G}{rant,} we beseech thee, O Lord: that we who have received these aids to our salvation may at all times and in all places be protected by the advocacy of blessed Mary ever Virgin; in whose honour we have made these offerings to thy majesty. (Through.) %MANUAL ADJUSTMENT:\\
    
    \newcolumn

    \subsection{Of Saints}\label{SPSaints}
    \collect
    \lett{D}{efend} us, O Lord, we beseech thee, from all dangers of mind and of body: and at the intercession of the blessed and glorious Mother of God, Mary ever Virgin, with blessed Joseph, thy blessed Apostles Peter and Paul, blessed \emph{N.}, and all the Saints, graciously grant us health and peace; that all adversities and errors being done away, thy Church may serve thee in freedom and quietness. (Through the same.)
    \secret
    \lett{G}{raciously} hear us, O God our Saviour: and by the power of this sacrament defend us from all enemies of mind and of body; granting us grace in this world, and glory in that which is to come. (Through.)
    \postcommunion
    \lett{W}{e} beseech thee, O Lord, that the gift now offered in this divine sacrament may cleanse and defend us: that at the intercession of blessed Mary the Virgin Mother of God, of blessed Joseph, of thy blessed Apostles Peter and Paul, of blessed \emph{N.} and of all the Saints; it may set us free from every perverse way, and deliver us from all adversities. (Through the same.)
    
   \subsection{Against the Persecutors of the Church}\label{SPAgainst}
   \collect
   \lett{W}{e} beseech thee, O Lord, mercifully to hear the prayers of thy Church: that all adversities and errors being done away, she may serve thee in freedom and quietness. Through.
   \secret
   \lett{D}{efend} us, O Lord, who wait upon thy mysteries: that we, cleaving fast to things heavenly, may serve thee both in body and soul. Through.
   \postcommunion
   \lett{W}{e} beseech thee, O Lord our God: that we, whom thou makest to rejoice in the partaking of heavenly things, may by thee be defended against all earthly perils. Through.
   
   %MANUAL ADJUSTMENT:
   \newcolumn

   \subsection{For the Chief Bishop}\label{SPChiefBishop}
   \collect
   \lett{O}{God,} the pastor and ruler of all the faithful, mercifully look upon thy servant \emph{N.}, whom thou hast been pleased to set as pastor over thy Church; grant him, we beseech thee, to be in word and conversation a wholesome example to the people committed to his charge; that he with them may attain unto everlasting life. Through.
   \secret
   \lett{L}{ook} favourably, we beseech thee, O Lord, upon the gifts which we offer: and guide with thy continual protection thy servant \textit{N.}, whom thou hast been pleased to set as pastor over thy Church. Through.
   \postcommunion
   \lett{W}{e} beseech thee, O Lord, that the divine sacrament which we have here received may be our protection: and together with the flock entrusted to his charge ever save and defend thy servant \textit{N.}; whom thou hast been pleased to set as pastor over thy Church. Through.
   
   %MANUAL ADJUSTMENT:
   \newcolumn

   \subsection{Of the Holy Ghost}\label{SPHolyGhost}
   \collect
   \lett{G}{od,} who didst teach the hearts of thy faithful people, by the sending to them the light of thy Holy Spirit: grant us by the same Spirit to have a right judgment in all things; and evermore to rejoice in his holy comfort. (Through . . . in the unity of the same Holy Spirit.)
   \secret
   \lett{S}{anctify,} we beseech thee, O Lord, the gifts which we offer: and cleanse our hearts by the enlightening of the Holy Spirit. (Through . . . in the unity of the same Holy Spirit.)
   \postcommunion
   \lett{P}{our} thy Holy Spirit upon us, O Lord, and cleanse our hearts: that they may be made fruitful by the inward sprinkling of his dew. (Through . . . in the unity of the same Holy Spirit.)

    \subsection{For the Living and Departed}\label{SPLivingDeparted}
   \collect
   \lett{A}{lmighty} and everlasting God, who hast dominion both of the living and of the dead, and hast mercy upon all who thou foreknowest will be thine in faith and works: we humbly beseech thee; that all those for whom we are minded to pour forth our prayers, whether in this present world they still be held in the flesh, or being delivered from the body have passed into that which is to come, may at the intercession of all thy Saints obtain of thy bountiful goodness the remission of all their sins. Through.
   \secret
   \lett{O}{God,} to whom alone is known the number of the elect, whom thou hast appointed unto heavenly felicity: grant, we beseech thee: that at the intercession of all thy Saints, the names of all those who have been commended to our prayers, and of all thy faithful people, may ever remain written in the book of those that are predestined to everlasting blessedness. Through.
   \postcommunion
   \lett{A}{lmighty} and merciful God, cleanse us, we pray thee, by the sacraments which we have received: and, at the intercession of all thy Saints, grant; that this thy sacrament may not bring upon us guilt to our punishment but pardon to our salvation: that it may cleanse away sins, that it may confirm the feeble, that it may defend us against every danger in this world: that it may obtain for thy faithful people both living and dead the remission of all their sins. Through.
 \end{multicols}

%MANUAL ADJUSTMENT:
\clearpage
\begin{rubric}
    The following propers may also be used when \blackrubric{ad libitum} prayers are permitted.
\end{rubric}
\vspace{-1ex}
\begin{multicols}{2}
\subsection{To Ask the Prayers of the Saints}
\collect
\lett{G}{rant,} we beseech thee, almighty God: that the intercession of holy Mary, Mother of God, and of all the holy Apostles, Martyrs, Confessors and Virgins, and of all thy elect, may everywhere cause us to rejoice; that while we call to mind their merits, we may perceive their advocacy. Through the same.
\secret
\lett{A}{ccept,} O Lord, the gifts we offer: and, at the intercession of blessed Mary, ever Virgin, with all thy Saints, defend us from all perils. Through.
\postcommunion
\lett{G}{rant,} we beseech thee, O Lord: that we who have received these heavenly sacraments, in veneration of the memory of blessed Mary, ever Virgin, and of all thy Saints; may attain in everlasting joys the fulfilment of our service in this life. Through.

%\begin{inhead}
%    or
%\end{inhead}
%
%\collect
%\lett{D}{efend} us, O Lord, we beseech thee, from all dangers of mind and of body: and at the intercession of the blessed and glorious Mother of God, Mary, ever Virgin, with blessed Joseph, thy blessed Apostles, Peter and Paul, blessed \textit{N.}, and all the Saints, graciously grant us health and peace; that, all adversities and errors being done away, thy Church may serve thee in freedom and quietness. Through the same.
%\secret
%\lett{G}{raciously} hear us, O God our Saviour: and by the power of this sacrament defend us from all enemies of mind and of body; granting us grace in this world, and glory in that which is to come. Through.
%\postcommunion
%\lett{W}{e} beseech thee, O Lord, that the gift now offered in this divine sacrament may cleanse and defend us: that at the intercession of blessed Mary the Virgin Mother of God, of blessed Joseph, of thy blessed Apostles Peter and Paul, of blessed \textit{N.}, and of all the Saints; it may set us free from every perverse way and deliver us from all adversities. Through the same.

   %MANUAL ADJUSTMENT:
   \newcolumn

\subsection{For all Estates of Men in the Church}
\collect
\lett{A}{lmighty} and everlasting God, by whose Spirit the whole body of the Church is governed and sanctified: receive our supplications and prayers, which we offer before thee for all estates of men thy holy Church; that every member of the same, in his vocation and ministry, may truly and godly serve thee. Through . . . in the unity of the same Holy Spirit.
\secret
\lett{G}{rant} unto thy servants, O Lord, the pardon of their sins, comfort in their life, and thy continual governance: so that they, serving thee aright, may continually obtain thy mercy. Through.
%MANUAL ADJUSTMENT:
\postcommunion
\lett{D}{eliver,} O Lord, we beseech thee, thy servants who call upon thee from all sin and from every foe: that they, walking in godly conversation, may be hurt by no adversity. Through.

%For the Chief Bishop

   %MANUAL ADJUSTMENT:
   \newcolumn

\subsection{For the Emperor}
\collect
\lett{O}{God,} the protector of all kingdoms, and specially of the Christian empire: grant unto thy servant, our emperor, \textit{N.}, to labour with wisdom for the triumph of virtue, that he who by thy institution is a prince, may ever by thy gift be powerful. Through.
\begin{rubric}
    \textsc{Note,} If he be not crowned, say \blackrubric{Emperor-elect}.
\end{rubric}

\secret
\lett{R}{eceive,} O Lord, the prayers and offerings of thy Church who calleth upon thee for the safety of thy servant: and for the protection of the faithful peoples perform with thy mighty arm thy wondrous works of old, that they may overcome the enemies of peace, and serve thee in freedom and safety. Through.
\postcommunion
\lett{O}{God,} who didst prepare the Roman empire for the preaching of the Gospel of the eternal King: defend thy servant, our emperor, \textit{N.}, with thy heavenly armour, that the peace of the churches may be troubled by no war nor tumult. Through.

   %MANUAL ADJUSTMENT:
   \newcolumn

\subsection{For the King}
\collect
\lett{W}{e} beseech thee, almighty God: that thy servant \textit{N.}, our King, who by thy mercy hath received the government of this realm, may also obtain an increase of every virtue; that, being meetly adorned therewith, he may flee from sin and iniquity, and by thy grace may attain unto thee who art the way, the truth, and the life. Through.
\secret
\lett{S}{anctify,} we beseech thee, O Lord, the gifts which we offer: that they may be made unto us the Body and Blood of thine only-begotten Son; and of thy great goodness grant that they may at all times be profitable unto our King, both for the obtaining of health of body and soul, and for the fulfilling of the charge committed unto him. Through.
\postcommunion
\lett{M}{ay} this saving oblation, O Lord, keep thy servant, \textit{N.}, our King, from all adversity: that he may preserve peace and tranquillity for thy Church; and after the course of this life may attain unto his eternal inheritance. Through.

   %MANUAL ADJUSTMENT:
   \newcolumn

\subsection{For Prelates \& Congregations Committed to their Charge}
\collect
\lett{A}{lmighty} and everlasting God, who alone workest great marvels: send down upon our Bishops, and all Congregations committed to their charge, the healthful spirit of thy grace; and that they may truly please thee, pour upon them the continual dew of thy blessing. Through.
\secret
\lett{L}{ook} down, O Lord, in mercy on the sacrifices of thy servants, which we devoutly celebrate for them unto the honour of thy name: that they may perceive the same to be profitable for the healing of their souls. Through.
\postcommunion
\lett{A}{ssist,} O Lord, with thy continual protection those whom thou dost refresh with this heavenly gift: and grant that, as thou never failest to comfort them, so thou wouldest make them worthy of everlasting redemption. Through.

   %MANUAL ADJUSTMENT:
   \newcolumn

\subsection{For the Congregation and Family}
\collect
\lett{W}{e} beseech thee, O Lord, at the intercession of blessed Mary, ever Virgin, to defend this thy family from all adversities: that they, devoutly serving thee with all their hearts, may by thy mercy be protected from the assaults of their enemies. Through.
\secret
\lett{A}{lmighty} God, we beseech thee to accept this offering of our devotion: that by the power of this sacrament we thy servants may be protected against all adversities. Through.
\postcommunion
\lett{G}{rant,} we beseech thee, O merciful God: that we who have here received the gifts of our redemption; may through the celebration of the same obtain the help of thy defence against all adversities. Through.


   %MANUAL ADJUSTMENT:
   \newcolumn

\subsection{For the Preservation of Concord in the Congregation}
\collect
\lett{O}{God,} who art the author of peace and lover of charity: grant unto thy servants that true unity which is according to thy will; that we may be delivered from all the temptations which beset us. Through.
\secret
\lett{W}{e} beseech thee, O Lord, mercifully to accept these sacrifices: and grant that we who pray to be absolved from our own offences may not be burdened with the sins of other men. Through.
\postcommunion
\lett{P}{our} forth upon us, O Lord, the Spirit of thy charity: that as thou hast fulfilled us with one heavenly bread, so of thy goodness thou wouldest make us to be of one heart and mind. Through . . . in the unity of the same.

   %MANUAL ADJUSTMENT:
   \newcolumn

%Against the Persecutors of the Church

\subsection{Against Persecutors \& Malefactors}
\collect
\lett{W}{e} beseech thee, O Lord, to bring down the pride of our enemies: and by the right hand of thy power destroy their obstinacy. Through.
\secret
\lett{G}{rant,} O Lord, that by the power of this mystery, we may be cleansed from our secret faults, and delivered from the snares of our enemies. Through.
\postcommunion
\lett{B}{ehold,} O God our defender, and protect us against the dangers of our enemies: that we may be delivered from all distress, and serve thee in freedom of spirit. Through.

\subsection{For any Necessity}
\collect
\lett{O}{God,} our refuge and strength, who art the author of all godliness: be ready, we beseech thee, to hear the devout prayers of thy Church; and grant that those things which we ask faithfully we may obtain effectually. Through.

   %MANUAL ADJUSTMENT:
   \newcolumn

\secret
\lett{G}{rant,} O merciful God: that this saving oblation may continually deliver us from our sins, and shield us from all adversities. Through.
\postcommunion
\lett{W}{e} who have partaken, O Lord, of the gifts of this sacred mystery, humbly beseech thee: that those things which thou hast commanded us to do in remembrance of thee may be profitable for the assistance of our infirmity. Who livest.

\subsection{In any Calamity}
\collect
\lett{A}{lmighty} God, despise not thy people who cry to thee in their affliction: but for the glory of thy name mercifully assist them in their tribulation. Through.
\secret
\lett{M}{ercifully} accept, O Lord, these sacrifices which thou hast ordained to be the propitiation of our sins, and the means whereby of thy loving mercy we are restored to salvation. Through.
\postcommunion
\lett{W}{e} beseech thee, O Lord, mercifully look upon our tribulation: and turn from us the fury of thine anger, which we most righteously deserve. Through.

\subsection{In time of Famine}
\collect
\lett{G}{rant} to us, we beseech thee, O Lord, the fulfilment of our devout supplications: and graciously turn away this famine; that the hearts of men may know that in thine anger thou sendest forth such scourges, and in thy mercy makest them to cease. Through.

\secret
\lett{O}{God,} who dost sustain the twofold nature of man by the food of the gifts we offer, and by this sacrament dost renew the same: grant, we beseech thee; that these thy gifts may never be lacking for our assistance both in body and soul. Through.

\postcommunion
\lett{O}{Lord,} support us, we beseech thee, with temporal food: even as thou dost vouchsafe to dwell in us in everlasting mysteries. Through.

\subsection{In time of Earthquake}
\collect
\lett{A}{lmighty} and everlasting God, who lookest upon the earth, and makest it to tremble: spare them that fear thee, be merciful to them that call upon thee; that whereas we are sore afraid for thy wrath that shaketh the foundations of the earth, we may likewise feel thy mercy when thou healest the sores thereof. Through.
\secret
\lett{O}{God,} who hast laid the foundations of the earth that they should never be moved, receive the prayers and oblations of thy people: put far from us the present perils of earthquake, and turn the terrors of thy divine anger into a wholesome medicine for the safety of mankind; that they who are of the earth, and shall return to earth, may rejoice to be made citizens of heaven by holy conversation. Through.

\postcommunion
\lett{P}{rotect} us, O Lord, we beseech thee, who receive thy holy things: and by thy heavenly gift stablish the earth which we see trembling for our sins; that the hearts of men may know that in thine anger thou sendest forth thy scourges, and in thy mercy makest them to cease. Through.

   %MANUAL ADJUSTMENT:
   \newcolumn

\subsection{For Rain}
\collect
\lett{O}{God,} in whom we live, and move, and have our being: grant us such a seasonable rain; that we, receiving those things that be requisite for our necessities in this life, may the more confidently seek the things that are eternal. Through.

\secret
\lett{W}{e} beseech thee, O Lord, mercifully to receive the gifts which we offer: and send unto us a seasonable rain that may suffice for our necessities. Through.

\postcommunion
\lett{G}{rant} us, O Lord, we beseech thee, a refreshing rain: and graciously water the parched face of the earth by showers from heaven. Through.

\subsection{For Fine Weather}

\collect
\lett{G}{raciously} hear us, O Lord, who call upon thee: and grant fair weather in answer to our prayers; that we, who are justly afflicted for our sins, may by thy preventing mercy feel thy pity. Through.

   %MANUAL ADJUSTMENT:
   \newcolumn

\secret
\lett{L}{ord,} we pray thee, that thy grace may alway prevent and follow us: and graciously accept these oblations which for our sins we bring for consecration unto thy name; that, through the intercession of thy Saints, they may be profitable for the safety of us all. Through.

\postcommunion
\lett{A}{lmighty} God, we humbly beseech thy mercy: that thou wouldest assuage this plague of immoderate rains, and vouchsafe to shew forth on us the gladness of thy countenance. Through.


\subsection{Against Storms}

\collect
\lett{W}{e} beseech thee, O Lord, that all spiritual wickedness may be driven far from thy house: and that the malignity of the tempests of the air be likewise dispelled. Through.

\secret
\lett{W}{e} offer unto thee, O Lord, our praises and gifts, rendering thanks for the blessings bestowed upon us, and ever humbly praying for the continuance of the same. Through.

\postcommunion
\lett{A}{lmighty} and everlasting God, who dost heal us by thy chastisement, and preserve us by thy pardon: grant unto thy suppliant people; that in tranquillity of weather we may rejoice in the comfort we have ever desired, and may alway make use of thy bountiful goodness. Through.

\subsection{Against Cattle Disease}

\collect
\lett{O}{God,} who even by the dumb beasts hast brought comfort to man's labour: we humbly beseech thee; that forasmuch as without them the sustenance of man cannot but fail, so thou wouldest preserve them for our enjoyment. Through.

\secret
\lett{L}{ook} mercifully, O Lord, upon the sacrifices which we offer: and of thy mercy grant us thy help in our time. Through.

\postcommunion
\lett{G}{rant,} O Lord, that thy faithful people, being filled with thy blessing, may be preserved both in body and soul: that they may render unto thee worthy service, and evermore receive the benefits of thy mercy. Through.


   %MANUAL ADJUSTMENT:
   \newcolumn

\subsection{For the Priest Himself}

\collect
\lett{A}{lmighty} and merciful God, graciously give heed to my humble prayers: and make me thy servant, to whom for no merits of my own but by the boundless compassion of thy bounty thou hast vouchsafed to serve thee in these heavenly mysteries, a worthy minister at the sacred altars; that the words which I utter with my lips may by thy sanctifying power be brought to good effect. Through.

\secret
\lett{C}{leanse} me, O Lord, by the power of this sacrament, from all the stains of my sins: and grant; that with the help of thy grace I may be made worthy thereby to fulfil the ministry committed to my charge. Through.

\postcommunion
\lett{A}{lmighty} and everlasting God, who hast vouchsafed to me, a sinner, to stand before thy sacred altars, and to praise the power of thy holy name: mercifully grant to me by the mystery of this sacrament the pardon of my sins; that I may offer unto thy majesty true and worthy service. Through.


   %MANUAL ADJUSTMENT:
   \newcolumn


\subsection{To ask for Tears}
\collect
\lett{A}{lmighty} and most merciful God, who didst bring forth living water from the rock for thy thirsting people: draw out from our stony hearts the tears of repentance; that we, worthily lamenting our sins, may of thy mercy obtain forgiveness of the same. Through.
\secret
\lett{O}{Lord} God, we beseech thee, look down in mercy on this oblation, which we offer to thy majesty for our sins: and draw forth from our eyes such streams of tears, as shall quench the burning flames which we most righteously deserve. Through.
\postcommunion
\lett{O}{Lord} God, mercifully pour the grace of thy Holy Spirit into our heart: that we may be enabled thereby with tears of sorrow to wash away the defilements of our sins; that like as we do desire thy merciful pardon, so of thy bountiful goodness we may effectually receive the same. Through. . . in the unity of the same.

   %MANUAL ADJUSTMENT:
   \newcolumn

\subsection{For the Remission of Sins}
\collect
\lett{O}{God,} who dost not cast out sinners, and of thy loving mercy sparest them that are penitent, howsoever they sin against thee: mercifully incline to the prayers of thy humble servants, and enlighten our hearts; that we may be enabled to fulfil all that thou dost command. Through.
\secret
\lett{L}{et} this present sacrifice, O Lord, which we offer to thee for our offences, be an acceptable offering unto thee: and avail for salvation both to the living and the dead. Through.
\postcommunion
\lett{A}{lmighty} God, graciously hear the prayers of this thy family: and grant; that, like as we have received these holy things from thee, so of thy bounty they may be preserved pure and undefiled within us. Through.

\subsection{For Public Penitents}
\collect
\lett{A}{lmighty} and everlasting God, of thy great goodness forgive the sins of thy servants who confess the same to thee: that they whose consciences by sin are accused may not be delivered unto punishment, but may obtain of thy loving mercy pardon and forgiveness. Through.
\secret
\lett{G}{rant,} we beseech thee, almighty and merciful God: that this saving oblation may continually deliver thy servants from their offences, and defend them against all adversities. Through.
\postcommunion
\lett{A}{lmighty} and merciful God, who wouldest not the death of sinners, but rather that they should repent, and confess to thee, and amend their lives: look down upon these thy servants; and by these holy sacraments which we have received turn from them thy wrathful indignation, and grant them forgiveness of all their sins. Through.

\subsection{For the Tempted and Troubled}
\collect
\lett{O}{God,} who dost justify the wicked, and wouldest not the death of a sinner, we humbly beseech thy majesty: that thy servants who put their trust and confidence in thy mercy may be defended by the loving-kindness of thy heavenly succour and preserved by thy continual help; that in all temptations they may steadfastly cleave unto thy service. Through.
\secret
\lett{W}{e} beseech thee, O Lord, that by the power of this mystery thou wouldest cleanse us from our faults, and absolve thy servants from all their offences. Through.
\postcommunion
\lett{G}{rant,} we beseech thee, O Lord, that the sacraments which we have received may purify our hearts: and that thy servants may be delivered from all transgression; that whereas their consciences are bound by reason of their guilt, they may receive with gladness the plenteous healing of thy heavenly power. Through.

\subsection{To repel Evil Thoughts}
\collect
\lett{A}{lmighty} and most merciful God, graciously regard our prayers: that, our hearts being delivered from the temptations of evil thoughts; we may be found worthy to be made a dwelling-place acceptable to thy Holy Spirit. Through . . . in the unity of the same Holy Spirit.
\secret
\lett{W}{e} here present unto thee, O Lord, these offerings for our salvation: that the thoughts of our hearts being purged from all uncleanness and preserved from harm, we may be enlightened by the grace of thy Holy Spirit. Through . . . in the unity of the same Holy Spirit.
\postcommunion
\lett{O}{God,} who lightenest every man that cometh into this world: lighten, we beseech thee, our hearts with the brightness of thy grace; that we may ever be enabled to think such things as be good and pleasing to thy majesty, and that we may love thee with an unfeigned heart. Through.

\subsection{To obtain the Grace of Continence}
\collect
\lett{K}{indle,} O Lord, with the fire of the Holy Spirit our reins and our hearts: that we, serving thee in pureness both of body and soul, may be found pleasing in thy sight. Through . . . in the unity of the same.
\secret
\lett{S}{mite} asunder, O Lord, the chains of our sins: that we who seek to offer unto thee the sacrifice of praise in perfect freedom and in innocency of heart, may be restored by thee unto our first estate; that like as thou hast saved us by thy grace, so by thy pardon thou wouldest restore us again unto salvation. Through.
\postcommunion
\lett{A}{ssist} us, O God, our helper and defender: that our heart and flesh being renewed by the strength of continence and the newness of chastity, we who have offered unto thy majesty this sacrifice may be cleansed thereby from all temptations. Through.

\subsection{To obtain the Grace of Humility}
\collect
\lett{O}{God,} who resistest the proud, and givest grace unto the humble: grant to us the virtue of true humility, after the pattern wherewith thine only-begotten Son manifested himself unto thy faithful people; that we may in no wise exalt ourselves to provoke thine anger, but may in lowliness of heart obtain thy grace. Through.
\secret
\lett{W}{e} beseech thee, O Lord, that this oblation may obtain for us the grace of true humility: and likewise take from our hearts the lust of the flesh, and the lust of the eyes, and the pride of life; that living soberly, righteously, and godly, we may attain unto the rewards of everlasting life. Through.
\postcommunion
\lett{M}{ay} this sacrament, O Lord, which we have received, cleanse us from every stain of sin: that we, shewing forth humility, may be brought unto thy heavenly kingdom. Through.

\subsection{To obtain the Grace of Patience}
\collect
\lett{O}{God,} who by the long-suffering of thine only-begotten Son didst tread down the pride of our ancient enemy: grant us, we beseech thee, so rightly to call to mind the sufferings which of his goodness he bare for our sakes; that after his example we may patiently endure all our troubles and adversities. Through the same.
\secret
\lett{W}{e} beseech thee, O Lord, graciously to receive our gifts and oblations: which we devoutly offer to thy majesty, beseeching thee that thou wouldest vouchsafe to bestow upon us the gift of patience. Through.
\postcommunion
\lett{W}{e} beseech thee, O Lord, that we who have received these sacred mysteries may be restored thereby unto the grace from which we have fallen: that by thy continual protection both here and ever we may be enabled to endure patiently in all adversities. Through.

\subsection{To obtain the Grace of Charity}
\collect
\lett{O}{God,} who makest all things work together for good to them that love thee: pour into our hearts the steadfast affection of thy charity; that as thou dost put into our minds good desires, so we may not be turned from the same by any temptations. Through.
\secret
\lett{O}{God,} who by the sacraments and commandments dost renew us after thy likeness: stablish our goings in thy paths; that like as thou hast vouchsafed unto us the hope of obtaining thy gift of charity, so by these sacrifices which we offer thou wouldest make us effectually to be partakers of the same. Through.
\postcommunion
\lett{W}{e} beseech thee, O Lord, that the grace of the Holy Spirit may enlighten our hearts: that they may be abundantly renewed in the sweetness of perfect charity. Through . . . in the unity of the same Holy Spirit.

   %MANUAL ADJUSTMENT:
   \newcolumn

\subsection{For Devoted Friends}
\collect
\lett{O}{God,} who by the grace of thy Holy Spirit hast endued the hearts of thy faithful people with the gifts of charity: grant to thy servants and handmaids, for whom we implore thy mercy, safety of mind and body; that they may love thee with all their strength, and with pure affection perform such things as are acceptable unto thee. Through . . . in the unity of the same.
\secret
\lett{W}{e} beseech thee, O Lord, to have compassion upon thy servants and handmaids, for whom we offer unto thy majesty this sacrifice of praise: that through these holy things they may obtain the grace of thy heavenly benediction, and attain unto the glory of everlasting felicity. Through.
\postcommunion
\lett{O}{Lord,} who dost make us to taste of these divine mysteries, we humbly beseech thee: that like as we have offered these sacraments of salvation for the advancement in charity of thy servants; so they may be profitable unto them for their prosperity and peace. Through.

   %MANUAL ADJUSTMENT:
   \newcolumn

\subsection{For Enemies}
\collect
\lett{O}{God,} who art the lover of peace, and preserver of charity: grant unto all our enemies true peace and charity: and vouchsafe unto them remission of all their sins, and by thy mighty power deliver us from their snares. Through.
\secret
\lett{W}{e} beseech thee, O Lord, graciously to accept the gifts which we offer: that by thy mercy we may be delivered from our enemies, and that they may receive of thee the pardon of their offences. Through.
\postcommunion
\lett{M}{ay} this communion, O Lord, deliver us from our offences: and defend us against the snares of our enemies. Through.

\subsection{For a Prisoner or Captive}
\collect
\lett{O}{God,} who didst deliver blessed Peter the Apostle from his chains, and suffer him to depart unhurt: loose the chains of thy servant who now lieth in captivity; and by the merits of the same thine Apostle suffer him to depart unhurt. Through.

   %MANUAL ADJUSTMENT:
   \newcolumn

\secret
\lett{L}{et} thy plenteous benediction, O Lord, descend upon these sacrifices: that they may loose the chains of this thy servant in captivity, and make us speedily to rejoice in his deliverance. Through.

\postcommunion
\lett{O}{Lord,} we beseech thee, graciously hear our prayers: and by these sacraments which we have received deliver thy servant from the chains of his captivity. Through.

\subsection{For Those at Sea}
\collect
\lett{O}{God,} who didst guide our fathers through the Red Sea, and didst bring them through the deep waters, singing praises to thy name: We humbly beseech thee; that thou wouldest defend thy servants now travelling by sea from all adversity, so that they, voyaging in tranquillity, may be brought by thee unto the haven where they would be. Through.
\secret
\lett{W}{e} beseech thee, O Lord, to receive the prayers of thy servants, with our sacrifices and oblations: that they for whom we celebrate these mysteries may be defended by thee against all adversities. Through.
\postcommunion
\lett{O}{Lord,} who hast sanctified us with this divine mystery, we humbly entreat and beseech thy majesty: that as thou dost suffer us to call to remembrance thy servants by means of these thy heavenly gifts; so through the wood of the holy Cross thou wouldest absolve them from their sins and deliver them from all perils and dangers. Through.

\subsection{For the Safety of the Living}
\collect
\lett{S}{tretch} forth, O Lord, upon thy faithful people the right hand of thy heavenly succour: that they, seeking thee with their whole heart, may be found worthy to obtain those things which they ask according to thy will. Through.
\secret
\lett{W}{e} beseech thee, O Lord, to have compassion on our prayers, and graciously to accept these oblations of thy faithful people, which we present before thee for their safety: and that none may ask amiss or fail in his petition; grant that those things which we ask faithfully, we may obtain effectually. Through.
\postcommunion
\lett{G}{rant,} we beseech thee, O Lord, unto thy faithful people steadfastness in faith and singleness of heart: that they, being stablished in divine charity, may not by any temptations be parted from the fulness of the same. Through.

\subsection{For the Living \& Departed}
\collect
\lett{A}{lmighty} and everlasting God, who hast dominion both of the living and of the dead, and hast mercy upon all whom thou foreknowest will be thine in faith and works: we humbly beseech thee; that all those for whom we are minded to pour forth our prayers, whether in this present world they still be held in the flsh, or being delivered from the body have passed into that which is to come, may at the intercession of all thy Saints obtain of thy bountiful goodness the remission of all their sins. Through.
\secret
\lett{O}{God,} to whom alone is known the number of the elect, whom thou hast appointed unto heavenly felicity: grant, we beseech thee; that at the intercession of all thy Saints, the names of all those who have been commended to our prayers, and of all thy faithful people, may ever remain written in the book of those that are predestined to everlasting blessedness. Through.

   %MANUAL ADJUSTMENT:
   \newcolumn

\postcommunion
\lett{A}{lmighty} and merciful God, cleanse us we pray thee, by the sacraments which we have received: and, at the intercession of all thy Saints, grant; that this thy sacrament may not bring upon us guilt to our punishment but pardon to our salvation: that it may cleanse away sins, that it may confirm the feeble, that it may defend us against every danger in this world: that it may obtain for thy
faithful people both living and dead the remission of all their sins. Through.
\end{multicols}

\begin{rubric}
    The following propers are used for the faithful departed.
\end{rubric}

\begin{multicols}{2}
\subsection{For a Chief Bishop Departed}\label{RequiemPrayers}
\collect
\lett{O}{God,} who of thy unspeakable providence didst vouchsafe to number thy servant \textit{N.} in the company of thy high Priests: grant, we beseech thee; that like as he did fulfil on earth the office of thine Only-begotten Son, so he may be numbered in the everlasting fellowship of the same thy Saints. Through the same.

\secret
\lett{R}{eceive,} O Lord, we beseech thee, these our sacrifices for the soul of thy servant \textit{N.}, thy high Priest: that forasmuch as in this life thou didst give him the dignity of a Bishop, so thou wouldest bid him to be joined to the company of the Saints in thy heavenly kingdom. Through.
\postcommunion%Pope changed to Metropolitan
\lett{M}{ay} thy merciful loving-kindness, which we implore, benefit, we beseech thee, O Lord, the soul of thy servant the \textit{Metropolitan} \textit{N.}: that, by thy mercy he may attain the eternal fellowship of him in whom he hoped and believed. Through.

\subsection{For a Departed Bishop}
\collect
\lett{O}{God,} who didst cause thy servant \textit{N.} to enjoy the dignity of a Bishop in the apostolic Priesthood: grant, we beseech thee; that he may evermore be joined unto the fellowship of the same. Through.
\secret
\lett{A}{ccept,} O Lord, we beseech thee, the sacrifice which we offer to thee for the soul of thy servant and Bishop \textit{N.}: that forasmuch as in this life thou didst bestow on him the dignity of a Bishop, so thou wouldest bid him to be joined unto the fellowship of thy Saints in the heavenly kingdom. Through.

   %MANUAL ADJUSTMENT:
   \newcolumn

\postcommunion
\lett{W}{e} beseech thee, O Lord, that thy merciful loving kindness, which we implore, may be profitable for the soul of thy servant and Bishop \textit{N.}: that by thy mercy he may attain unto the eternal fellowship of him in whom he hath hoped and believed. Through.
\begin{inhead}
or,
\end{inhead}
\vspace{-1\baselineskip}
\collect
\lett{G}{rant} unto us, O Lord, that the \textit{soul} of thy \textit{servant} and \textit{Bishop} \textit{N.} (and \textit{N.}), which thou hast taken out of the toil and conflict of this world, may be \textit{partaker} of the fellowship of thy Saints. Through.


\secret
\lett{G}{rant} unto us, we beseech thee, O Lord: that like as by the offering of this oblation thou didst grant unto the whole world remission of sins so the \textit{soul} of thy \textit{servant} and \textit{Bishop} \textit{N.} (and \textit{N.}) may be profited by the same. Through.

\postcommunion
\lett{W}{e} beseech thee, almighty God, that the \textit{soul} of thy \textit{servant} and \textit{Bishop} \textit{N.} (and \textit{N.}), being purified by these sacrifices, may be found worthy to obtain thy pardon and everlasting refreshment. Through.

   %MANUAL ADJUSTMENT:
   \newcolumn

\subsection{For a Departed Priest}
\collect
\lett{O}{God,} who didst cause thy \textit{servant} \textit{N.} (and \textit{N.}) to enjoy the dignity of a Priest in the apostolic Priesthood: grant, we beseech thee, that \textit{he} may evermore be joined unto the fellowship of the same. Through.
\secret
\lett{A}{ccept,} O Lord, we beseech thee, the sacrifice which we offer to thee for the \textit{soul} of thy \textit{servant} and \textit{Priest} \textit{N.} (and \textit{N.}): that forasmuch as in this life thou didst bestow on him the dignity of a Priest, so thou wouldest bid him to be joined unto the fellowship of thy Saints in the heavenly kingdom.
\postcommunion
\lett{W}{e} beseech thee, O Lord, that thy merciful loving kindness, which we implore, may be profitable for the \textit{soul} of thy \textit{servant} and \textit{Priest} \textit{N.} (and \textit{N.}): that by thy mercy \textit{he} may attain unto the eternal fellowship of him in whom in whom \textit{he} hath hoped and believed. Through.
\begin{inhead}
    or,
\end{inhead}
\vspace{-1\baselineskip}
\collect
\lett{G}{rant,} we beseech thee, O Lord, that the \textit{soul} of thy \textit{servant} \textit{N.} (and \textit{N.}): whom while \textit{he} lived in this world thou didst adorn with the holy gifts of Priesthood; may ever rejoice in the glory of thy heavenly abode. Through.
\secret
\lett{R}{eceive,} O Lord, we beseech thee, these sacrifices, which we offer for the \textit{soul} of \textit{N.} (and \textit{N.}) thy \textit{servant} and \textit{Priest}: that as thou hast given \textit{him} the dignity of the Priesthood, thou wouldest also give him the reward of the same. Through.
\postcommunion
\lett{G}{rant,} we beseech thee, almighty God: to the \textit{soul} of \textit{N.} (and \textit{N.}) thy \textit{servant} and \textit{Priest}, the fellowship of eternal blessedness in the congregation of the just. Through.

\subsection{For a Man Departed}
\collect
\lett{I}{ncline} thine ear, O Lord, unto the prayers wherewith we humbly entreat thy mercy: that thou wouldest set the \textit{soul} of thy \textit{servant} \textit{N.} (and \textit{N.}), which thou hast bidden depart this life, in a place of peace and light; and make \textit{him} to be joined unto the fellowship of thy Saints. Through.
\secret
\lett{G}{rant} to us, we beseech thee, O Lord: that like as by the offering of this oblation thou didst bestow upon the whole world remission of sins: so the \textit{soul} of thy \textit{servant} \textit{N.} (and \textit{N.}) may be profited by the same. Through.
\postcommunion
\lett{A}{bsolve,} O Lord, we beseech thee, the \textit{soul} of thy \textit{servant} \textit{N.} (and \textit{N.}), from every bond of sin: that in the glory of the resurrection \textit{he} may be raised up amid thy Saints and elect unto newness of life. Through.


\subsection{For a Woman Departed}
\collect
\lett{W}{e} beseech thee, O Lord, of thy loving-kindness to have mercy on the \textit{soul} of thy \textit{handmaid} \textit{N.} (and \textit{N.}): that being purged from the defilements of our mortal nature, \textit{she} may be restored to the portion of everlasting salvation. Through.
\secret
\lett{W}{e} beseech thee, O Lord, that by these sacrifices, whereby alone we have deliverance from our sins, the \textit{soul} of thy \textit{handmaid} \textit{N.} (and \textit{N.}) may be absolved from all \textit{her} offences: that through this our bounden duty of propitiation she may obtain thine everlasting mercy. Through.
\postcommunion
\lett{W}{e} beseech thee, O Lord, that forasmuch as the \textit{soul} of thine \textit{handmaid} \textit{N.} (and \textit{N.}) hath received this pledge of thy never-failing mercy: \textit{she} may be made partaker of thine everlasting brightness. Through.

\subsection{For a Father and Mother}
\collect
\lett{O}{God,} who hast bidden us to honour our father and our mother: of thy loving-kindness have mercy on the souls of \textit{my father and mother} and forgive them their sins; and grant that \textit{I} may behold them in the joy of the eternal brightness. Through.
\secret
\lett{R}{eceive,} O Lord, the sacrifice which I offer unto thee for the souls of \textit{my father and mother}: and grant them in the land of the living everlasting blessedness; and join \textit{me} with them in the felicity of thy Saints. Through.
\postcommunion
\lett{W}{e} beseech thee, O Lord, that this partaking of thy heavenly sacrament may obtain for the souls of \textit{my father and mother} rest and light perpetual: and that with them \textit{I} may receive the crown of thy everlasting grace. Through.

   %MANUAL ADJUSTMENT:
   \newcolumn

\subsection{For a Father Only}
\collect
\lett{O}{God,} who hast bidden us to honour our father and our mother: of thy loving-kindness have mercy on the soul of \textit{my} father and forgive him his sins; and grant that \textit{I} may behold him in the joy of the eternal brightness. Through.
\secret
\lett{R}{eceive,} O Lord, the sacrifice which I offer unto thee for the soul of \textit{my} father: and grant him in the land of the living everlasting blessedness; and join \textit{me} with him in the felicity of thy Saints. Through.
\postcommunion
\lett{W}{e} beseech thee, O Lord, that this partaking of thy heavenly sacrament may obtain for the soul of \textit{my} father rest and light perpetual; and that with him \textit{I} may receive the crown of thy everlasting grace. Through.

\subsection{For a Mother Only}
\collect
\lett{O}{God,} who hast bidden us to honour our father and our mother: of thy loving-kindness have mercy on the soul of \textit{my} mother and forgive her her sins; and grant that \textit{I} may behold her in the joy of the eternal brightness. Through.
\secret
\lett{R}{eceive,} O Lord, the sacrifice which \textit{I} offer unto thee for the soul of \textit{my} mother: and grant her in the land of the living everlasting blessedness; and join \textit{me} with her in the felicity of thy Saints. Through.
\postcommunion
\lett{W}{e} beseech thee, O Lord, that this partaking of thy heavenly sacrament may obtain for the soul of \textit{my} mother rest and light perpetual: and that with her \textit{I} may receive the crown of thy everlasting grace. Through.

\subsection{For Departed Brethren, Kinsfolk, \& Benefactors}
\collect
\lett{O}{God,} the giver of pardon and lover man's salvation: we beseech thee of thy mercy to grant; that the brethren, kinsfolk, and benefactors of our congregation who have passed out of this world, may at the intercession of blessed Mary, ever Virgin, and of all thy Saints attain unto the fellowship of eternal blessedness. Through.

   %MANUAL ADJUSTMENT:
   \newcolumn

\secret
\lett{O}{God,} whose mercies are with-out number, graciously receive our humble prayers: and through these sacraments of our salvation grant unto the souls of our brethren, kinsfolk, and benefactors, to whom thou hast given the confession of thy name, the remission of all their sins. Through.
\postcommunion
\lett{G}{rant,} we beseech thee, al-mighty and merciful God: that the souls of our brethren, kinsfolk, and benefactors, for whom we have offered unto thy majesty this sacrifice of praise; may by the power of this sacrament be cleansed from all their sins, and by thy mercy obtain the blessedness of everlasting light.
Through.

\subsection{For Those who Rest in a Cemetery}
\collect
\lett{O}{God,} through whose mercy the souls of the faithful are at rest: mercifully grant to thy servants and handmaids, and to all that here and in all places do rest in Christ, the remission of their sins; that being delivered from all their iniquities, they may rejoice with thee for evermore. Through the same.
\secret
\lett{G}{raciously} receive, O Lord, this sacrifice which we offer for the souls of thy servants and handmaids, and of all Christian people who here and in all places are asleep in Christ: that by this wondrous sacrifice they may be delivered from the dreadful bonds of death, and be found worthy to enter into everlasting life. Through the same.

%MANUAL ADJUSTMENT:
\vspace{-0.25\baselineskip}

\postcommunion
\lett{O}{God,} the light of the souls of the faithful, assist our supplications: and grant unto thy servants and handmaids, whose bodies both here and in all places do rest in Christ, a place of refreshing, the blessedness of thy rest, and the glory of everlasting light. Through the same.

%MANUAL ADJUSTMENT:
\vspace{-0.25\baselineskip}

\subsection{For Many Persons Departed}

%MANUAL ADJUSTMENT:
\vspace{-0.25\baselineskip}

\collect
\lett{O}{God,} whose property is ever to have mercy and to forgive: have compassion upon the souls of thy servants and handmaids, and forgive them all their sins; that being delivered from the bonds of this our mortal nature, they may be found worthy to enter into life. Through.

%MANUAL ADJUSTMENT:
\vspace{-0.25\baselineskip}

\secret
\lett{G}{rant} to us, we beseech thee, O Lord, that like as by the offering of this oblation thou didst bestow upon the whole world remission of sins: so the souls of thy servants and handmaids may be profited by the same. Through.

%MANUAL ADJUSTMENT:
\vspace{-0.25\baselineskip}

\postcommunion
\lett{O}{God,} to whom alone it belongeth to bestow healing after death: grant, we beseech thee; that the souls of thy servants and handmaids, being cleansed from the defilements of this earthly life, may be numbered in the portion of thy redeemed. Who livest and reignest.

%MANUAL ADJUSTMENT:
\vspace{-0.25\baselineskip}

\subsection{Other Prayers for Many Persons Departed}

%MANUAL ADJUSTMENT:
\vspace{-0.25\baselineskip}

\collect
\lett{G}{rant,} O Lord, we pray thee, to the souls of thy servants and handmaids thy perpetual mercy: that as they have hoped and believed in thee, so it may be profitable unto them for evermore. Through.

%MANUAL ADJUSTMENT:
\vspace{-0.25\baselineskip}

\secret
\lett{W}{e} beseech thee, O Lord, mercifully to look upon these our gifts: that these things which we humbly offer to the glory of thy name, may be profitable unto them that are departed for the forgiveness of their sins. Through.

%MANUAL ADJUSTMENT:
\vspace{-0.25\baselineskip}

\postcommunion
\lett{W}{e} humbly offer unto thee, O Lord, our prayers for the souls of thy servants and handmaids: beseeching thee; that whatsoever they may have contracted in their earthly conversation being pardoned by thy goodness, they may be made partakers of the gladness of thy redeemed. Through.
\end{multicols}
%Votive Masses
\cleardoublepage
\fancyhead[RE,LO]{}\fancyhead[RO,LE]{}
\fancyhead[C]{}\thispagestyle{empty}
\phantomsection
\addcontentsline{toc}{chapter}{Votive Masses}

\begin{figure}[H]
    \centering
    \includegraphics[width=98mm, height=140mm]{Votives/MassStGregory.eps}
    \caption{\textsc{\Huge{Votive Masses}}}
\end{figure}
\include{pages/Votives/Votive}
%Pastoral Rites
\cleardoublepage
\include{pages/Pastoral/PastoralCover}
\include{pages/Pastoral/Baptism}
\include{pages/Pastoral/AdultBaptism}
\include{pages/Pastoral/Confession}
\include{pages/Pastoral/Reception}
\include{pages/Pastoral/Confirmation}
\include{pages/Pastoral/Matrimony}
\include{pages/Pastoral/Purification}
\include{pages/Pastoral/Unction}
\include{pages/Pastoral/CommunionOfTheSick}
\fancyhead[RO,LE]{\textit{Burial of the Dead}}
\fancyhead[RE,LO]{\textit{}}
%This rite is revised, in accordance with the Observations, in light of A Manual for Priests.
\section{The Order for the Burial of the Dead}

\subsection{Reception of the Body}
%Clearer rubric from 1928 English BCP:
\begin{secrubric}
    The Priest, meeting the Body at the entrance of the church yard, and prior to going before it either into the church or towards the grave, shall say or sing,\par
    \textsc{Note,} The Priest should be vested in surplice and black stole, and also, on more solemn occasions, in black cope.
\end{secrubric}
\lett{I}{am} the resurrection and the life, saith the Lord: he that believeth in me, though he were dead, yet shall he live: and whosoever liveth and believeth in me, shall never die.\par
    \secondline{I know that my redeemer liveth, and that he shall stand at the latter day upon the earth: and though this body be destroyed, yet shall I see God: whom I shall see for myself, and mine eyes shall behold, and not as a stranger.}
    \thirdline{We brought nothing into this world, and it is certain we can carry nothing out. The \divineName{Lord} gave, and the \divineName{Lord} hath taken away; blessed be the name of the \divineName{Lord}.}

\subsection{Procession}
\begin{rubric}
	A Clerk carrying the Cross goes before (and, on more solemn occasions, with two acolytes with candles), accompanied by another bearing Holy Water. As the Ministers and Clerks proceed into the Church, the following Responsory is begun.
\end{rubric}

\subsubsection{Subvenite}
\lett{C}{ome} to \textit{his} aid, ye Saints of God, come to meet \textit{him}, ye Angels of the Lord. Receiving \textit{his} soul, offering it in the sight of the Most High.\par
\secondline{℣. May Christ receive thee, who hath called thee: and may the Angels lead thee unto Abraham's bosom.}
\thirdline{℟. Receiving \textit{his} soul, offering it in the sight of the Most High.}
℣. Rest eternal grant unto \textit{him}, O Lord: and let light perpetual shine upon \textit{him}.\par
℟. Offering it in the sight of the Most High.

\begin{rubric}
	The Coffin is set down in the midst of the Church, so that the feet of the departed, unless he be a Priest, are toward the High Altar; but if he be a Priest, his head is toward the Altar. Then Candles should be lighted about the Body.
\end{rubric}

\subsection{Service in the Church}
\begin{rubric}
%Addition of prayers for the dead and traditional replacement of the Gloria Patri in Requiem contexts.
    After they are come into the Church, and the Responsory hath been said, then shall be said one or more or the following Selections, taken from the Psalms. Instead of the \blackrubric{Gloria Patri} shall be said the following at the end of each Psalm.
\end{rubric}
℣. Rest eternal * grant unto them, O Lord.\par
℟. And let light perpetual * shine upon them.\par\noindent

\subsubsection{Psalm 39. \textit{Dixi, Custodiam}}
\lett{L}{ord,} let me know mine end, and the number of my days : that I may be certified how long I have to live.\par
\secondline{\psanum{6}Behold, thou hast made my days as it were a span long : and mine age is even as nothing in respect of thee; and verily every man living is altogether vanity.}
\thirdline{\psanum{7}For man walketh in a vain shadow, and disquieteth himself in vain : he heapeth up riches, and cannot tell who shall gather them.}
\psanum{8}And now, Lord, what is my hope : truly my hope is even in thee.\par
\psanum{9}Deliver me from all mine offences : and make me not a rebuke unto the foolish.\par
\psanum{10}I became dumb, and opened not my mouth : for it was thy doing.\par
\psanum{11}Take thy plague away from me : I am even consumed by the means of thy heavy hand.\par
\psanum{12}When thou with rebukes dost chasten man for sin, thou makest his beauty to consume away, like as it were a moth fretting a garment : every man therefore is but vanity.\par
\psanum{13}Hear my prayer, O Lord, and with thine ears consider my calling : hold not thy peace at my tears.\par
\psanum{14}For I am a stranger with thee : and a sojourner, as all my fathers were.\par
\psanum{15}O spare me a little, that I may recover my strength : before I go hence, and be no more seen.\par

\subsubsection{Psalm 90. \textit{Domine, refugium}}
\lett{L}{ord,} thou hast been our refuge : from one generation to another.\par
%Third line too long:
\secondline{\psanum{2}Before the mountains were brought forth, or ever the earth and the world were made : thou art God from everlasting, and world without}
\noindent
end.\par
\psanum{3}Thou turnest man to destruction : again thou sayest, Come again, ye children of men.\par
\psanum{4}For a thousand years in thy sight are but as yesterday : seeing that is past as a watch in the night.\par
\psanum{5}As soon as thou scatterest them they are even as a sleep : and fade away suddenly like the grass.\par
\psanum{6}In the morning it is green, and groweth up : but in the evening it is cut down, dried up, and withered.\par
\psanum{7}For we consume away in thy displeasure : and are afraid at thy wrathful indignation.\par
\psanum{8}Thou hast set our misdeeds before thee : and our secret sins in the light of thy countenance.\par
\psanum{9}For when thou art angry all our days are gone : we bring our years to an end, as it were a tale that is told.\par
\psanum{10}The days of our age are threescore years and ten; and though men be so strong that they come to fourscore years : yet is their strength then but labour and sorrow; so soon passeth it away, and we are gone.\par
\psanum{11}But who regardeth the power of thy wrath : for even thereafter as a man feareth, so is thy displeasure.\par
\psanum{12}So teach us to number our days : that we may apply our hearts unto wisdom.\par

\subsubsection{Psalm 27. \textit{Dominus illuminatio}}
\lett{T}{he} Lord is my light and my salvation ; whom then shall I fear : the Lord is the strength of my life; of whom then shall I be afraid?\par
\secondline{\psanum{2}When the wicked, even mine enemies and my foes, came upon me to eat up my flesh : they stumbled and fell.}
\thirdline{\psanum{3}Though an host of men were laid against me, yet shall not my heart be afraid : and though there rose up war against me, yet will I put my trust in him.}
\psanum{4}One thing have I desired of the Lord, which I will require : even that I may dwell in the house of the Lord all the days of my life, to behold the fair beauty of the Lord, and to visit his temple.\par
\psanum{5}For in the time of trouble he shall hide me in his tabernacle : yea, in the secret place of his dwelling shall he hide me, and set me up upon a rock of stone.\par
\psanum{6}And now shall he lift up mine head : above mine enemies round about me.\par
\psanum{7}Therefore will I offer in his dwelling an oblation with great gladness : I will sing, and speak praises unto the Lord.\par
\psanum{8}Hearken unto my voice, O Lord, when I cry unto thee : have mercy upon me, and hear me.\par
\psanum{9}My heart hath talked of thee, Seek ye my face : Thy face, Lord, will I seek.\par
\psanum{10}O hide not thou thy face from me : nor cast thy servant away in displeasure.\par
\psanum{11}Thou hast been my succour : leave me not, neither forsake me, O God of my salvation.\par
\psanum{12}When my father and my mother forsake me : the Lord taketh me up.\par
\psanum{13}Teach me thy way, O Lord : and lead me in the right way, because of mine enemies.\par
\psanum{14}Deliver me not over into the will of mine adversaries : for there are false witnesses risen up against me, and such as speak wrong.\par
\psanum{15}I should utterly have fainted : but that I believe verily to see the goodness of the Lord in the land of the living.\par
\psanum{16}O tarry thou the Lord's leisure : be strong, and he shall comfort thine heart; and put thou thy trust in the Lord.\par

\subsubsection{Psalm 46. \textit{Deus noster refugium}}
\lett{G}{od} is our hope and strength : a very present help in trouble.\par
\secondline{\psanum{2}Therefore will we not fear, though the earth be moved : and though the hills be carried into the midst of the sea;}
\thirdline{\psanum{3}Though the waters thereof rage and swell : and though the mountains shake at the tempest of the same.}
\psanum{4}The rivers of the flood thereof shall make glad the city of God : the holy place of the tabernacle of the most Highest.\par
\psanum{5}God is in the midst of her, therefore shall she not be removed : God shall help her, and that right early.\par
\psanum{6}The heathen make much ado, and the kingdoms are moved : but God hath shewed his voice, and the earth shall melt away.\par
\psanum{7}The Lord of hosts is with us : the God of Jacob is our refuge.\par
\psanum{8}O come hither, and behold the works of the Lord : what destruction he hath brought upon the earth.\par
\psanum{9}He maketh wars to cease in all the world : he breaketh the bow, and knappeth the spear in sunder, and burneth the chariots in the fire.\par
\psanum{10}Be still then, and know that I am God : I will be exalted among the heathen, and I will be exalted in the earth.\par
\psanum{11}The Lord of hosts is with us : the God of Jacob is our refuge.\par

\subsubsection{Psalm 121. \textit{Levavi oculus}}
\lett{I}{will} lift up mine eyes unto the hills : from whence cometh my help.\par
\secondline{\psanum{2}My help cometh even from the Lord : who hath made heaven and earth.}
\thirdline{\psanum{3}He will not suffer thy foot to be moved : and he that keepeth thee will not sleep.}
\psanum{4}Behold, he that keepeth Israel : shall neither slumber nor sleep.\par
\psanum{5}The Lord himself is thy keeper : the Lord is thy defence upon thy right hand;\par
\psanum{6}So that the sun shall not burn thee by day : neither the moon by night.\par
\psanum{7}The Lord shall preserve thee from all evil : yea, it is even he that shall keep thy soul.\par
\psanum{8}The Lord shall preserve thy going out, and thy coming in : from this time forth for evermore.\par

\subsubsection{Psalm 130. \textit{De profundis}}
\lett{O}{ut} of the deep have I called unto thee, O Lord : Lord, hear my voice.\par
\secondline{\psanum{2}O let thine ears consider well : the voice of my complaint.}
\thirdline{\psanum{3}If thou, Lord, wilt be extreme to mark what is done amiss : O Lord, who may abide it?}
\psanum{4}For there is mercy with thee : therefore shalt thou be feared.\par
\psanum{5}I look for the Lord; my soul doth wait for him : in his word is my trust.\par
\psanum{6}My soul fleeth unto the Lord : before the morning watch, I say, before the morning watch.\par
\psanum{7}O Israel, trust in the Lord, for with the Lord there is mercy : and with him is plenteous redemption.\par
\psanum{8}And he shall redeem Israel : from all his sins.\par

%Verses 29-34 restored, as in 1892 BCP:
\readingcitation{Lesson}{1 Corinthians 15:20}
\begin{rubric}
	Then shall follow the Lesson.
\end{rubric}
\begin{rubric}
	\textsc{Note,} Either Romans 8:14-39 or John 14:1-6 may be said instead.
\end{rubric}
\lett{N}{ow} is Christ risen from the dead, and become the first-fruits of them that slept. For since by man came death, by man came also the resurrection of the dead. For as in Adam all die, even so in Christ shall all be made alive. But every man in his own order: Christ the first-fruits; afterward they that are Christ's, at his coming. Then cometh the end, when he shall have delivered up the kingdom to God, even the Father; when he shall have put down all rule, and all authority, and power. For he must reign, till he hath put all enemies under his feet. The last enemy that shall be destroyed is death. For he hath put all things under his feet. But when he saith, all things are put under him, it is manifest that he is excepted, which did put all things under him. And when all things shall be subdued unto him, then shall the Son also himself be subject unto Him that put all things under him that God may be all in all.\par
Else what shall they do which are baptized for the dead, if the dead rise not at all? Why are they then baptized for the dead? and why stand we in jeopardy every hour? I protest by your rejoicing, which I have in Christ Jesus our Lord, I die daily. If after the manner of men I have fought with beasts at Ephesus, what advantageth it me, if the dead rise not? let us eat and drink, for tomorrow we die. Be not deceived: evil communications corrupt good manners. Awake to righteousness, and sin not; for some have not the knowledge of God. I speak this to your shame. But some man will say, How are the dead raised up? and with what body do they come? Thou fool! that which thou sowest is not quickened, except it die. And that which thou sowest, thou sowest not that body that shall be, but bare grain, it may chance of wheat, or of some other grain. But God giveth it a body as it hath pleased him, and to every seed his own body. All flesh is not the same flesh; but there is one kind of flesh of men, another flesh of beasts, another of fishes, and another of birds. There are also celestial bodies, and bodies terrestrial; but the glory of the celestial is one, and the glory of the terrestrial is another. There is one glory of the sun, and another glory of the moon, and another glory of the stars; for one star differeth from another star in glory. So also is the resurrection of the dead. It is sown in corruption; it is raised in incorruption: it is sown in dishonour; it is raised in glory: it is sown in weakness; it is raised in power: it is sown a natural body; it is raised a spiritual body. There is a natural body, and there is a spiritual body. And so it is written, The first man Adam was made a living soul; the last Adam was made a quickening spirit. Howbeit, that was not first which is spiritual, but that which is natural; and afterward that which is spiritual. The first man is of the earth, earthy: the second man is the Lord from heaven. As is the earthy, such are they also that are earthy: and as is the heavenly, such are they also that are heavenly. And as we have borne the image of the earthy, we shall also bear the image of the heavenly.\par
Now this I say, brethren, that flesh and blood cannot inherit the kingdom of God; neither doth corruption inherit incorruption Behold, I show you a mystery: we shall not all sleep, but we shall all be changed, in a moment, in the twinkling of an eye, at the last trump: for the trumpet shall sound, and the dead shall be raised incorruptible, and we shall be changed. For this corruptible must put on incorruption, and this mortal must put on immortality. So when this corruptible shall have put on incorruption, and this mortal shall have put on immortality; then shall be brought to pass the saying that is written, Death is swallowed up in victory. O death, where is thy sting? O grave, where is thy victory? The sting of death is sin; and the strength of sin is the Law. But thanks be to God, which giveth us the victory through our Lord Jesus Christ. Therefore, my beloved brethren, be ye steadfast, unmoveable, always abounding in the work of the Lord, forasmuch as ye know that your labour is not in vain in the Lord.
 
 %Requiring the absolution:
\begin{rubric}
	Here may be sung a Hymn or Anthem; and at the discretion of the Minister, the Apostles' Creed, the Lord's Prayer, and such other fitting Prayers as are elsewhere provided in this Book, ending with the prayers which followeth and the Blessing; the Minister, before the Prayers, first pronouncing,
\end{rubric}

℣. The Lord be with you.

℟. And with thy spirit.

\letuspray
\lett{A}{bsolve,} we beseech thee, O Lord, the soul of thy \textit{servant} \textit{N.} from every bond of sin: that in the glory of the resurrection \textit{he} may be raised up amid thy Saints and elect unto newness of life. Through Christ, our Lord. \textit{Amen.}

\lett{R}{emember} thy servant, O Lord, according to the favour which thou bearest unto thy people, and grant that, increasing in knowledge and love of thee, \textit{he} may go from strength to strength, in the life of perfect service, in thy heavenly kingdom; through Jesus Christ our Lord, who liveth and reigneth with thee and the Holy Ghost, ever, one God, world without end. \textit{Amen.}

\lett{U}{nto} God's gracious mercy and protection we commit you. The Lord bl {\ding{64}} ess you and keep you. The Lord make his face to shine upon you, and be gracious unto you. The Lord lift up his countenance upon you, and give you peace, both now and evermore. \textit{Amen.}

%Moved here, contrary to the Manual for Priests in order to create a definite marker and conclusion for the Office prior.
\begin{rubric}
    If a Requiem Mass be said, it may be offered here (p. \pageref{RequiemMasses}).
\end{rubric}

\vfill

  \begin{figure}[H]
  	\centering
  	\includegraphics[scale=0.2]{tailpiece/MadonnaChildMoon.eps}
  \end{figure}


%MANUAL ADJUSTMENT:
\clearpage
\subsection{Absolution of the Dead}     
% \begin{rubric}
%     At the end of the Requiem Mass, the Celebrant retires to the Epistle Corner, where he takes off his chasuble and maniple, and puts on a black cope. The Subdeacon, between two Acolytes with lighted candles, carries the cross as in Processions, two other Acolytes preceding them, one with the thurible and incense-boat, the other with the aspersorium and aspergilium. The Celebrant follows, having first made a reverence to the Altar, with the Deacon on his left. The Subdeacon with the Cross takes his stand at the feet of the bier or catafalque opposite the altar, between the aforesaid Acolytes holding candles, while the Celebrant stands on the other side at the head of the place between the Altar and the bier, turning a little towards the Epistle corner, so that he looks towards the Subdeacon's Cross; on his left the Deacon, and near him the two other Acolytes carrying the thurible and aspersorium. Then, an Acolyte or Clerk holding the book, the Celebrant says at once the following Collect no change of number or sex being made, even if it is said for many persons or a woman: which Collect, however, is omitted if the body be not present.
% \end{rubric}\par
\begin{rubric}
	If this office follow a Requiem Mass, the Priest takes off his chasuble and maniple, and puts on a black cope. He then proceeds to stand at the foot of the bier, where he says the following prayer.
\end{rubric}


\lett{E}{nter} not into judgment with thy servant, O Lord, for in thy sight shall no man living be justified, except thou grant unto him remission of all his sins. Therefore, we beseech thee, let not the sentence of thy judgment fall upon him, whom the faithful prayer of Christian people commendeth unto thee: but by the succour of thy grace let him who while he lived was sealed with the sign of the Holy Trinity be found worthy to escape the avenging judgment: Who livest and reignest world without end. \textit{Amen.}

\begin{rubric}
    Then, the Cantor beginning, the Clerks standing around sing the following Responsory:
\end{rubric}\par\noindent
Deliver me, O Lord, * from death eternal in that fearful day: * When the heavens and the earth shall be shaken: * When thou shalt come to judge the world by fire. ℣. I am in fear and trembling till the sifting be upon us, and the wrath to come, * When the heavens and the earth shall be shaken. ℣. O that day that day of anger, of calamity and misery, a great day and exceeding bitter. * When thou shalt come to judge the world by fire. ℣. Rest eternal grant unto them, O Lord: and let light perpetual shine upon them.\par
℟. Deliver me, O Lord, * from death eternal in that fearful day: * When the heavens and the earth shall be shaken: * When thou shalt come to judge the world by fire.

\begin{rubric}
    Towards the end of the Responsory, the Celebrant puts incense into the thurible, blessing it as usual, the Deacon ministering the boat.
\end{rubric}
\begin{rubric}
	The Responsory ended, the Priest says,
\end{rubric}

℣. Lord, have mercy upon us.

℟. Christ, have mercy upon us.

℣. Lord, have mercy upon us.

%\begin{rubric}
%    Then the Priest says in a loud voice: \emph{Our Father}, the rest being said secretly by all. Meanwhile the Priest receives the aspergilium from the hand of the Deacon, and makes a reverence to the Altar; the Deacon accompanying him on his right, and holding the fore-edge of the Cope, he makes the circuit of the bier, sprinkling it with Holy Water, thrice on the right side and thrice on the left. When he passes in front of the Cross be bows low, while the Deacon genuflects: afterwards be receives the thurible from the Deacon, and censes the bier in the same manner as he has sprinkled it. Having returned to his first position, the Deacon holding the book, he says with joined hands:
%\end{rubric}
\begin{rubric}
	While the Lord's Prayer is said silently, the Priest goeth about the bier, sprinkling it with Holy Water, thrice on each side; and then in like manner he censes it, thrice on each side. He then says aloud,
\end{rubric}
℣. And lead us not into temptation.

℟. But deliver us from evil.\\

℣. From the gate of hell.

℟. Deliver \textit{his} soul, O Lord.

℣. May \textit{he} rest in peace.

℟. Amen.

℣. O Lord, hear my prayer.

℟. And let my cry come unto thee.

℣. The Lord be with you.

℟. And with thy spirit.

℣. Let us pray.
\par\noindent
	\centerline{\rule{0.5\textwidth}{0.4pt}}
\par\noindent
\begin{rubric}
	If the Body be present, the following Collect shall be said.
\end{rubric}
\lett{O}{God,} whose property is ever to have mercy and to forgive, we humbly entreat thee for the soul of thy \textit{servant} \textit{N.}, which thou hast this day commanded to depart from this world: deliver it not into the hands of the enemy, neither forget it at the last, but command that it be received by thy holy Angels, and brought unto the fatherland of paradise; that forasmuch as \textit{he} hoped and believed in thee, \textit{he} may not undergo the pains of hell, but may obtain everlasting felicity. Through Christ, our Lord. \textit{Amen.}
\par\noindent
	\centerline{\rule{0.5\textwidth}{0.4pt}}
\par\noindent
\begin{rubric}
	If the Body be not present, the following shall be said.
\end{rubric}
\lett{A}{bsolve,} O Lord, we beseech thee, the \textit{soul} of thy \textit{servant} \textit{N.} from every bond of sin: that in the glory of the resurrection, \textit{he} may be raised up amid thy Saints and elect unto newness of life. Through Christ, our Lord. \textit{Amen.}

\begin{rubric}
    Then, the Priest making a cross with his right hand over the bier, says,
\end{rubric}

℣. Rest eternal grant unto \textit{him,} O Lord.

℟. And let light perpetual shine upon \textit{him}.

℣. May \textit{he} rest in peace.

℟. Amen.\par

\begin{rubric}
    The Celebrant, again making a cross over the bier, says without inflexion:
\end{rubric}
℣. May \textit{his} soul and the souls of all the faithful departed through the mercy of God rest in peace.

℟. Amen.
\par\noindent
	\centerline{\rule{0.5\textwidth}{0.4pt}}
\par\noindent

%MANUAL ADJUSTMENT:
\clearpage
\subsection{Service at the Grave}

\begin{rubric}
	While the Body is being borne from the Church to the Grave, the following shall be sung,
\end{rubric}

 \lett{I}{nto} Paradise * may the Angels lead thee: at thy coming may the Martyrs receive thee, and bring thee into the holy city Jerusalem. May the Choirs of Angels receive thee, and with Lazarus, once poor, mayest thou have eternal rest.

\subsubsection{Blessing of the Grave}
\begin{rubric}
    If the ground be unconsecrated, when the Priest and People arrive, the following blessing is said.
\end{rubric}
\letuspray
%From the 1929 Scottish BCP:
\lett{O}{Lord} Jesu Christ, who wast laid in the new tomb of Joseph, and didst thereby sanctify the grave to be a bed of hope to thy people: Vouchsafe, we beseech thee, to bless, ha {\ding{64}} llow, and consecrate this grave, that it may be a resting-place, peaceful and secure, for the body of thy servant which we are about to commit to thy gracious keeping; who art the resurrection and the life, and who livest and reignest with the Father and the Holy Ghost, one God, world without end. \textit{Amen.}
\begin{rubric}
    The Priest then sprinkles the grave with Holy Water, and then censes the Body and the Grave.
\end{rubric}

\subsubsection{At the Grave}
\begin{rubric}
    When they come to the Grave, while the Body is made ready to be laid into the earth, shall be sung or said,
\end{rubric}\par

\lett{M}{an,} that is born of a woman, hath but a short time to live, and is full of misery. He cometh up, and is cut down, like a flower; he fleeth as it were a shadow, and never continueth in one stay.\par
\secondline{In the midst of life we are in death; of whom may we seek for succour, but of thee, O Lord, who for our sins art justly displeased?}
\thirdline{Yet, O Lord God most holy, O Lord most mighty, O holy and most merciful Saviour, deliver us not into the bitter pains of eternal death.}
    Thou knowest, Lord, the secrets of our hearts; shut not thy merciful ears to our prayer; but spare us, Lord most holy, O God most mighty, O holy and merciful Saviour, thou most worthy Judge eternal, suffer us not, at our last hour, for any pains of death, to fall from thee.
\begin{inhead}
    or,
\end{inhead}
\lett{A}{ll} that the Father giveth me shall come to me: and him that cometh to me I will in no wise cast out.\par
    \secondline{He that raised up Jesus from the dead: will also quicken your mortal bodies by the spirit which dwelleth in you.}
    \thirdline{Wherefore my heart is glad, and my glory rejoiceth: my flesh also shall rest in hope.}
    Thou shalt show me the path of life; in thy presence is the fulness of joy: and at thy right hand there is pleasure for evermore.

\begin{rubric}
    Then, while the earth shall be cast upon the Body by some standing by, the Priest shall say,
\end{rubric}

\lett{U}{nto} Almighty God we commend the soul of our \textit{brother} departed, and we commit \textit{his} body to the ground; earth to earth, ashes to ashes, dust to dust; in sure and certain hope of the Resurrection unto eternal life, through our Lord Jesus Christ, at whose coming in glorious majesty to judge the world, the earth and the sea shall give up their dead; and the corruptible bodies of those who sleep in him shall be changed, and made like unto his own glorious body; according to the mighty working whereby he is able to subdue all things unto himself.

\begin{rubric}
    Then shall be said or sung,
\end{rubric}\par\noindent
\antiphon{}{I heard a voice from heaven, saying unto me, Write, From henceforth blessed are the dead who die in the Lord: even so saith the Spirit; for they rest from their labours.}

\lett{B}{lessed} {\ding{64}} be the Lord God of Israel; * for he hath visited and redeemed his people;\par
    \secondline{And hath raised up a mighty salvation for us, * in the house of his servant David;}
    \thirdline{As he spake by the mouth of his holy Prophets, * which have been since the world began;}
    That we should be saved from our enemies, * and from the hand of all that hate us.\par
    To perform the mercy promised to our forefathers, * and to remember his holy covenant;\par
    To perform the oath which he sware to our forefather Abraham, * that he would give us;\par
    That we being delivered out of the hand of our enemies * might serve him without fear;\par
    In holiness and righteousness before him, * all the days of our life.\par
    And thou, child, shalt be called the prophet of the Highest: * for thou shalt go before the face of the Lord to prepare his ways;\par
    To give knowledge of salvation unto his people * for the remission of their sins,\par
    Through the tender mercy of our God; * whereby the day-spring from on high hath visited us;\par
    To give light to them that sit in darkness, and in the shadow of death, * and to guide our feet into the way of peace.

℣. Rest eternal * grant unto them, O Lord.\par
℟. And let light perpetual * shine upon them.

\antiphon{}{I heard a voice from heaven, saying unto me, Write, From henceforth blessed are the dead who die in the Lord: even so saith the Spirit; for they rest from their labours.}\\

℣. The Lord be with you.

℟. And with thy spirit.

℣. Let us pray.

℣. Lord, have mercy upon us.

℟. Christ, have mercy upon us.

℣. Lord, have mercy upon us.

\lett{O}{ur} Father, who art in heaven, Hallowed be thy Name. Thy kingdom come. Thy will be done on earth, As it is in heaven. Give us this day our daily bread. And forgive us our trespasses, As we forgive those who trespass against us. And lead us not into temptation; But deliver us from evil. Amen.

\begin{rubric}
    Meanwhile, the Priest sprinkles the body with Holy Water.
\end{rubric}
%Taking from the English Ritual and providing a rubric for it:
\begin{rubric}
	Then the Priest shall say that which followeth.
\end{rubric}
\lett{G}{rant,} we beseech thee, O Lord, this mercy to thy \textit{servant} departed; that, forasmuch as \textit{he} desired to do thy will, \textit{he} may not suffer the punishment of \textit{his} misdeed: that, as true faith joined \textit{him} to the company of the faithful here below, so in heaven thy mercy may number \textit{him} among the angelic choirs. Through Christ, our Lord. \textit{Amen.}
\begin{rubric}
    The Priest may say one or more of the following Prayers, at his discretion.
\end{rubric}\par

\lett{O}{God,} whose mercies cannot be numbered; Accept our prayers on behalf of the soul of thy \textit{servant} departed, and grant \textit{him} an entrance into the land of light and joy, in the fellowship of thy saints; through Jesus Christ our Lord. \textit{Amen.}\par

\lett{A}{lmighty} God, with whom do live the spirits of those who depart hence in the Lord, and with whom the souls of the faithful, after they are delivered from the burden of the flesh, are in joy and felicity; We give thee hearty thanks for the good examples of all those thy servants, who, having finished their course in faith, do now rest from their labours. And we beseech thee, that we, with all those who are departed in the true faith of thy holy Name, may have our perfect consummation and bliss, both in body and soul, in thy eternal and everlasting glory; through Jesus Christ our Lord. \textit{Amen.}\par

\lett{O}{merciful} God, the Father of our Lord Jesus Christ, who is the Resurrection and the Life; in whom whosoever believeth, shall live, though he die; and whosoever liveth, and believeth in him, shall not die eternally; who also hath taught us, by his holy Apostle Saint Paul, not to be sorry, as men without hope, for those who sleep in him; We humbly beseech thee, O Father, to raise us from the death of sin unto the life of righteousness; that, when we shall depart this life, we may rest in him; and that, at the general Resurrection in the last day, we may be found acceptable in thy sight; and receive that blessing, which thy well-beloved Son shall then pronounce to all who love and fear thee, saying, Come, ye blessed children of my Father, receive the kingdom prepared for you from the beginning of the world. Grant this, we beseech thee, O merciful Father, through Jesus Christ, our Mediator and Redeemer. \textit{Amen.}\par

\lett{T}{he} God of peace, who brought again from the dead our Lord Jesus Christ, the great Shepherd of the sheep, through the blood of the everlasting covenant; Make you perfect in every good work to do his will, working in you that which is well pleasing in his sight; through Jesus Christ, to whom be glory for ever and ever. \textit{Amen.}

\lett{O}{almighty} God, the God of the spirits of all flesh, who by a voice from heaven didst proclaim, Blessed are the dead who die in the Lord; Multiply, we beseech thee, to those who rest in Jesus, the manifold blessings of thy love, that the good work which thou didst begin in them may be perfected unto the day of Jesus Christ. And of thy mercy, O heavenly Father, vouchsafe that we, who now serve thee here on earth, may at last, together with them, be found meet to be partakers of the inheritance of the saints in light; for the sake of the same thy Son Jesus Christ our Lord. \textit{Amen.}\par

\lett{M}{ost} merciful Father, who hast been pleased to take unto thyself the soul of this thy servant; Grant to us who are still in our pilgrimage, and who walk as yet by faith, that having served thee with constancy on earth, we may be joined hereafter with thy blessed saints in glory everlasting; through Jesus Christ our Lord. \textit{Amen.}\par

\lett{O}{Lord} Jesus Christ, who by thy death didst take away the sting of death; Grant unto us thy servants so to follow in faith where thou hast led the way, that we may at length fall asleep peacefully in thee, and awake up after thy likeness; through thy mercy, who livest with the Father and the Holy Ghost, one God, world without end. \textit{Amen.}\par

\lett{A}{lmighty} and everliving God, we yield unto thee most high praise and hearty thanks, for the wonderful grace and virtue declared in all thy saints, who have been the choice vessels of thy grace, and the lights of the world in their several generations; most humbly beseeching thee to give us grace so to follow the example of their stedfastness in thy faith, and obedience to thy holy commandments, that at the day of the general Resurrection, we, with all those who are of the mystical body of thy Son, may be set on his right hand, and hear that his most joyful voice: Come, ye blessed of my Father, inherit the kingdom prepared for you from the foundation of the world. Grant this, O Father, for the sake of the same, thy Son Jesus Christ, our only Mediator and Advocate. \textit{Amen.}\par
\begin{rubric}
	The Priest then continues,
\end{rubric}
℣. Rest eternal grant unto \textit{him}, O Lord.\par
℟. And let light perpetual shine upon \textit{him}.\par
℣. May \textit{he} rest in peace.\par
℟. Amen.

\lett{M}{ay} \textit{his} soul, {\ding{64}} and the souls of all the faithful departed, through the mercy of God, rest in peace. \textit{Amen.}

\lettspace

\begin{rubric}
    The Priest then sprinkles the grave with Holy Water.
\end{rubric}
\par\noindent
	\centerline{\rule{0.5\textwidth}{0.4pt}}
\par\noindent
\subsection{At the Burial of the Dead at Sea}
\begin{rubric}
    The same office may be used; but instead of the Sentence of Committal, the Minister shall say,
\end{rubric}\par

\lett{U}{nto} Almighty God we commend the soul of our brother departed, and we commit his body to the deep; in sure and certain hope of the Resurrection unto eternal life, through our Lord Jesus Christ; at whose coming in glorious majesty to judge the world, the sea shall give up her dead; and the corruptible bodies of those who sleep in him shall be changed, and made like unto his glorious body; according to the mighty working whereby he is able to subdue all things unto himself.
\include{pages/Pastoral/BurialChild}
%Family Prayers
\cleardoublepage
\fancyhead[RE,LO]{}\fancyhead[RO,LE]{}
\fancyhead[C]{}\thispagestyle{empty}
\phantomsection
\addcontentsline{toc}{chapter}{Forms of Prayer to be used in Families}

\begin{figure}[H]
    \centering
    \includegraphics[width=98mm, height=140mm]{Family/Family.eps}
    \caption{\textsc{\Huge{Forms of Prayer to be used in Families}}}
\end{figure}
\include{pages/Family/Family}
%Psalter of David
\cleardoublepage
\fancyhead[RE,LO]{}\fancyhead[RO,LE]{}
\fancyhead[C]{}\thispagestyle{empty}
\phantomsection
\addcontentsline{toc}{chapter}{Psalter of David}

\begin{figure}[H]
    \centering
    \includegraphics[width=98mm, height=140mm]{Psalter/David.eps}
    \caption{\textsc{\Huge{Psalter of David}}}
\end{figure}
\include{pages/Psalter/Psal}
%Back Cover
\include{pages/Back}
\end{document}